% Lesson 58 — Grammar: General revision — Anglais Tle A — Cours signé OnBuch+
% Programme officiel MINESEC. Compiler avec : tectonic EN58-grammar-general-revision.tex
\newcommand{\DOCMATIERE}{Anglais}
\newcommand{\DOCNIVEAU}{Tle A}
\newcommand{\DOCMODULE}{Chapter 5 : Media and Communication}
\newcommand{\DOCLECON}{EN58}
\newcommand{\DOCTITRE}{Lesson 58 — Grammar: General revision}
% OnBuch+ — préambule « cours signés » ENRICHI (compilé avec tectonic / XeLaTeX)
% Système visuel « Pop » d'OnBuch+ : fond crème (jamais blanc), bordures encre
% épaisses, ombres « dures » décalées, titres Archivo Black en capitales.
% Version MATHÉMATIQUES : graphes (pgfplots/tikz), boîtes pédagogiques
% (définition, propriété, méthode, attention, le savais-tu, exercice résolu,
% exercice, TP, bilan), page de garde et signature OnBuch+ ancrée en bas.
%
% Macros attendues dans main.tex AVANT \input{preamble} :
%   \DOCMATIERE  \DOCNIVEAU  \DOCMODULE  \DOCLECON (numéro)  \DOCTITRE
\documentclass[11pt]{article}

\usepackage{fontspec}
\usepackage[english,french]{babel} % français = langue principale ; l'anglais via \eng{..} / textebox
\usepackage[a4paper,margin=2cm,top=3.4cm,bottom=2.8cm,headheight=1.4cm,footskip=1.3cm]{geometry}
\usepackage{amsmath,amssymb,amsfonts}
\usepackage{mathtools} % \xmapsto, \coloneqq... souvent utilisés par les modèles
\usepackage{cancel} % simplifications barrées (bilans de forces)
% \abs{z} (module, valeur absolue) et \norm{u} (norme), utilisés par les modèles
\providecommand{\abs}{}\let\abs\relax\DeclarePairedDelimiter{\abs}{\lvert}{\rvert}
\providecommand{\norm}{}\let\norm\relax\DeclarePairedDelimiter{\norm}{\lVert}{\rVert}
\usepackage[table]{xcolor} % [table] : fournit aussi \rowcolors (alternance de lignes)
\usepackage{pagecolor}
\usepackage{tikz}
\usetikzlibrary{arrows.meta,positioning,calc,shapes.geometric,shapes.misc,shapes.arrows,decorations.pathmorphing,decorations.pathreplacing,decorations.markings,patterns,shadows,mindmap,trees,fit,backgrounds,matrix,intersections,angles,quotes}
\usepackage{pgfplots}
\usepackage[version=4]{mhchem} % équations nucléaires \ce{^{235}_{92}U}
\usepackage[european,siunitx]{circuitikz} % schémas électriques
\pgfplotsset{compat=1.18}
\usepgfplotslibrary{fillbetween,groupplots} % aires entre courbes (calcul intégral)
\usepackage{enumitem}
\usepackage{fancyhdr}
% [most] charge toutes les bibliothèques tcolorbox (listings, minted, raster,
% documentation...) et gonfle inutilement la mémoire TeX sur un document long
% (~300 boîtes) : on ne charge que ce qui est réellement utilisé.
\usepackage[skins,breakable]{tcolorbox}
\usepackage{graphicx}
\usepackage{colortbl}
\usepackage{booktabs}
\usepackage{tabularx}
\usepackage{diagbox} % case d'en-tête diagonale des tableaux à double entrée
% Colonnes extensibles centrée / gauche / droite (C, L, R), souvent utilisées par les modèles
\newcolumntype{C}{>{\centering\arraybackslash}X}
\newcolumntype{L}{>{\raggedright\arraybackslash}X}
\newcolumntype{R}{>{\raggedleft\arraybackslash}X}
\usepackage{array}
\usepackage{multicol}
\usepackage{multirow}
\usepackage{siunitx}
% parse-numbers=false : accepte aussi \SI{2,5 \times 10^{-3}}{...}
\sisetup{locale=FR,output-decimal-marker={,},group-minimum-digits=4,parse-numbers=false}
\DeclareSIUnit\ohm{\ensuremath{\symup{\Omega}}} % U+2126 absent de la police
\DeclareSIUnit\division{div} % réglages d'oscilloscope (ms/div, V/div)
\DeclareSIUnit\tr{tr} % tours (vitesse de rotation, ex. \SI{1500}{\tr\per\minute})
\DeclareSIUnit\molar{M} % concentration molaire (ex. \SI{3}{\molar}, \SI{1}{\micro\molar})
\DeclareSIUnit\an{an} % année (le modèle confond parfois avec \year, un compteur TeX) — ex. \SI{5}{\kilo\meter\per\an}
\DeclareSIUnit\FCFA{FCFA} % franc CFA (ex. \SI{500}{\FCFA\per\kilo\gram})
\DeclareSIUnit\degreeBrix{\textdegree Brix} % teneur en sucre d'un sirop/jus (réfractométrie)
\DeclareSIUnit\month{mois} % le modèle utilise parfois \month, un compteur TeX, comme unité de durée
\usepackage{qrcode}
% \degree : commande fréquemment utilisée par les modèles pour un angle ou une
% température en dehors de \SI{}{} (non fournie par les paquets chargés).
\providecommand{\degree}{^{\circ}}
% \qed (fin de démonstration), utilisé par les modèles sans amsthm
\providecommand{\qed}{\hfill$\square$}
% \barre : double barre des valeurs interdites dans un tableau de variations (tkz-tab)
\providecommand{\barre}{\ensuremath{\|}}
% environnement proof (amsthm non chargé) utilisé par les modèles
\ifdefined\proof\else\newenvironment{proof}[1][Démonstration]{\par\noindent\textit{#1.}\ }{\qed\par}\fi
\usepackage{lastpage}
\usepackage{hyperref}
\hypersetup{hidelinks}

% ============================================================
% Palette « Pop » OnBuch+ — mêmes hex que la classe P (plus_ui.dart)
% ============================================================
\definecolor{poporange}{HTML}{F59321}
\definecolor{popdark}{HTML}{E07A0C}
\definecolor{popcream}{HTML}{FFF6E8}
\definecolor{popink}{HTML}{1C1714}
\definecolor{poppurple}{HTML}{7A5AE0}
\definecolor{popgreen}{HTML}{2E9E6A}
\definecolor{poppink}{HTML}{E0457B}
\definecolor{popblue}{HTML}{2F7DE1}
\definecolor{popgold}{HTML}{C98A12}
\definecolor{popmuted}{HTML}{8A7F76}
\definecolor{popline}{HTML}{EADFD2}
\definecolor{popgray}{HTML}{8A7F76}
\colorlet{popgrey}{popgray}
% Alias tolérants (le modèle nomme parfois les couleurs différemment)
\colorlet{popwhite}{white}
\colorlet{popblack}{popink}
\colorlet{popred}{poppink}
\colorlet{popyellow}{popgold}
% Teintes claires (fonds de tableaux/encadrés) — le modèle nomme parfois
% les couleurs avec un suffixe L (ex. popblueL) sans les avoir définies.
\colorlet{poporangeL}{poporange!20}
\colorlet{popdarkL}{popdark!20}
\colorlet{poppurpleL}{poppurple!20}
\colorlet{popgreenL}{popgreen!20}
\colorlet{poppinkL}{poppink!20}
\colorlet{popblueL}{popblue!20}
\colorlet{popgoldL}{popgold!20}
\colorlet{popredL}{popred!20}
\colorlet{popgrayL}{popgray!20}
% Teintes claires pour fonds de tableaux / zones
\colorlet{popblueL}{popblue!10!white}
\colorlet{poporangeL}{poporange!14!white}
\colorlet{popgreenL}{popgreen!12!white}
\colorlet{poppurpleL}{poppurple!12!white}
\colorlet{poppinkL}{poppink!12!white}
\colorlet{popgoldL}{popgold!14!white}

\pagecolor{popcream}

% ============================================================
% Typographie
% ============================================================
\defaultfontfeatures{Ligatures=TeX}
\setmainfont{PlusJakartaSans-Medium.ttf}[Path=../../../fonts/,BoldFont=PlusJakartaSans-Bold.ttf]
% Maths en sans-serif (Fira Math) avec chiffres et lettres droites de Plus
% Jakarta, pour que les formules se fondent dans le texte.
\usepackage[math-style=ISO,bold-style=ISO]{unicode-math}
\setmathfont{FiraMath-Regular.otf}[Path=../../../fonts/]
\setmathfont{PlusJakartaSans-Medium.ttf}[Path=../../../fonts/,range={up/num,up/{Latin,latin},\mathup}]
\usepackage{icomma} % virgule décimale sans espace en mode math
% Caractères mathématiques Unicode absents des polices de texte (Plus Jakarta,
% Archivo Black) : rendus par leur équivalent mathématique, en texte comme en
% formule — sinon ils s'affichent en carré vide (ex. « Arithmétique dans ℤ »).
\usepackage{newunicodechar}
\newunicodechar{ℝ}{\ensuremath{\mathbb{R}}}
\newunicodechar{ℤ}{\ensuremath{\mathbb{Z}}}
\newunicodechar{ℂ}{\ensuremath{\mathbb{C}}}
\newunicodechar{ℕ}{\ensuremath{\mathbb{N}}}
\newunicodechar{ℚ}{\ensuremath{\mathbb{Q}}}
\newunicodechar{𝔻}{\ensuremath{\mathbb{D}}}
\newunicodechar{α}{\ensuremath{\alpha}}
\newunicodechar{β}{\ensuremath{\beta}}
\newunicodechar{γ}{\ensuremath{\gamma}}
\newunicodechar{δ}{\ensuremath{\delta}}
\newunicodechar{ε}{\ensuremath{\varepsilon}}
\newunicodechar{θ}{\ensuremath{\theta}}
\newunicodechar{λ}{\ensuremath{\lambda}}
\newunicodechar{μ}{\ensuremath{\mu}}
\newunicodechar{σ}{\ensuremath{\sigma}}
\newunicodechar{τ}{\ensuremath{\tau}}
\newunicodechar{φ}{\ensuremath{\varphi}}
\newunicodechar{ψ}{\ensuremath{\psi}}
\newunicodechar{ω}{\ensuremath{\omega}}
\newunicodechar{Σ}{\ensuremath{\Sigma}}
% majuscules produites par \MakeUppercase dans les titres (α-aminés -> Α)
\newunicodechar{Α}{\ensuremath{\alpha}}
\newunicodechar{Β}{\ensuremath{\beta}}
\newunicodechar{Γ}{\ensuremath{\Gamma}}
\newunicodechar{Δ}{\ensuremath{\Delta}}
\newunicodechar{Ε}{\ensuremath{\varepsilon}}
\newunicodechar{Θ}{\ensuremath{\Theta}}
\newunicodechar{Λ}{\ensuremath{\Lambda}}
\newunicodechar{Μ}{\ensuremath{\mu}}
\newunicodechar{Τ}{\ensuremath{\tau}}
\newunicodechar{Φ}{\ensuremath{\Phi}}
\newunicodechar{Ψ}{\ensuremath{\Psi}}
\newunicodechar{Ω}{\ensuremath{\Omega}}
\newunicodechar{↦}{\ensuremath{\mapsto}}
\newunicodechar{∈}{\ensuremath{\in}}
\newunicodechar{∉}{\ensuremath{\notin}}
\newunicodechar{∧}{\ensuremath{\wedge}}
\newunicodechar{⇔}{\ensuremath{\Leftrightarrow}}
\newunicodechar{⟺}{\ensuremath{\Longleftrightarrow}}
\newunicodechar{⇒}{\ensuremath{\Rightarrow}}
\newunicodechar{⟹}{\ensuremath{\Longrightarrow}}
\newunicodechar{∘}{\ensuremath{\circ}}
\newunicodechar{∪}{\ensuremath{\cup}}
\newunicodechar{∩}{\ensuremath{\cap}}
\newunicodechar{≡}{\ensuremath{\equiv}}
\newunicodechar{⊂}{\ensuremath{\subset}}
\newunicodechar{⊥}{\ensuremath{\perp}}
\newunicodechar{∃}{\ensuremath{\exists}}
\newunicodechar{∀}{\ensuremath{\forall}}
\newunicodechar{∛}{\ensuremath{\sqrt[3]{\,}}}
\newunicodechar{∜}{\ensuremath{\sqrt[4]{\,}}}
\newunicodechar{⃗}{\ensuremath{{}^{\to}}}
\newunicodechar{ⁿ}{\ifmmode^{n}\else\textsuperscript{\ensuremath{n}}\fi}
\newunicodechar{ˣ}{\ifmmode^{x}\else\textsuperscript{\ensuremath{x}}\fi}
\newunicodechar{ᵇ}{\ifmmode^{b}\else\textsuperscript{\ensuremath{b}}\fi}
\newunicodechar{ʳ}{\ifmmode^{r}\else\textsuperscript{\ensuremath{r}}\fi}
\newunicodechar{ᵘ}{\ifmmode^{u}\else\textsuperscript{\ensuremath{u}}\fi}
\newunicodechar{ʸ}{\ifmmode^{y}\else\textsuperscript{\ensuremath{y}}\fi}
\newunicodechar{ᵃ}{\ifmmode^{a}\else\textsuperscript{\ensuremath{a}}\fi}
\newunicodechar{ᵉ}{\ifmmode^{e}\else\textsuperscript{\ensuremath{e}}\fi}
\newunicodechar{ᵏ}{\ifmmode^{k}\else\textsuperscript{\ensuremath{k}}\fi}
\newunicodechar{ᵖ}{\ifmmode^{p}\else\textsuperscript{\ensuremath{p}}\fi}
\newunicodechar{ᴺ}{\ifmmode^{N}\else\textsuperscript{\ensuremath{N}}\fi}
\newunicodechar{ᶜ}{\ifmmode^{c}\else\textsuperscript{\ensuremath{c}}\fi}
\newunicodechar{ʰ}{\ifmmode^{h}\else\textsuperscript{\ensuremath{h}}\fi}
\newunicodechar{ᵠ}{\ifmmode^{q}\else\textsuperscript{\ensuremath{q}}\fi}
\newunicodechar{ₙ}{\ifmmode_{n}\else\textsubscript{\ensuremath{n}}\fi}
\newunicodechar{ₐ}{\ifmmode_{a}\else\textsubscript{\ensuremath{a}}\fi}
\newunicodechar{ᵢ}{\ifmmode_{i}\else\textsubscript{\ensuremath{i}}\fi}
\newunicodechar{ᵣ}{\ifmmode_{r}\else\textsubscript{\ensuremath{r}}\fi}
\newunicodechar{ₖ}{\ifmmode_{k}\else\textsubscript{\ensuremath{k}}\fi}
\newunicodechar{ᵦ}{\ifmmode_{\beta}\else\textsubscript{\ensuremath{\beta}}\fi}
\newunicodechar{ₓ}{\ifmmode_{x}\else\textsubscript{\ensuremath{x}}\fi}
\newfontfamily\popdisplay{ArchivoBlack-Regular.ttf}[Path=../../../fonts/]
\newfontfamily\poplogo{SpaceGrotesk-Bold.ttf}[Path=../../../fonts/]
\newfontfamily\popsemi{PlusJakartaSans-SemiBold.ttf}[Path=../../../fonts/]
\newfontfamily\popxbold{PlusJakartaSans-ExtraBold.ttf}[Path=../../../fonts/]
\newfontfamily\popxboldls{PlusJakartaSans-ExtraBold.ttf}[Path=../../../fonts/,LetterSpace=3.0]

\setlist[enumerate]{leftmargin=1.7em,itemsep=5pt,topsep=4pt}
\setlist[itemize]{leftmargin=1.7em,itemsep=4pt,topsep=4pt,label={\color{poporange}\textbullet}}
\setlength{\parindent}{0pt}
\setlength{\parskip}{5pt}
\renewcommand{\arraystretch}{1.25}

% pgfplots : style Pop par défaut
\pgfplotsset{
  every axis/.append style={axis line style={popink,line width=1pt},tick style={popink},
    grid style={popline,line width=0.5pt},label style={font=\small},tick label style={font=\footnotesize}},
  popcurve/.style={poporange,line width=1.8pt,no markers,smooth},
}
% Même style disponible dans un \draw TikZ simple (hors pgfplots)
\tikzset{popcurve/.style={poporange,line width=1.8pt}}
\tikzset{popgrid/.style={popline,very thin}}
\colorlet{popcurve}{poporange} % le nom du style est parfois utilisé comme couleur

% ============================================================
% Pill / squiggle
% ============================================================
\newcommand{\poppill}[3]{%
  \tikz[baseline=-0.55ex]\node[draw=popink,line width=1.2pt,rounded corners=2.6pt,%
    fill=#2,inner xsep=6.5pt,inner ysep=3pt]{%
    \popxboldls\fontsize{7}{7}\selectfont\color{#3}\MakeUppercase{#1}};%
}
\newcommand{\popsquiggle}[1]{%
  \tikz[baseline=-0.4ex]{
    \draw[#1,line width=2.6pt,line cap=round]
      plot [smooth] coordinates {(0,0.09) (0.34,0.01) (0.68,0.21) (1.02,0.01) (1.36,0.10)};
  }%
}
% Mot-clé surligné (marqueur orange sous le texte)
\newcommand{\cle}[1]{{\popxbold\color{popdark}#1}}

% ============================================================
% En-tête / pied
% ============================================================
\pagestyle{fancy}
\fancyhf{}
\renewcommand{\headrulewidth}{0pt}
\renewcommand{\footrulewidth}{0pt}
\lhead{\raisebox{-0.15ex}{\poplogo\fontsize{13}{13}\selectfont\color{popink}OnBuch\color{poporange}+}}
\rhead{\poppill{\DOCMATIERE\ \textbullet\ \DOCNIVEAU\ \textbullet\ Leçon \DOCLECON}{white}{popink}}
\lfoot{\popxbold\fontsize{7.6}{7.6}\selectfont\color{popmuted}\MakeUppercase{Cours signé OnBuch+}\ \textbullet\ {\popsemi\DOCTITRE}}
\rfoot{\tikz[baseline=-0.6ex]\node[rounded corners=8.5pt,fill=popink,minimum height=17pt,minimum width=17pt,inner xsep=5pt,inner sep=0]{\popxbold\fontsize{8}{8}\selectfont\color{popcream}\thepage\,/\,\pageref*{LastPage}};}

% ============================================================
% Titres
% ============================================================
\newcommand{\chaptitre}[2]{%
  \begin{tcolorbox}[enhanced,colback=poporange,colframe=popink,boxrule=2.6pt,arc=11pt,
      drop shadow=popink,left=15pt,right=15pt,top=12pt,bottom=11pt,before skip=3pt,after skip=18pt]
    \poppill{#2}{popink}{popcream}\par\vspace{9pt}
    {\raggedright\hyphenpenalty=10000\exhyphenpenalty=10000\popdisplay\fontsize{22}{24}\selectfont\color{popcream}\MakeUppercase{#1}\par}
  \end{tcolorbox}
}
% Section numérotée : pastille orange avec le numéro + titre Archivo
\newcounter{popsec}
\newcommand{\coursec}[1]{%
  \stepcounter{popsec}\setcounter{popsub}{0}%
  \par\vspace{16pt}\noindent
  \begin{minipage}{\linewidth}
  \tikz[baseline=-0.7ex]\node[draw=popink,line width=1.6pt,fill=poporange,rounded corners=4pt,minimum size=22pt,inner sep=0]{\popdisplay\fontsize{12}{12}\selectfont\color{popcream}\thepopsec};%
  \hspace{8pt}{\popdisplay\fontsize{15}{17}\selectfont\color{popink}\MakeUppercase{#1}}\par
  \vspace{4pt}\noindent{\color{poporange}\rule{36pt}{3.6pt}}
  \end{minipage}\par\nopagebreak\vspace{8pt}
}
\newcounter{popsub}
\newcommand{\courssub}[1]{%
  \stepcounter{popsub}%
  \par\vspace{9pt}\noindent{\popxbold\fontsize{11.5}{13}\selectfont\color{popink}{\color{poporange}\thepopsec.\thepopsub}\hspace{6pt}#1}\par\nopagebreak\vspace{4pt}
}

% ============================================================
% Boîtes pédagogiques — même recette : fond blanc, bordure épaisse colorée,
% ombre dure encre, badge plein. Titre optionnel [#1].
% ============================================================
\tcbset{popbase/.style={breakable,enhanced,colback=white,boxrule=2.3pt,arc=9pt,
  shadow={3pt}{-3pt}{0pt}{popink},left=14pt,right=14pt,top=11pt,bottom=11pt,before skip=11pt,after skip=15pt,
  fontupper=\popsemi\fontsize{10.6}{14.5}\selectfont\color{popink},
  fontlower=\popsemi\fontsize{10.6}{14.5}\selectfont\color{popink}}}
% Coche dessinée (le glyphe ✓ n'existe pas dans les polices du document)
\newcommand{\popcheck}[1]{\tikz[baseline=-0.5ex]\draw[#1,line width=1.6pt,line cap=round,line join=round](-0.13,0.0)--(-0.04,-0.09)--(0.13,0.1);}
\providecommand{\coche}{\popcheck{popgreen}}
\providecommand{\resultat}[1]{\par\smallskip\noindent\fcolorbox{popgreen}{popgreenL}{\parbox{\dimexpr\linewidth-2\fboxsep-2\fboxrule\relax}{\textbf{Résultat.} #1}}\par\smallskip}
\newcommand{\popboxhead}[3]{\poppill{#1}{#2}{white}\ifx\relax#3\relax\else\hspace{6pt}{\popxbold\fontsize{10.6}{12}\selectfont\color{popink}#3}\fi\par\vspace{7pt}}

\newtcolorbox{aretenir}[1][]{popbase,colframe=popgreen,before upper={\popboxhead{À retenir}{popgreen}{#1}}}
% Texte en anglais (document de lecture, dialogue, extrait) : cadre
% d'étude. Un titre optionnel (source, auteur) s'affiche dans l'en-tête.
\newtcolorbox{textebox}[1][]{popbase,colframe=popblue,before upper={\popboxhead{Text}{popblue}{#1}\begin{otherlanguage}{english}},after upper={\end{otherlanguage}}}
% Marques ✓ / ✗ (forme correcte / fautive) souvent utilisées par les modèles
\providecommand{\cross}{\textcolor{poppink}{\ensuremath{\times}}}
\providecommand{\xmark}{\cross}\providecommand{\crossmark}{\cross}\providecommand{\cmark}{\ensuremath{\checkmark}}
% Ligne de réponse / justification à compléter (inventée par les modèles)
\providecommand{\justification}[1]{\par\noindent\textit{Justification~:} \dotfill\par}
% \textipa (package tipa non chargé) : la police ne couvre pas l'API -> simple passage
\providecommand{\textipa}[1]{#1}
% Soulignement (ulem non chargé) : \uline, \ul
\providecommand{\uline}[1]{\underline{#1}}\providecommand{\ul}[1]{\underline{#1}}
% \glossaire{\item ...} inventée par les modèles : liste de glossaire
\providecommand{\glossaire}[1]{\begin{itemize}#1\end{itemize}}
% \alert (beamer) inventée par les modèles : texte mis en évidence
\providecommand{\alert}[1]{\textbf{\textcolor{poppink}{#1}}}
% variantes ASCII d'ordinaux écrites en macro
\providecommand{\iere}{\textsuperscript{re}}\providecommand{\ieme}{\textsuperscript{e}}\providecommand{\ere}{\textsuperscript{re}}\providecommand{\eme}{\textsuperscript{e}}\providecommand{\er}{\textsuperscript{er}}\providecommand{\ie}{\textsuperscript{e}}
% Ordinaux français écrits en macro par les modèles : 1\ière, 3\ième
\newcommand{\ière}{\textsuperscript{re}}\newcommand{\ième}{\textsuperscript{e}}
% \exemple (inventée par les modèles) : étiquette « Exemple : »
\providecommand{\exemple}{\textbf{Exemple~:}\ }
% \square utilisable aussi hors mode math (cases à cocher)
\AtBeginDocument{\let\squareMath\square\renewcommand{\square}{\ensuremath{\squareMath}}}
% Mot ou phrase en anglais à l'intérieur d'un paragraphe français
\newcommand{\eng}[1]{\ifmmode\text{\textit{#1}}\else\begingroup\selectlanguage{english}\itshape #1\endgroup\fi}
\let\esp\eng % alias toléré
% flèches d'intonation inventées par les modèles
\providecommand{\textsearrow}{}\renewcommand{\textsearrow}{\ensuremath{\searrow}}\providecommand{\textnearrow}{}\renewcommand{\textnearrow}{\ensuremath{\nearrow}}
% \fr{..} : mot français (inventé par les modèles)
\providecommand{\fr}[1]{\textit{#1}}
\newtcolorbox{exemplebox}[1][]{popbase,colframe=poporange,before upper={\popboxhead{Exemple}{poporange}{#1}}}
\newtcolorbox{definition}[1][]{popbase,colframe=popblue,before upper={\popboxhead{Définition}{popblue}{#1}}}
\newtcolorbox{propriete}[1][]{popbase,colframe=poppurple,before upper={\popboxhead{Propriété}{poppurple}{#1}}}
\newtcolorbox{methode}[1][]{popbase,colframe=popgold,before upper={\popboxhead{Méthode}{popgold}{#1}}}
\newtcolorbox{attention}[1][]{popbase,colframe=poppink,before upper={\popboxhead{Attention}{poppink}{#1}}}
\newtcolorbox{experience}[1][]{popbase,colframe=popblue,colback=popblueL,before upper={\popboxhead{Expérience}{popblue}{#1}}}
% « Le savais-tu ? » : carte violette pleine (contexte, culture, Cameroun)
\newtcolorbox{savaistu}[1][]{popbase,colback=poppurple,colframe=popink,
  fontupper=\popsemi\fontsize{10.4}{14.2}\selectfont\color{white},
  before upper={\poppill{Le savais-tu\,?}{white}{poppurple}\ifx\relax#1\relax\else\hspace{6pt}{\popxbold\color{white}#1}\fi\par\vspace{7pt}}}
% Exercice résolu : énoncé en haut, \tcblower puis la solution sur fond crème
\newcounter{popexo}\newcounter{popexr}
\newtcolorbox{exoresolu}[1][]{popbase,colframe=poporange,colbacklower=popcream,
  segmentation style={popink,line width=1.2pt,dashed},
  before upper={\stepcounter{popexr}\popboxhead{Exercice résolu \thepopexr}{poporange}{#1}},
  before lower={\poppill{Solution}{popink}{popcream}\par\vspace{6pt}}}
% Exercice (non résolu en place) — la correction est dans la partie Corrigés
\newtcolorbox{exercice}[1][]{popbase,colframe=popink,
  before upper={\stepcounter{popexo}\popboxhead{Exercice \thepopexo}{popink}{#1}}}
% Corrigé
\newtcolorbox{corrige}[1][]{popbase,colframe=popgreen,colback=popgreenL,
  before upper={\popboxhead{Corrigé}{popgreen}{#1}}}
% Objectifs / prérequis / bilan
\newtcolorbox{objectifs}{popbase,colframe=popink,colback=poporangeL,
  before upper={\popboxhead{Objectifs de la leçon}{popink}{}}}
\newtcolorbox{prerequis}{popbase,colframe=popink,colback=popblueL,
  before upper={\popboxhead{Prérequis}{popblue}{}}}
\newtcolorbox{bilan}{popbase,colframe=popink,colback=white,boxrule=2.8pt,
  before upper={\popboxhead{Fiche bilan}{poporange}{L'essentiel à connaître pour l'examen}}}
% Figure encadrée avec légende
\newtcolorbox{popfigure}[1][]{enhanced,breakable=false,colback=white,colframe=popline,boxrule=1.4pt,arc=8pt,
  left=10pt,right=10pt,top=10pt,bottom=8pt,before skip=10pt,after skip=12pt,halign=center,
  lower separated=false,halign lower=center,
  fontlower=\popsemi\fontsize{9}{11.5}\selectfont\color{popmuted},
  #1}
\newcommand{\legende}[1]{\par\vspace{6pt}{\popsemi\fontsize{9}{11.5}\selectfont\color{popmuted}#1}}

% ============================================================
% Signature OnBuch+
% ============================================================
\newcommand{\poplogoinline}[1]{{\poplogo\fontsize{#1}{#1}\selectfont\color{popink}OnBuch\color{poporange}+}}
\newcommand{\signatureonbuch}{%
  \begin{tcolorbox}[enhanced,colback=popink,colframe=popink,boxrule=2.2pt,arc=10pt,
      drop shadow=poporange,left=14pt,right=14pt,top=10pt,bottom=10pt,before skip=0pt,after skip=0pt]
    \tikz[baseline=-0.7ex]\node[circle,fill=popgreen,draw=popcream,line width=1.3pt,minimum size=22pt,inner sep=0]{\popcheck{white}};%
    \hspace{9pt}\begin{minipage}[c]{0.62\linewidth}
      {\popxbold\fontsize{12}{13}\selectfont\color{popcream}Signé\ \textbullet\ {\poplogo OnBuch\color{poporange}+}}\par\vspace{2pt}
      {\raggedright\popsemi\fontsize{8}{10.5}\selectfont\color{popline}\MakeUppercase{Équipe pédagogique OnBuch+}\\\MakeUppercase{Conforme au programme officiel MINESEC}\par}
    \end{minipage}\hfill
    {\poplogo\fontsize{22}{22}\selectfont\color{popcream}OnBuch\color{poporange}+}
  \end{tcolorbox}%
}

% ============================================================
% Page de garde
% #1 titre · #2 matière · #3 niveau · #4 module · #5 n° leçon · #6 durée ·
% #7 description / objectifs (texte ou itemize)
% ============================================================
\newcommand{\pagegarde}[7]{%
  \thispagestyle{empty}
  \newgeometry{a4paper,margin=1.7cm,top=1.6cm,bottom=1.4cm}%
  \begin{tikzpicture}[remember picture,overlay]
    \fill[poporange!18] ([xshift=-3.2cm,yshift=-3.6cm]current page.north east) circle (4.6cm);
    \fill[poppurple!14] ([xshift=2.4cm,yshift=7.6cm]current page.south west) circle (3.8cm);
  \end{tikzpicture}%
  {\poplogo\fontsize{30}{30}\selectfont\color{popink}OnBuch\color{poporange}+}\hfill
  \raisebox{0.6ex}{\poppill{Cours signé \textbullet\ Premium}{popink}{popcream}}\par
  \vspace{1pt}\popsquiggle{poppurple}\par
  \vspace{8pt}
  \begin{tcolorbox}[enhanced,colback=poporange,colframe=popink,boxrule=2.8pt,arc=13pt,
      drop shadow=popink,left=18pt,right=18pt,top=12pt,bottom=13pt,after skip=10pt]
    \poppill{#2 \textbullet\ #3}{popink}{popcream}\hspace{5pt}\poppill{#4}{white}{popink}\par\vspace{8pt}
    {\popdisplay\fontsize{14}{14}\selectfont\color{popink}LEÇON #5}\par\vspace{4pt}
    {\raggedright\hyphenpenalty=10000\exhyphenpenalty=10000\popdisplay\fontsize{25}{27}\selectfont\color{popcream}\MakeUppercase{#1}\par}
    \vspace{8pt}
    {\popxbold\fontsize{9.5}{9.5}\selectfont\color{popink}\MakeUppercase{Durée conseillée : #6}}
  \end{tcolorbox}
  \begin{tcolorbox}[enhanced,colback=white,colframe=popink,boxrule=2.2pt,arc=10pt,
      drop shadow=popink,left=16pt,right=16pt,top=10pt,bottom=10pt,after skip=8pt]
    \poppill{Dans cette leçon}{poppurple}{white}\par\vspace{5pt}
    \popsemi\fontsize{9.3}{12.6}\selectfont\color{popink}#7
  \end{tcolorbox}
  \noindent\begin{minipage}[t]{0.487\linewidth}
    \begin{tcolorbox}[enhanced,colback=popgreen,colframe=popink,boxrule=2.2pt,arc=10pt,
        drop shadow=popink,left=11pt,right=11pt,top=10pt,bottom=10pt,height=3.1cm,valign=center]
      \tcbox[colback=white,colframe=popink,boxrule=1.3pt,arc=5pt,boxsep=0pt,nobeforeafter,
        left=3pt,right=3pt,top=3pt,bottom=3pt,valign=center]{\qrcode[height=1.9cm]{https://play.google.com/store/apps/details?id=com.luvvix.onbuch&hl=fr}}%
      \hspace{8pt}\begin{minipage}[c]{0.5\linewidth}
        {\popdisplay\fontsize{11}{12.5}\selectfont\color{white}\MakeUppercase{Télécharge OnBuch}}\par\vspace{4pt}
        {\raggedright\popsemi\fontsize{8.4}{11}\selectfont\color{white}Gratuit \textbullet\ Google Play\\Tuteur IA, annales, concours, résultats\par}
      \end{minipage}
    \end{tcolorbox}
  \end{minipage}\hfill
  \begin{minipage}[t]{0.487\linewidth}
    \begin{tcolorbox}[enhanced,colback=poppurple,colframe=popink,boxrule=2.2pt,arc=10pt,
        drop shadow=popink,left=11pt,right=11pt,top=10pt,bottom=10pt,height=3.1cm,valign=center]
      \tikz[baseline=-0.7ex]\node[circle,fill=white,draw=popink,line width=1.3pt,minimum size=1.6cm,inner sep=0]{\popdisplay\fontsize{24}{24}\selectfont\color{poppurple}+};%
      \hspace{8pt}\begin{minipage}[c]{0.6\linewidth}
        {\popdisplay\fontsize{11}{12.5}\selectfont\color{white}\MakeUppercase{Passe à OnBuch+}}\par\vspace{4pt}
        {\raggedright\popsemi\fontsize{8.4}{11}\selectfont\color{white}Tous les cours signés, Tuteur IA illimité, fiches PDF — onglet OnBuch+ de l'app\par}
      \end{minipage}
    \end{tcolorbox}
  \end{minipage}
  \vspace{8pt}
  \popcontenu
  \vfill
  \signatureonbuch
  \restoregeometry
  \newpage
}

% Rangée « Ce cours contient » (page de garde)
\newcommand{\poptile}[3]{% #1 couleur #2 gros texte #3 légende
  \begin{tcolorbox}[enhanced,colback=#1,colframe=popink,boxrule=2pt,arc=9pt,shadow={2.5pt}{-2.5pt}{0pt}{popink},
      left=6pt,right=6pt,top=9pt,bottom=9pt,halign=center,valign=center,height=1.75cm,width=\linewidth,nobeforeafter]
    {\popdisplay\fontsize{12.5}{13}\selectfont\color{white}\MakeUppercase{#2}}\par\vspace{3pt}
    {\centering\popsemi\fontsize{7.8}{9.5}\selectfont\color{white}#3\par}
  \end{tcolorbox}}
\newcommand{\popcontenu}{%
  \par\noindent{\popxboldls\fontsize{7.5}{7.5}\selectfont\color{popmuted}CE COURS SIGNÉ CONTIENT}\par\vspace{7pt}
  \noindent\begin{minipage}[t]{0.235\linewidth}\poptile{popblue}{Cours}{complet, illustré, pas à pas}\end{minipage}\hfill
  \begin{minipage}[t]{0.235\linewidth}\poptile{poporange}{Méthodes}{et exercices résolus}\end{minipage}\hfill
  \begin{minipage}[t]{0.235\linewidth}\poptile{poppink}{Exercices}{type examen + corrigés}\end{minipage}\hfill
  \begin{minipage}[t]{0.235\linewidth}\poptile{popgreen}{Fiche bilan}{l'essentiel à retenir}\end{minipage}\par}

% Sommaire
\newtcolorbox{sommaire}{enhanced,colback=white,colframe=popink,boxrule=2.2pt,arc=10pt,
  drop shadow=popink,left=18pt,right=18pt,top=13pt,bottom=13pt,before skip=6pt,after skip=20pt,
  before upper={\poppill{Sommaire}{poppurple}{white}\par\vspace{9pt}\popsemi\fontsize{10.6}{14.5}\selectfont\color{popink}}}

% Signature de fin de cours — poussée tout en bas de la dernière page
\newcommand{\findecours}{%
  \par\vfill\nopagebreak
  \begin{center}\popsquiggle{poporange}\end{center}\vspace{4pt}
  \signatureonbuch
}
\AtBeginDocument{\def\llbracket{\mathopen{[\mkern-3.5mu[}}\def\rrbracket{\mathclose{]\mkern-3.5mu]}}} % intervalles d'entiers (glyphe ⟦ absent de la police)
% Glyphes absents de Fira Math : redéfinis à partir de glyphes disponibles
\AtBeginDocument{%
  \def\setminus{\mathbin{\backslash}}\let\smallsetminus\setminus
  \def\vdots{\mathord{\vbox{\baselineskip3.2pt\lineskiplimit0pt\kern4pt\hbox{.}\hbox{.}\hbox{.}}}}%
  \def\ddots{\mathinner{\mkern1mu\raise6.5pt\hbox{.}\mkern2mu\raise3.6pt\hbox{.}\mkern2mu\raise0.7pt\hbox{.}\mkern1mu}}%
  \def\triangle{\mathord{\scalebox{0.9}{$\symup{\Delta}$}}}%
  \def\nabla{\mathord{\raisebox{\depth}{\scalebox{1}[-1]{$\symup{\Delta}$}}}}%
  \def\checkmark{\popcheck{popdark}}%
  \def\star{\mathbin{\ast}}\def\bigstar{\mathord{\ast}}%
  \def\diamond{\mathbin{\scalebox{0.6}{$\Diamond$}}}%
  \def\bigcup{\mathop{\mathchoice{\raisebox{-0.3em}{\scalebox{1.7}{$\cup$}}}{\scalebox{1.25}{$\cup$}}{\cup}{\cup}}\displaylimits}%
  \def\bigcap{\mathop{\mathchoice{\raisebox{-0.3em}{\scalebox{1.7}{$\cap$}}}{\scalebox{1.25}{$\cap$}}{\cap}{\cap}}\displaylimits}%
  \def\lesseqqgtr{\mathrel{\lessgtr}}\def\bigoplus{\mathop{\oplus}\displaylimits}%
}
\providecommand{\ssi}{si et seulement si} % abréviation fréquente
\AtBeginDocument{\providecommand{\wideparen}{\overparen}} % arc de cercle

%% FIGFIX-BEGIN v4 — correctifs globaux des figures (docs/onbuch-generation/tools/figfix)
\usepackage{adjustbox}\usepackage{etoolbox}
% toute figure TikZ/pgfplots plus large que la ligne est réduite à la largeur disponible
\BeforeBeginEnvironment{tikzpicture}{\begin{adjustbox}{max width=\linewidth}}
\AfterEndEnvironment{tikzpicture}{\end{adjustbox}}
% légendes pgfplots lisibles par-dessus les courbes
\pgfplotsset{every axis/.append style={legend style={fill=white,fill opacity=0.92,text opacity=1,draw=black!15,inner sep=3pt}}}
% étiquettes d'arêtes (syntaxe edge["..."]) avec fond blanc
\tikzset{every edge quotes/.append style={fill=white,inner sep=1.2pt}}
%% FIGFIX-END
\begin{document}

\pagegarde{Lesson 58 — Grammar: General revision}{Anglais}{Tle A}{Chapter 5 : Media and Communication}{EN58}{2 h}{Leçon de révision générale de grammaire pour la Terminale A : synthèse des structures clés du programme (temps, passif, discours rapporté, conditionnels, relatifs, modaux), avec focus sur les erreurs typiques des francophones et entraînement aux exercices de transformation type Bac.
\vspace{4pt}
\begin{multicols}{2}\small
\begin{itemize}
\item À la fin de cette leçon, je sais identifier rapidement la structure grammaticale demandée dans un exercice de transformation
\item Je sais conjuguer correctement les temps principaux (present simple/continuous, past simple/continuous, present perfect, future) en contexte
\item Je sais transformer une phrase active en phrase passive à tous les temps
\item Je sais rapporter un discours direct au discours indirect avec les changements de temps, de pronoms et de marqueurs temporels
\item Je sais construire les trois types de phrases conditionnelles et le souhait avec 'wish'
\item Je sais employer les pronoms relatifs (who, which, that, whose, where) et les modal auxiliaries à bon escient
\end{itemize}\end{multicols}}

\begin{sommaire}
\begin{multicols}{2}
\begin{enumerate}
\item Diagnostic et carte mentale des structures
\item Révision des temps verbaux
\item Voix passive et discours rapporté
\item Conditionnels, souhaits et hypothèses
\item Relatifs, modaux et comparaison
\item Gérondif/infinitif et pièges des francophones
\item Entraînement type Bac et production guidée
\item Activité d'intégration
\item Méthodes et exercices résolus
\item Exercices
\item Corrigés des exercices
\item Fiche bilan
\end{enumerate}
\end{multicols}
\end{sommaire}

\begin{objectifs}
\begin{itemize}
\item À la fin de cette leçon, je sais identifier rapidement la structure grammaticale demandée dans un exercice de transformation
\item Je sais conjuguer correctement les temps principaux (present simple/continuous, past simple/continuous, present perfect, future) en contexte
\item Je sais transformer une phrase active en phrase passive à tous les temps
\item Je sais rapporter un discours direct au discours indirect avec les changements de temps, de pronoms et de marqueurs temporels
\item Je sais construire les trois types de phrases conditionnelles et le souhait avec 'wish'
\item Je sais employer les pronoms relatifs (who, which, that, whose, where) et les modal auxiliaries à bon escient
\item Je sais repérer et corriger les erreurs fréquentes des francophones (faux amis, calques, concordance des temps)
\item Je sais rédiger un court paragraphe en mobilisant plusieurs structures correctement
\end{itemize}
\end{objectifs}

\begin{prerequis}
\begin{itemize}
\item Les temps de base étudiés en Première (present simple et continuous, past simple, present perfect)
\item La voix passive (forme de base : be + participe passé)
\item Le discours rapporté (reported speech) : principes généraux
\item Les phrases conditionnelles types 0, 1, 2 et 3
\item Les pronoms relatifs et les modal auxiliaries (can, must, should, may...)
\end{itemize}
\end{prerequis}

\newpage

\coursec{Diagnostic et carte mentale des structures}

\courssub{Situation d'accroche et problématique}

Imaginez : vous êtes en pleine épreuve d'anglais au Baccalauréat à Yaoundé. La salle est silencieuse, le surveillant annonce qu'il reste trente minutes. Le sujet vous demande de réécrire des phrases au passif, de transformer un discours direct en discours indirect et de corriger des erreurs de concordance des temps. Un candidat de Bamenda hésite : \eng{If I would have known...} ou \eng{If I had known...} ? Une seule réponse est correcte — \eng{If I had known...} — et elle peut valoir des points précieux. Une autre candidate de Douala transforme \eng{The journalist wrote the article} en \eng{The article was wrote by the journalist} : une faute de participe passé (\eng{written}, pas \eng{wrote}) qui coûte la moitié du point.

Cette leçon de révision générale vous donne les clés pour mobiliser rapidement et sans faute toutes les structures étudiées depuis la Première. La problématique est simple : \textbf{comment identifier en quelques secondes la structure grammaticale demandée par un exercice, et l'appliquer correctement ?} La réponse tient en trois étapes : \cle{diagnostiquer} ses forces et ses faiblesses, \cle{cartographier} les structures à maîtriser, et adopter une \cle{méthode d'identification} systématique. C'est le programme de cette première section.

\courssub{Test diagnostique : huit phrases à corriger}

Avant toute révision, il faut savoir où l'on en est. Le test suivant contient huit phrases, chacune comportant \textbf{une seule erreur} portant sur un point majeur du programme de Terminale. Corrigez-les sans consulter vos cahiers, puis comparez avec le corrigé.

\begin{exercice}[Diagnostic — Find and correct the mistake]
Chaque phrase contient une erreur. Soulignez-la et réécrivez la phrase correctement.
\begin{enumerate}
\item \eng{The news are very surprising today.}
\item \eng{She said me that she was tired.}
\item \eng{If I would have money, I would buy a bendskin.}
\item \eng{The match was win by our school team.}
\item \eng{He suggested to go to the market early.}
\item \eng{This is the boy which father works at the radio station.}
\item \eng{I am living in Buea since 2020.}
\item \eng{You must to attend the press conference yesterday.}
\end{enumerate}
\end{exercice}

\begin{corrige}[Test diagnostique]
\begin{enumerate}
\item \eng{The news \textbf{is} very surprising today.} — \eng{news} est un nom singulier indénombrable, malgré son \eng{-s} final. (Le journal est très surprenant aujourd'hui.)
\item \eng{She \textbf{told} me that she was tired.} — on dit \eng{say something} mais \eng{tell somebody}. (Elle m'a dit qu'elle était fatiguée.)
\item \eng{If I \textbf{had} money, I would buy a bendskin.} — jamais \eng{would} dans la proposition introduite par \eng{if} au conditionnel de type 2. (Si j'avais de l'argent, j'achèterais une moto-taxi.)
\item \eng{The match was \textbf{won} by our school team.} — participe passé irrégulier : \eng{win, won, won}. (Le match a été gagné par notre équipe.)
\item \eng{He suggested \textbf{going} to the market early.} — \eng{suggest} est suivi du gérondif, jamais de l'infinitif. (Il a suggéré d'aller au marché tôt.)
\item \eng{This is the boy \textbf{whose} father works at the radio station.} — \eng{whose} exprime la possession ; \eng{which} ne peut pas remplacer un possessif humain. (C'est le garçon dont le père travaille à la station de radio.)
\item \eng{I \textbf{have been living} in Buea since 2020.} — avec \eng{since} + point de départ, l'anglais exige le present perfect (continu). (J'habite à Buea depuis 2020.)
\item \eng{You \textbf{had to attend} the press conference yesterday.} \textbf{ou} \eng{You \textbf{should have attended} the press conference yesterday.} — \eng{must} n'a ni infinitif avec \eng{to}, ni forme passée ; au passé, on utilise \eng{had to} (obligation passée) ou \eng{should have} + participe (reproche/conseil non suivi). (Tu devais assister / Tu aurais dû assister à la conférence de presse hier.)
\end{enumerate}
\textbf{Auto-évaluation :} 7–8 bonnes réponses : excellentes bases, visez la perfection. 4–6 : révision ciblée nécessaire. 0–3 : reprenez chaque chapitre de grammaire méthodiquement.
\end{corrige}

\courssub{Recensement : les huit familles de structures à réviser}

Le test diagnostique révèle les huit grandes familles de structures qui structurent tout le programme de grammaire du cycle et qui tombent régulièrement aux examens. Les voici, avec pour chacune la question-clé à se poser devant une phrase :

\begin{center}
\begin{tabularx}{\linewidth}{|l|X|X|}
\hline
\rowcolor{popblueL} \textbf{Structure} & \textbf{Question-clé} & \textbf{Exemple} \\
\hline
\cle{Tenses} (temps) & Quand l'action se situe-t-elle ? Est-elle finie, en cours, antérieure ? & \eng{I have lived here since 2020.} \\
\hline
\cle{Passive voice} (passif) & Qui subit l'action ? Le sujet est-il agent ou patient ? & \eng{The article was written.} \\
\hline
\cle{Reported speech} (discours rapporté) & Qui rapporte les paroles de qui, et à quel moment ? & \eng{He said he was busy.} \\
\hline
\cle{Conditionals} (conditionnels) & L'hypothèse est-elle réelle, irréelle au présent ou au passé ? & \eng{If I had known, I would have come.} \\
\hline
\cle{Relatives} (relatifs) & Quel mot relie les deux propositions : personne, chose, possession, lieu ? & \eng{The girl who won the prize.} \\
\hline
\cle{Modals} (modaux) & Quelle nuance : obligation, permission, capacité, probabilité, conseil ? & \eng{You should revise daily.} \\
\hline
\cle{Comparison} (comparaison) & Compare-t-on deux éléments ou plus de deux ? & \eng{The most interesting article.} \\
\hline
\cle{Gerund / Infinitive} (gérondif / infinitif) & Le verbe exige-t-il \eng{-ing}, \eng{to} + base, ou les deux ? & \eng{I enjoy reading. / I want to read.} \\
\hline
\end{tabularx}
\end{center}

\begin{aretenir}[La règle d'or de la révision]
Une structure grammaticale se maîtrise en trois temps : \textbf{forme} (comment elle se construit), \textbf{emploi} (quand on l'utilise), \textbf{exceptions} (les cas qui piègent). Pour chaque famille du tableau ci-dessus, vous devez pouvoir réciter ces trois éléments sans hésiter.
\end{aretenir}

\courssub{Construction collective de la carte mentale}

Au tableau, la classe construit ensemble une carte mentale (\eng{mind map}) : le professeur écrit le nœud central, puis chaque rangée propose une branche, sa forme canonique et un exemple. Voici le résultat attendu, enrichi de la branche \eng{Wish / If only} :

\begin{popfigure}
\begin{tikzpicture}[mindmap, grow cyclic, every node/.style={concept}, concept color=popblue,
  root concept/.append style={concept color=popdark, font=\bfseries},
  level 1/.append style={level distance=3.4cm, sibling angle=45},
  level 2/.append style={level distance=2.2cm, sibling angle=40},
  text=white, scale=0.82, transform shape]
\node[root concept]{Grammar revision}
  child[concept color=popblue] { node {Tenses} child { node {since / for} } child { node {past perfect} } }
  child[concept color=popgreen] { node {Passive} child { node {be + past participle} } }
  child[concept color=poporange] { node {Reported speech} child { node {tense shift} } }
  child[concept color=poppurple] { node {Conditionals} child { node {if + past, would} } child { node {Wish / If only} } }
  child[concept color=poppink] { node {Relatives} child { node {who / which / whose} } }
  child[concept color=popgold] { node {Modals} child { node {must / should / might} } }
  child[concept color=popblue!70] { node {Comparison} child { node {-er / more / most} } }
  child[concept color=popgreen!70!popdark] { node {Gerund / Infinitive} child { node {enjoy + -ing} } };
\end{tikzpicture}
\legende{Carte mentale des huit familles de structures à réviser pour le Baccalauréat (avec Wish/If only).}
\end{popfigure}

\begin{experience}[Activité de classe — Build the mind map]
Par groupes de quatre, reproduisez la carte mentale sur une feuille et ajoutez à chaque branche : (1) la forme canonique, (2) un exemple original situé au Cameroun (radio, presse, football, marché), (3) l'erreur typique d'un francophone. Chaque groupe présente ensuite une branche à la classe en anglais, en une minute maximum.
\end{experience}

\courssub{Méthode d'identification : repérer les indices}

Devant un exercice de grammaire, le candidat efficace ne se précipite pas : il \textbf{lit d'abord la consigne}, qui contient presque toujours le nom de la structure attendue ou un verbe d'action révélateur.

\begin{methode}[Identifier la structure demandée en trois étapes]
\begin{enumerate}
\item \textbf{Soulignez le mot-clé de la consigne.} \eng{Turn into the passive voice} annonce le passif ; \eng{Report the following sentences} annonce le discours indirect ; \eng{Rewrite using \emph{if}} annonce le conditionnel ; \eng{Join the sentences with a relative pronoun} annonce les relatifs ; \eng{Correct the mistakes} exige de diagnostiquer soi-même la structure fautive.
\item \textbf{Repérez le temps et les marqueurs de la phrase d'origine.} Cherchez les indices : \eng{since}, \eng{for}, \eng{already}, \eng{yesterday}, \eng{the day before}, un verbe déclencheur comme \eng{said}, \eng{told}, \eng{asked}. Ces marqueurs déterminent la concordance des temps.
\item \textbf{Appliquez la transformation sans changer le sens.} Seule la forme change ; le sens, le temps de référence et les personnes logiques doivent être conservés. Relisez ensuite votre réponse en vérifiant la concordance des temps et l'accord sujet–verbe.
\end{enumerate}
\end{methode}

\begin{exemplebox}[Observer la transformation — actif vers passif]
Phrase d'origine (voix active) : \eng{The journalist wrote the article.} « Le journaliste a écrit l'article. »

Après transformation (voix passive) : \eng{The article was written by the journalist.} « L'article a été écrit par le journaliste. »

Observez le mécanisme : le complément d'objet direct (\eng{the article}) devient sujet ; le verbe prend la forme \eng{be} + participe passé, avec \eng{be} conjugué au temps du verbe actif (\eng{wrote} → prétérit → \eng{was}) ; l'agent est introduit par \eng{by}. Le sens est strictement identique.
\end{exemplebox}

\begin{exoresolu}[Transformation guidée — consigne : \eng{Turn into the passive voice}]
Énoncé : \eng{The students have read the newspaper article.}

\tcblower
Solution détaillée :
\begin{enumerate}
\item Mot-clé de la consigne : \eng{passive voice} → structure visée : le passif.
\item Temps de la phrase d'origine : \eng{have read} → present perfect. Le verbe \eng{be} du passif devra donc être au present perfect : \eng{have/has been}.
\item Le COD \eng{the newspaper article} devient sujet (singulier → \eng{has}) ; participe passé de \eng{read} : \eng{read} (orthographe identique, prononciation « rèd »).
\end{enumerate}
Réponse : \eng{The newspaper article has been read by the students.} « L'article de journal a été lu par les élèves. »

Vérification : sens inchangé, temps conservé (present perfect), accord correct (\eng{has} avec un sujet singulier).
\end{exoresolu}

\begin{attention}[Ne jamais changer le sens lors d'une transformation]
L'erreur la plus grave — et la plus sanctionnée — est de modifier le sens ou le temps de la phrase. Transformez la \textbf{forme}, jamais le \textbf{fond}. Pièges fréquents des francophones :
\begin{itemize}
\item Changer le temps : \eng{She said she \textbf{is} tired} au lieu de \eng{She said she \textbf{was} tired} (concordance des temps obligatoire après un verbe au passé).
\item Oublier le participe passé : \eng{The article was \textbf{wrote}} au lieu de \eng{was \textbf{written}}.
\item Traduire mot à mot : \eng{I am here since Monday} calqué sur « je suis ici depuis lundi » — il faut \eng{I \textbf{have been} here since Monday}.
\item Mettre \eng{would} dans la proposition en \eng{if} : \eng{If I \textbf{would have} known} est toujours faux ; écrivez \eng{If I \textbf{had} known}.
\end{itemize}
\end{attention}

\courssub{Stratégie d'examen : la routine en trois vérifications}

Le jour de l'épreuve, adoptez une routine fixe pour chaque phrase à transformer ou à corriger :

\begin{enumerate}
\item \textbf{Lire la consigne} et nommer mentalement la structure : « c'est du discours rapporté », « c'est un conditionnel de type 3 ».
\item \textbf{Identifier les indices} : temps du verbe principal, verbe déclencheur (\eng{said}, \eng{told}), marqueurs temporels (\eng{yesterday} → \eng{the day before} ; \eng{now} → \eng{then}).
\item \textbf{Vérifier la concordance} : après avoir écrit la réponse, relisez-la en contrôlant trois points : le temps du verbe, l'accord sujet–verbe, et l'orthographe des participes passés irréguliers (\eng{written}, \eng{won}, \eng{told}, \eng{read}).
\end{enumerate}

\begin{center}
\begin{tabularx}{\linewidth}{|X|X|X|}
\hline
\rowcolor{popblueL} Français & English & Remarque \\
\hline
« Il m'a dit qu'il était fatigué. » & \eng{He told me (that) he was tired.} & \eng{tell} + personne ; recul du présent au prétérit \\
\hline
« L'article a été écrit par le journaliste. » & \eng{The article was written by the journalist.} & \eng{be} au prétérit + participe passé irrégulier \\
\hline
« Si j'avais su, je serais venu. » & \eng{If I had known, I would have come.} & jamais \eng{would} après \eng{if} \\
\hline
« J'habite ici depuis 2020. » & \eng{I have lived here since 2020.} & present perfect avec \eng{since}, pas de présent simple \\
\hline
« C'est le garçon dont le père est journaliste. » & \eng{This is the boy whose father is a journalist.} & \eng{whose} pour la possession \\
\hline
\end{tabularx}
\end{center}

\begin{savaistu}[La grammaire au Baccalauréat camerounais]
Au Baccalauréat, les exercices de grammaire portent presque toujours sur cinq points : le \cle{passif}, le \cle{discours rapporté}, le \cle{conditionnel}, les \cle{relatifs} et la \cle{correction d'erreurs} (concordance des temps, verbes irréguliers, gérondif ou infinitif). Autrement dit, la carte mentale construite dans cette leçon couvre l'essentiel de ce qui vous sera demandé. Les candidats qui réussissent ne sont pas ceux qui « savent » la grammaire, mais ceux qui savent \emph{reconnaître} en cinq secondes la structure attendue : c'est exactement la compétence que vous venez d'entraîner.
\end{savaistu}

\begin{aretenir}[À retenir de cette section]
\begin{itemize}
\item Huit familles de structures dominent le programme : temps, passif, discours rapporté, conditionnels (y compris \eng{Wish/If only}), relatifs, modaux, comparaison, gérondif/infinitif.
\item La méthode d'identification repose sur trois gestes : souligner le mot-clé de la consigne, repérer les indices temporels, transformer sans changer le sens.
\item La concordance des temps et les participes passés irréguliers sont les deux sources principales de fautes chez les francophones.
\end{itemize}
Les sections suivantes reprendront chaque branche de la carte mentale, une par une, avec règles, exemples et entraînement.
\end{aretenir}

\coursec{Révision des temps verbaux}

Après avoir dressé notre \emph{diagnostic} et visualisé la \emph{carte mentale} des structures clés du programme, nous entrons maintenant dans le vif du sujet : la maîtrise opérationnelle des temps verbaux. En Terminale, on ne se contente plus de « connaître » les formes ; il faut \textbf{choisir} le temps juste en fonction de la perspective temporelle (aspect) et des marqueurs contextuels. Cette section revoit les six piliers du système verbal anglais en les confrontant systématiquement aux interférences du français.

\courssub{Présent simple vs présent continu : habitude vs instantanéité}

\begin{textebox}[Observation : une journée type à Yaoundé]
Every morning, Pierre \textbf{wakes up} at 5:30 a.m. He \textbf{takes} a bendskin to the Lycée Leclerc. He usually \textbf{buys} a puff-puff on the way. Today, however, it \textbf{is raining} heavily. He \textbf{is waiting} under a shelter because the road \textbf{is flooding}. Look! A taxi \textbf{is coming}. He \textbf{runs} to catch it.
\end{textebox}

\begin{propriete}[Formes et emplois : Present Simple vs Present Continuous]
\begin{itemize}
    \item \textbf{Present Simple} (Base verbale + \textit{-s} à la 3\textsuperscript{e} pers. sg.) \\
    \emph{Emploi :} Habitudes, vérités générales, faits permanents, horaires fixes. \\
    \emph{Marqueurs :} \cle{always}, \cle{usually}, \cle{often}, \cle{sometimes}, \cle{never}, \cle{every day/week/year}, \cle{on Mondays}.
    \item \textbf{Present Continuous} (\textit{be} + \textit{-ing}) \\
    \emph{Emploi :} Action en cours \textbf{au moment de parler} (Now), situation temporaire, futur programmé (avec indicateur de temps futur). \\
    \emph{Marqueurs :} \cle{now}, \cle{at the moment}, \cle{currently}, \cle{this week}, \cle{tonight} (futur), \cle{Look!}, \cle{Listen!}.
\end{itemize}
\end{propriete}

\begin{center}
\begin{tabularx}{\linewidth}{|X|X|X|}
\hline
\rowcolor{popblueL} \textbf{Contexte} & \textbf{Present Simple} & \textbf{Present Continuous} \\
\hline
Habitude & \eng{She teaches English at GBHS.} (Elle enseigne l'anglais au GBHS.) & \eng{She is teaching a difficult class this term.} (Elle donne un cours difficile ce trimestre.) \\
\hline
Vérité générale & \eng{The sun rises in the East.} (Le soleil se lève à l'est.) & -- \\
\hline
Action maintenant & -- & \eng{Look! It is raining in Douala.} (Regarde ! Il pleut à Douala.) \\
\hline
Futur programmé & \eng{The train leaves at 6 p.m.} (Horaires officiels) & \eng{I am meeting the principal tomorrow.} (RDV fixé) \\
\hline
Verbes d'état (stative) & \eng{I know Yaoundé well.} (Je connais bien Yaoundé.) & \textbf{Interdit} : *I am knowing... (sauf sens dynamique : \eng{I am seeing the doctor} = je vais chez le médecin) \\
\hline
\end{tabularx}
\end{center}

\begin{attention}[Piège francophone : « J'habite ici depuis... »]
En français, le présent exprime une action commencée dans le passé et qui se poursuit (\textit{« J'habite à Buea depuis 2010 »}). En anglais, \textbf{c'est impossible} avec le present simple ou continuous.
\begin{itemize}
    \item \textbf{Faux :} \eng{*I live here since 2010.} / \eng{*I am living here since 2010.}
    \item \textbf{Correct :} \eng{I \textbf{have lived} here \textbf{since} 2010.} (Present Perfect -- voir infra).
\end{itemize}
\end{attention}

\begin{exoresolu}[Choix du présent]
\textbf{Énoncé :} Mettez le verbe entre parenthèses au présent simple ou continu.
\begin{enumerate}
    \item Listen! Somebody \_\_\_\_\_\_\_\_ (play) the balafon next door.
    \item My father usually \_\_\_\_\_\_\_\_ (drive) to work, but this week he \_\_\_\_\_\_\_\_ (take) the bus.
    \item Water \_\_\_\_\_\_\_\_ (boil) at 100°C.
    \item The match \_\_\_\_\_\_\_\_ (start) at 4 p.m. tomorrow (official schedule).
\end{enumerate}
\tcblower
\textbf{Solution détaillée :}
\begin{enumerate}
    \item \textbf{is playing} (Action en cours \textit{now} / \textit{Listen!}).
    \item \textbf{drives} (habitude) / \textbf{is taking} (situation temporaire \textit{this week}).
    \item \textbf{boils} (vérité scientifique générale).
    \item \textbf{starts} (horaire officiel $\rightarrow$ present simple). Si c'était un arrangement personnel : \eng{is starting}.
\end{enumerate}
\end{exoresolu}

\courssub{Passé simple vs passé continu : achevé vs en cours}

\begin{propriete}[Règle : Past Simple vs Past Continuous]
\begin{itemize}
    \item \textbf{Past Simple} (Ved / 2\textsuperscript{e} colonne irréguliers) : Action \textbf{terminée}, datée, ponctuelle ou série d'actions successives dans le passé. Marqueurs : \cle{yesterday}, \cle{last week}, \cle{in 2010}, \cle{two days ago}, \cle{then}.
    \item \textbf{Past Continuous} (\textit{was/were} + \textit{-ing}) : Action \textbf{en cours} à un moment précis du passé, ou action longue interrompue par une action courte (past simple).
    \item \textbf{Articulation \cle{when} / \cle{while}} :
    \begin{itemize}
        \item \eng{When} + \textbf{Past Simple} (l'interruption) / \eng{While} + \textbf{Past Continuous} (le fond).
        \item \eng{While I \textbf{was walking} down the street, I \textbf{met} an old friend.}
        \item \eng{When the phone \textbf{rang}, I \textbf{was doing} my homework.}
    \end{itemize}
\end{itemize}
\end{propriete}

\begin{exemplebox}[Exemples contrastés]
\eng{Last night, I \textbf{read} a novel.} (Action terminée, peut-être fini le livre).
\eng{Last night at 10 p.m., I \textbf{was reading} a novel.} (Action en cours à 22h, pas fini).
\eng{The students \textbf{wrote} the exam yesterday.} (Fait accompli).
\eng{The students \textbf{were writing} the exam when the lights \textbf{went} out.} (Interruption).
\end{exemplebox}

\courssub{Present Perfect : le pont passé-présent (le grand piège)}

C'est le temps qui n'a \textbf{pas d'équivalent direct en français}. Il ne situe pas l'action dans le passé (comme le \textit{passé composé}), mais relie le passé au \textbf{moment présent} ( NOW ).

\begin{textebox}[Texte support : Inauguration à Yaoundé]
Last week, while I \textbf{was listening} to the radio, the presenter \textbf{announced} that a new media centre \textbf{had just opened} in Yaoundé. Many young people \textbf{have already visited} it. The director \textbf{said} they \textbf{would organise} workshops next month. I \textbf{have never been} there, but I \textbf{am going to go} this weekend.
\end{textebox}

\begin{propriete}[Present Perfect : Forme, Emplois, Marqueurs, Interdits]
\textbf{Forme :} \textit{have/has} + \textbf{Participe Passé (V3)}.
\begin{center}
\begin{tabular}{|l|l|l|}
\hline
\rowcolor{popblueL} \textbf{Emploi} & \textbf{Exemple} & \textbf{Traduction / Nuance} \\
\hline
Expérience de vie (jamais datée) & \eng{I \textbf{have visited} the Kribi falls.} & J'ai visité les chutes de Kribi (une fois dans ma vie). \\
\hline
Bilan / Résultat présent & \eng{She \textbf{has lost} her phone.} & Elle a perdu son téléphone (elle ne l'a plus \textbf{maintenant}). \\
\hline
Action récente (\cle{just}) & \eng{The bell \textbf{has just rung}.} & La sonnerie vient de retentir. \\
\hline
Non-accomplissement (\cle{yet}, \cle{not... yet}) & \eng{I \textbf{haven't finished} \textbf{yet}.} & Je n'ai pas fini \textbf{encore}. \\
\hline
Durée depuis le passé (\cle{since}, \cle{for}) & \eng{We \textbf{have known} each other \textbf{for} 10 years.} & Nous nous connaissons \textbf{depuis} 10 ans (et encore maintenant). \\
\hline
\end{tabular}
\end{center}

\textbf{Marqueurs typiques :} \cle{ever}, \cle{never}, \cle{just}, \cle{already}, \cle{yet}, \cle{since} (+ point de départ), \cle{for} (+ durée), \cle{recently}, \cle{so far}, \cle{up to now}.

\textbf{\cle{INTERDICTION MAJEURE} :} \textbf{Jamais} de marqueur de temps \textbf{passé révolu} (\cle{yesterday}, \cle{last week}, \cle{in 2015}, \cle{ago}, \cle{when I was young}) avec le Present Perfect. $\rightarrow$ \textbf{Past Simple obligatoire}.
\end{propriete}

\begin{exemplebox}[Le duel Since / For / Ago -- Traduction de « Depuis »]
\begin{itemize}
    \item \eng{I \textbf{have lived} in Douala \textbf{since} 2015.} (Point de départ $\rightarrow$ Present Perfect).
    \item \eng{I \textbf{have lived} in Douala \textbf{for} eight years.} (Durée $\rightarrow$ Present Perfect).
    \item \eng{I \textbf{lived} in Bafoussam \textbf{in} 2015 / \textbf{two years ago}.} (Date passée révolue $\rightarrow$ \textbf{Past Simple}).
    \item \eng{I \textbf{arrived} \textbf{ten minutes ago}.} (\cle{ago} = il y a... $\rightarrow$ Past Simple).
\end{itemize}
\end{exemplebox}

\begin{attention}[Erreur classique : Calque du français « J'habite ici depuis »]
\begin{itemize}
    \item \textbf{Faux :} \eng{*I am living here since 2015.} (Present continuous + since).
    \item \textbf{Faux :} \eng{*I live here since 2015.} (Present simple + since).
    \item \textbf{Correct :} \eng{I \textbf{have lived} here \textbf{since} 2015.} (Present Perfect simple pour état/durée).
    \item \textbf{Nuance :} \eng{I \textbf{have been living} here \textbf{since} 2015.} (Present Perfect Continuous -- insiste sur la continuité de l'action, souvent temporaire).
\end{itemize}
\end{attention}

\begin{exoresolu}[Present Perfect vs Past Simple]
\textbf{Énoncé :} Choisissez le temps correct.
\begin{enumerate}
    \item \eng{Shakespeare \_\_\_\_\_\_\_\_ (write) many plays.} 
    \item \eng{My brother \_\_\_\_\_\_\_\_ (write) three letters this morning.} (Il est encore le matin).
    \item \eng{I \_\_\_\_\_\_\_\_ (see) that movie yesterday.}
    \item \eng{I \_\_\_\_\_\_\_\_ (never / eat) Ndolé before.}
\end{enumerate}
\tcblower
\textbf{Solution :}
\begin{enumerate}
    \item \textbf{wrote} (Shakespeare = personne morte, période révolue $\rightarrow$ Past Simple).
    \item \textbf{has written} (Période \textit{this morning} non terminée $\rightarrow$ Present Perfect + bilan).
    \item \textbf{saw} (\cle{yesterday} = marqueur passé révolu $\rightarrow$ Past Simple).
    \item \textbf{have never eaten} (Expérience de vie + \cle{never} $\rightarrow$ Present Perfect).
\end{enumerate}
\end{exoresolu}

\courssub{Expression du futur : Will vs Going to vs Present Continuous}

\begin{methode}[Comment choisir son futur ? -- Arbre de décision]
\begin{enumerate}
    \item \textbf{Y a-t-il une intention/projet décidé \textbf{AVANT} de parler ?}
    \begin{itemize}
        \item OUI $\rightarrow$ \cle{BE GOING TO + BV} (\eng{I am going to study medicine.} -- Projet mûri).
        \item NON (décision spontanée \textbf{AU MOMENT} de parler) $\rightarrow$ \cle{WILL + BV} (\eng{The phone rings... "I \textbf{will answer} it!"}).
    \end{itemize}
    \item \textbf{Est-ce un arrangement fixe, un rendez-vous programmé (agenda) ?}
    \begin{itemize}
        \item OUI $\rightarrow$ \cle{PRESENT CONTINUOUS} (\eng{I am meeting the Dean at 10 a.m. tomorrow.}).
    \end{itemize}
    \item \textbf{Est-ce une prédiction ?}
    \begin{itemize}
        \item Basée sur des indices présents/preuves $\rightarrow$ \cle{GOING TO} (\eng{Look at those clouds! It \textbf{is going to rain}.}).
        \item Basée sur l'opinion, l'espoir, la certitude subjective $\rightarrow$ \cle{WILL} (\eng{I think Cameroon \textbf{will win} the CAN.}).
    \end{itemize}
    \item \textbf{Promesse, offre, menace, refus spontané ?} $\rightarrow$ \cle{WILL} (\eng{I \textbf{will help} you.} / \eng{The car \textbf{won't start}.}).
\end{enumerate}
\end{methode}

\begin{center}
\begin{tabularx}{\linewidth}{|X|X|X|}
\hline
\rowcolor{poporangeL} \textbf{Futur} & \textbf{Nuance principale} & \textbf{Exemple camerounais} \\
\hline
\cle{Will} & Décision instantanée, prédiction subjective, promesse & \eng{"It's hot." "I \textbf{will open} the window."} \\
\hline
\cle{Be going to} & Intention préétablie, prédiction basée sur preuve & \eng{I \textbf{am going to register} for the Baccalaureate.} \\
\hline
\cle{Present Continuous} & Programme fixe, rendez-vous pris (agenda) & \eng{The Minister \textbf{is visiting} our school next Friday.} \\
\hline
\cle{Present Simple} & Horaires, calendrier officiel (transport, cours) & \eng{The exam \textbf{starts} at 8 a.m. sharp.} \\
\hline
\end{tabularx}
\end{center}

\courssub{Concordance des temps (Backshifting) : le report au passé}

Dans le discours rapporté (\textit{reported speech}) ou les subordonnées introduites par un verbe de parole/pensée au passé (\textit{said, thought, knew, asked}), les temps « reculent » d'un cran vers le passé.

\begin{propriete}[Tableau de concordance standard (Backshifting)]
\begin{center}
\begin{tabularx}{\linewidth}{|X|X|X|}
\hline
\rowcolor{poppurpleL} \textbf{Direct Speech (Présent/Passé)} & \textbf{Reported Speech (Passé)} & \textbf{Exemple} \\
\hline
Present Simple & Past Simple & \eng{"I \textbf{like} fish."} $\rightarrow$ He said he \textbf{liked} fish. \\
\hline
Present Continuous & Past Continuous & \eng{"I \textbf{am reading}."} $\rightarrow$ He said he \textbf{was reading}. \\
\hline
Present Perfect / Past Simple & Past Perfect & \eng{"I \textbf{have finished} / I \textbf{finished}."} $\rightarrow$ He said he \textbf{had finished}. \\
\hline
Will & Would & \eng{"I \textbf{will} come."} $\rightarrow$ He said he \textbf{would} come. \\
\hline
Can & Could & \eng{"I \textbf{can} swim."} $\rightarrow$ He said he \textbf{could} swim. \\
\hline
Must & Had to & \eng{"I \textbf{must} go."} $\rightarrow$ He said he \textbf{had to} go. \\
\hline
\cle{Now} & \cle{Then} / \cle{At that moment} & \\
\cle{Today} & \cle{That day} & \\
\cle{Tomorrow} & \cle{The next day} / \cle{The following day} & \\
\cle{Yesterday} & \cle{The day before} / \cle{The previous day} & \\
\cle{Here} & \cle{There} & \\
\hline
\end{tabularx}
\end{center}

\end{propriete}

\begin{attention}[Exceptions : Pas de backshifting si...]
\begin{enumerate}
    \item Le verbe introducteur est au \textbf{présent} : \eng{He \textbf{says} he \textbf{is} tired.}
    \item La déclaration est une \textbf{vérité générale} intemporelle : \eng{The teacher said the earth \textbf{revolves} around the sun.}
    \item L'action est \textbf{toujours vraie au moment du rapport} (choix stylistique) : \eng{She said she \textbf{loves} / \textbf{loved} Cameroon.}
    \item Modaux déjà « passés » : \cle{could, would, should, might, must (obligation), used to, ought to} ne changent pas.
\end{enumerate}
\end{attention}

\courssub{Synthèse visuelle : La frise chronologique des temps}

\begin{popfigure}
\begin{tikzpicture}[
    timeline/.style={very thick, popdark},
    tick/.style={thick, popblue},
    labelnode/.style={font=\small, text=popdark, align=center},
    pastzone/.style={fill=popblue!15},
    presentzone/.style={fill=popgreen!15},
    futurezone/.style={fill=poporange!15},
    arrowstyle/.style={-{Stealth[length=3mm]}, thick}
]
% Ligne principale
\draw[timeline] (-7,0) -- (7,0);
% Zones
\fill[pastzone] (-7,-0.4) rectangle (-0.5,0.4);
\fill[presentzone] (-0.5,-0.4) rectangle (0.5,0.4);
\fill[futurezone] (0.5,-0.4) rectangle (7,0.4);

% Ticks et labels temps
\foreach \x/\lbl/\dsc in {
    -6/Past Simple/Action termin\'ee ; dat\'ee (yesterday),
    -4/Past Continuous/Action en cours dans le pass\'e (while),
    -2/Present Perfect/Pont pass\'e-pr\'esent (since ; for ; just),
    0/Now/\textbf{NOW},
    2/Present Continuous/Futur programm\'e (tomorrow),
    4/Be going to/Intention ; preuve (look!),
    6/Will/D\'ecision instantan\'ee ; pr\'ediction
} {
    \draw[tick] (\x,0.4) -- (\x,-0.4);
    \node[labelnode, anchor=south] at (\x,0.6) {\textbf{\lbl}};
    \node[labelnode, anchor=north, text width=2.5cm] at (\x,-0.7) {\dsc};
}

% Flèches aspectuelles
\draw[arrowstyle, popblue] (-6.5,1.2) -- (-1.5,1.2) node[midway, above, font=\tiny, popblue] {Pass\'e r\'evolu};
\draw[arrowstyle, popgreen] (-1.5,1.2) -- (0.5,1.2) node[midway, above, font=\tiny, popgreen] {Lien au pr\'esent};
\draw[arrowstyle, poporange] (0.5,1.2) -- (6.5,1.2) node[midway, above, font=\tiny, poporange] {Projection future};

% Marqueurs NOW
\node[font=\bfseries\large, popgreen, anchor=south] at (0,0.5) {NOW};
\end{tikzpicture}
\legende{Frise chronologique : positionnement des temps verbaux par rapport à l'instant de parole (NOW). Zone bleue = Passé révolu ; Verte = Présent / Lien présent ; Orange = Futur.}
\end{popfigure}

\courssub{Entraînement ciblé : Le « Mix » Bac}

\begin{exercice}[Exercices de transformation et de choix -- Type Bac]
\textbf{Exercice 1 :} Mettez les verbes entre parenthèses au temps qui convient (Present Simple, Continuous, Past Simple, Continuous, Present Perfect, Futurs).
\begin{enumerate}
    \item Every year, thousands of tourists \_\_\_\_\_\_\_\_ (visit) Mount Cameroon.
    \item Be quiet! The baby \_\_\_\_\_\_\_\_ (sleep) in the next room.
    \item While we \_\_\_\_\_\_\_\_ (have) dinner, the lights \_\_\_\_\_\_\_\_ (go) out.
    \item I \_\_\_\_\_\_\_\_ (not / see) my cousin \_\_\_\_\_\_\_\_ (since) he \_\_\_\_\_\_\_\_ (leave) for France in 2018.
    \item \_\_\_\_\_\_\_\_ you \_\_\_\_\_\_\_\_ (do) anything tonight? Yes, I \_\_\_\_\_\_\_\_ (meet) friends at the Cultural Centre. (Arrangement fixe).
    \item "I \_\_\_\_\_\_\_\_ (buy) a new phone," she said. (Discours rapporté).
\end{enumerate}

\textbf{Exercice 2 :} Traduisez en anglais (attention aux pièges \cle{depuis}, \cle{venir de}, futur).
\begin{enumerate}
    \item J'habite à Bamenda depuis trois ans.
    \item Il vient de manger un morceau de gâteau.
    \item Quand j'étais jeune, je jouais au football tous les soirs.
    \item Je pense qu'il va pleuvoir (regarde les nuages !).
    \item Le train part dans dix minutes (horaire).
\end{enumerate}
\end{exercice}

\begin{corrige}[Exercices de transformation et de choix]
\textbf{Exercice 1 :}
\begin{enumerate}
    \item \textbf{visit} (Habitude + \textit{Every year} $\rightarrow$ Present Simple).
    \item \textbf{is sleeping} (Action en cours + \textit{Be quiet!} $\rightarrow$ Present Continuous).
    \item \textbf{were having} / \textbf{went} (Action longue interrompue par action courte $\rightarrow$ Past Continuous / Past Simple).
    \item \textbf{haven't seen} / \textbf{since} / \textbf{left} (\textit{since} + date passée $\rightarrow$ Present Perfect principal ; subordonnée \textit{since he left} $\rightarrow$ Past Simple car action ponctuelle passée).
    \item \textbf{Are you doing} / \textbf{am meeting} (Futur programmé/arrangement $\rightarrow$ Present Continuous).
    \item \textbf{would buy} (Discours rapporté : \textit{will} $\rightarrow$ \textit{would}).
\end{enumerate}

\textbf{Exercice 2 :}
\begin{enumerate}
    \item \eng{I \textbf{have lived} in Bamenda \textbf{for} three years.} (ou \eng{I \textbf{have been living}...}). \textbf{Pas} \eng{*I live here since...}
    \item \eng{He \textbf{has just eaten} a piece of cake.} (\textit{Just} + Present Perfect = « venir de »).
    \item \eng{When I \textbf{was young}, I \textbf{played} / \textbf{used to play} football every evening.} (Habitude passée).
    \item \eng{I think it \textbf{is going to rain} (look at those clouds!).} (Prédiction basée sur preuve $\rightarrow$ Going to).
    \item \eng{The train \textbf{leaves} in ten minutes.} (Horaire officiel $\rightarrow$ Present Simple).
\end{enumerate}
\end{corrige}

\begin{savaistu}[Le « Present Perfect » : une exception culturelle ?]
Saviez-vous que l'anglais est l'une des rares langues européennes à posséder un temps grammaticalisé (Present Perfect) qui \textbf{interdit} explicitement les marqueurs de temps passé révolu ? En français, le \textit{passé composé} a largement remplacé le \textit{passé simple} à l'oral et accepte « hier » (\textit{« J'ai mangé hier »}). En allemand standard (\textit{Perfekt}) aussi. Mais en anglais standard (UK/US), \eng{*I have seen him yesterday} reste une erreur grave. Au Cameroun, dans l'anglais camerounais (CamE), on entend parfois des calques comme \eng{*I have seen him yesterday} sous l'influence du français, mais à l'écrit scolaire et au Bac, \textbf{seule la norme standard (Past Simple avec \textit{yesterday}) est acceptée}.
\end{savaistu}

\coursec{Voix passive et discours rapporté}

Dans la section précédente, nous avons révisé les temps verbaux anglais. Or ces temps ne suffisent pas : pour bien s'exprimer, il faut aussi savoir \emph{qui fait l'action} et \emph{comment rapporter les paroles de quelqu'un}. C'est précisément l'objet de cette section : la \cle{voix passive} (\eng{passive voice}) et le \cle{discours rapporté} (\eng{reported speech}). Ces deux structures sont omniprésentes dans les médias — le thème de ce chapitre — et tombent très régulièrement à l'épreuve d'anglais du Baccalauréat.

\courssub{Observer : un article de presse}

\begin{textebox}[A press article — read and observe]
\textbf{Minister Arrested in Corruption Scandal}

Yaoundé — A senior government official \textbf{was arrested} yesterday morning at his residence in the Bastos neighbourhood. The arrest \textbf{was carried out} by officers of the special anti-corruption unit. According to a police spokesman, several documents \textbf{have been seized} and large sums of money \textbf{were found} hidden in the house.

The suspect \textbf{is being questioned} at the central police station. His lawyer told journalists: “My client is innocent. The truth will be established very soon.” The lawyer later \textbf{said that} his client \textbf{was} innocent and that the truth \textbf{would be established} very soon.

A judicial investigation \textbf{has been opened}, and the case \textbf{will be examined} by the High Court next month. “Justice must be done,” declared the Attorney General. He \textbf{added that} justice \textbf{had to be done}.

The minister's resignation \textbf{is expected} before the end of the week. Meanwhile, civil society organisations have asked that the whole affair \textbf{be made} public. “We want transparency,” their spokesperson said. She \textbf{told reporters that} they \textbf{wanted} transparency.

\emph{Glossaire :} to arrest (v., arrêter) ; to carry out (v., mener, effectuer) ; to seize (v., saisir) ; to question (v., interroger) ; to establish (v., établir) ; resignation (n., démission) ; transparency (n., transparence).
\end{textebox}

\textbf{Questions d'observation.}
\begin{enumerate}
\item Souligne tous les verbes en gras construits avec \eng{be} + participe passé. Qui fait l'action dans chaque cas ? Est-ce précisé ?
\item Repère les passages en italique de type \eng{He said that...}. Compare avec les paroles directes : qu'est-ce qui a changé ?
\end{enumerate}

On constate deux phénomènes : (1) l'article insiste sur \emph{ce qui s'est passé} plutôt que sur \emph{qui l'a fait} — c'est le passif ; (2) les paroles des personnes sont \emph{rapportées} avec des transformations — c'est le discours indirect.

\courssub{La voix passive : formation}

\begin{propriete}[Règle d'or du passif]
Le passif se forme avec : \textbf{sujet + be (conjugué au temps de la phrase active) + participe passé} (+ \eng{by} + agent).

\cle{Seul le verbe \eng{be} change de temps} ; le participe passé reste invariable.

\eng{They broadcast the news.} (actif) $\rightarrow$ \eng{The news is broadcast.} « On diffuse les informations. / Les informations sont diffusées. »
\end{propriete}

Dans la transformation, \textbf{l'objet} de la phrase active devient le \textbf{sujet} de la phrase passive, et l'ancien sujet devient le \emph{complément d'agent} introduit par \eng{by} (souvent omis).

\begin{popfigure}
\begin{center}
\begin{tikzpicture}[node distance=0.4cm and 1.6cm]
\node[draw=popblue, fill=popblueL, rounded corners, align=center, inner sep=6pt] (act) {\textbf{ACTIVE}\\ They \quad broadcast \quad the news.\\ \footnotesize sujet \quad verbe \quad objet};
\node[draw=popgreen, fill=popgreenL, rounded corners, align=center, inner sep=6pt, right=of act] (pass) {\textbf{PASSIVE}\\ The news \quad is broadcast \quad (by them).\\ \footnotesize sujet \quad be + p.p. \quad agent};
\draw[-{Stealth[length=3mm]}, line width=1.5pt, poporange] (act) -- (pass) node[midway, above, align=center, font=\footnotesize\bfseries, text=popdark]{objet $\to$ sujet\\ sujet $\to$ by + agent};
\end{tikzpicture}
\legende{Transformation actif $\rightarrow$ passif : inversion objet/sujet}
\end{center}
\end{popfigure}

\begin{center}
\begin{tabularx}{\linewidth}{|l|X|X|}
\hline
\rowcolor{popblueL} \textbf{Temps} & \textbf{Forme passive} & \textbf{Exemple} \\
\hline
Present simple & am/is/are + p.p. & The news \textbf{is broadcast} every evening. \\
\hline
Present continuous & am/is/are being + p.p. & The suspect \textbf{is being questioned}. \\
\hline
Past simple & was/were + p.p. & The minister \textbf{was arrested} yesterday. \\
\hline
Past continuous & was/were being + p.p. & The house \textbf{was being searched}. \\
\hline
Present perfect & have/has been + p.p. & Documents \textbf{have been seized}. \\
\hline
Past perfect & had been + p.p. & The money \textbf{had been hidden}. \\
\hline
Future (will) & will be + p.p. & The case \textbf{will be examined} next month. \\
\hline
Future (be going to) & am/is/are going to be + p.p. & A new studio \textbf{is going to be built} soon. \\
\hline
Modal & modal + be + p.p. & Justice \textbf{must be done}. \\
\hline
\end{tabularx}
\end{center}

\begin{attention}[Participes passés irréguliers]
Le participe passé doit être exact : \eng{broadcast -- broadcast -- broadcast}, \eng{write -- wrote -- written}, \eng{hide -- hid -- hidden}, \eng{do -- did -- done}, \eng{make -- made -- made}. Erreur fréquente : \eng{The news is \textbf{broadcasted}} — faux, le participe est \eng{broadcast}.
\end{attention}

\courssub{Emplois du passif}

\begin{definition}[Quand utilise-t-on le passif ?]
On choisit le passif quand l'\textbf{agent} (celui qui fait l'action) est :
\begin{itemize}
\item \textbf{inconnu} : \eng{My bicycle has been stolen.} « Mon vélo a été volé. » (on ne sait pas par qui) ;
\item \textbf{évident} : \eng{The thief was arrested.} (évidemment par la police) ;
\item \textbf{sans importance} : \eng{Cocoa is exported from Douala.} « Le cacao est exporté depuis Douala. »
\end{itemize}
Le passif est aussi la marque du \textbf{style journalistique et scientifique}, car il paraît objectif et impersonnel.
\end{definition}

\begin{savaistu}[Le passif dans les titres de la presse anglophone]
La presse anglophone (BBC, \emph{The Guardian}, \emph{The New York Times}) utilise massivement le passif dans ses titres, souvent en supprimant même le verbe \eng{be} pour gagner de la place : \eng{Minister arrested in corruption scandal} (= \eng{A minister \textbf{has been} arrested...}), \eng{Two killed in road accident}, \eng{New school opened in Buea}. Au Cameroun, les journaux anglophones comme \emph{The Guardian Post} ou \emph{Cameroon Tribune} (édition anglaise) suivent la même convention.
\end{savaistu}

Comparaison avec le français : le français préfère souvent le pronom \emph{on} ou une tournure active là où l'anglais impose le passif.

\begin{center}
\begin{tabularx}{\linewidth}{|X|X|}
\hline
\rowcolor{popblueL} \textbf{Français} & \textbf{English} \\
\hline
On parle anglais ici. & English \textbf{is spoken} here. \\
\hline
On m'a volé mon téléphone. & My phone \textbf{has been stolen}. \\
\hline
On construit un nouveau pont à Douala. & A new bridge \textbf{is being built} in Douala. \\
\hline
On lui a offert un poste. & She \textbf{was offered} a job. \\
\hline
\end{tabularx}
\end{center}

\begin{attention}[Verbes à deux compléments]
Avec \eng{give, offer, tell, send, show}, l'anglais peut rendre passive la \emph{personne} : \eng{She was offered a job} (impossible en français : « Elle a été offerte un poste » est incorrect). Ne jamais dire \eng{I was said...} pour « on m'a dit » : on dit \eng{I was told...}.
\end{attention}

\courssub{Le discours rapporté : principes généraux}

Rapporter les paroles de quelqu'un, c'est les intégrer dans notre propre phrase. Quand le verbe introducteur est au passé (\eng{said, told, asked}), tout le discours \emph{recule d'un cran dans le temps}.

\begin{methode}[Rapporter un discours en 4 étapes]
\begin{enumerate}
\item Choisir le \textbf{verbe introducteur} au passé : \eng{said (that)}, \eng{told me (that)}, \eng{asked}, \eng{advised}, \eng{invited}.
\item \textbf{Adapter les pronoms et possessifs} selon le nouveau point de vue : \eng{I} $\to$ \eng{he/she}, \eng{my} $\to$ \eng{his/her}, \eng{we} $\to$ \eng{they}.
\item \textbf{Reculer le temps} du verbe : present simple $\to$ past simple, etc.
\item \textbf{Adapter les marqueurs de temps et de lieu} : \eng{now} $\to$ \eng{then}, \eng{today} $\to$ \eng{that day}, \eng{here} $\to$ \eng{there}, \eng{tomorrow} $\to$ \eng{the next day}.
\end{enumerate}
\end{methode}

\begin{propriete}[Le recul des temps — backshift]
\begin{center}
\begin{tabular}{|l|l|}
\hline
\rowcolor{popblueL} \textbf{Discours direct} & \textbf{Discours rapporté} \\
\hline
present simple & past simple \\
\hline
present continuous & past continuous \\
\hline
present perfect & past perfect \\
\hline
past simple & past perfect \\
\hline
will & would \\
\hline
can & could \\
\hline
must & had to \\
\hline
may & might \\
\hline
shall & should / would \\
\hline
\end{tabular}
\end{center}
\emph{Note :} Les modaux \eng{could, would, should, might} ne changent pas. Pas de recul non plus pour une vérité générale encore vraie (\eng{“The earth is round” $\to$ He said the earth is round}) ou si le verbe introducteur est au présent (\eng{He says he is tired}).
\end{propriete}

\begin{exemplebox}[Exemples traduits]
\eng{“I am tired,” she said.} $\to$ \eng{She said (that) she was tired.} « Elle a dit qu'elle était fatiguée. »

\eng{“We have finished our work,” the students said.} $\to$ \eng{The students said they had finished their work.} « Les élèves ont dit qu'ils avaient fini leur travail. »

\eng{The teacher said: “You must work hard.”} $\to$ \eng{The teacher told us (that) we had to work hard.} « Le professeur nous a dit que nous devions travailler dur. »

\eng{“I will call you tomorrow,” he said.} $\to$ \eng{He said he would call me the next day.} « Il a dit qu'il m'appellerait le lendemain. »
\end{exemplebox}

\begin{popfigure}
\begin{center}
\begin{tikzpicture}[node distance=0.35cm and 2.2cm]
\node[draw=popblue, fill=popblueL, rounded corners, align=left, inner sep=6pt] (dir) {\textbf{Direct speech}\\ “I \textbf{am} happy \textbf{here} \textbf{today},”\\ he said.};
\node[draw=poppurple, fill=poppurpleL, rounded corners, align=left, inner sep=6pt, right=of dir] (rep) {\textbf{Reported speech}\\ He said (that) he \textbf{was}\\ happy \textbf{there} \textbf{that day}.};
\draw[-{Stealth[length=3mm]}, line width=1.5pt, poporange] (dir) -- (rep) node[midway, above, align=center, font=\footnotesize\bfseries, text=popdark]{backshift\\ am $\to$ was};
\node[below=0.4cm of dir.south east, align=center, text=popdark, font=\footnotesize] (markers) {here $\to$ there\\ today $\to$ that day};
\draw[-{Stealth[length=2mm]}, popgreen] (dir.south east) -- (markers.north east);
\end{tikzpicture}
\legende{Du discours direct au discours rapporté : recul des temps et des marqueurs}
\end{center}
\end{popfigure}

\begin{center}
\begin{tabularx}{\linewidth}{|X|X|}
\hline
\rowcolor{popblueL} \textbf{Direct speech} & \textbf{Reported speech} \\
\hline
now & then / at that moment \\
\hline
today & that day \\
\hline
yesterday & the day before / the previous day \\
\hline
tomorrow & the next day / the following day \\
\hline
here & there \\
\hline
this / these & that / those \\
\hline
ago & before \\
\hline
last week & the week before \\
\hline
\end{tabularx}
\end{center}

\courssub{Verbes introducteurs, questions et ordres rapportés}

\begin{propriete}[Say ou tell ?]
\begin{itemize}
\item \eng{say (that)...} — sans complément de personne : \eng{He said that he was busy.} « Il a dit qu'il était occupé. »
\item \eng{tell + complément de personne + (that)...} : \eng{He told \textbf{me} that he was busy.} « Il m'a dit qu'il était occupé. »
\end{itemize}
Autres verbes utiles : \eng{ask} (demander), \eng{advise} (conseiller), \eng{invite} (inviter), \eng{warn} (avertir), \eng{promise} (promettre), \eng{explain} (expliquer).
\end{propriete}

\begin{attention}[Erreur fréquente : oublier le complément après tell]
On dit \eng{He told \textbf{me} that...} mais \eng{He said that...}. Jamais \eng{He said me that...} ! De même : \eng{He told that...} sans complément est incorrect. Retiens : \eng{\textbf{say} something} mais \eng{\textbf{tell} somebody}.
\end{attention}

\textbf{Questions rapportées.} Deux cas :

\begin{itemize}
\item \textbf{Question fermée} (oui/non) : on introduit \eng{if} ou \eng{whether} + ordre déclaratif. \eng{“Do you like football?”} $\to$ \eng{He asked me \textbf{if} I liked football.} « Il m'a demandé si j'aimais le football. »
\item \textbf{Question ouverte} (mot interrogatif) : on garde le mot interrogatif + ordre déclaratif. \eng{“Where do you live?”} $\to$ \eng{He asked me \textbf{where I lived}.} « Il m'a demandé où j'habitais. »
\end{itemize}

\begin{attention}[Supprimer l'inversion et l'auxiliaire do/does/did]
Dans une question rapportée, la phrase redevient \textbf{déclarative} : sujet + verbe, sans inversion et sans auxiliaire. On dit \eng{He asked me where I \textbf{lived}} et jamais \eng{where did I live}. De même : \eng{She asked if I \textbf{was}} (pas \eng{if was I}).
\end{attention}

\textbf{Ordres et conseils rapportés.} Structure : \eng{tell / ask / advise / invite + complément + (not) to + infinitif}.

\begin{exemplebox}[Ordres rapportés — exemples traduits]
\eng{“Open your books,” the teacher said.} $\to$ \eng{The teacher told us \textbf{to open} our books.} « Le professeur nous a dit d'ouvrir nos livres. »

\eng{“Don't be late,” my mother said.} $\to$ \eng{My mother told me \textbf{not to be} late.} « Ma mère m'a dit de ne pas être en retard. »

\eng{“You should see a doctor,” she said.} $\to$ \eng{She \textbf{advised} me \textbf{to see} a doctor.} « Elle m'a conseillé de voir un médecin. »
\end{exemplebox}

\begin{exoresolu}[Mise au passif et discours rapporté]
\textbf{Énoncé.} 1) Mets à la voix passive : \eng{Journalists interview the president every year.} 2) Rapporte au discours indirect : \eng{“I have never been to Bamenda,” Mary said.} 3) Rapporte : \eng{“Will you help me tomorrow?” Peter asked Jane.}
\tcblower
\textbf{Solution détaillée.}
1) L'objet \eng{the president} devient sujet ; le verbe est au present simple, donc \eng{be} prend la forme \eng{is} : \eng{The president \textbf{is interviewed} by journalists every year.} « Le président est interviewé par des journalistes chaque année. »

2) Verbe introducteur au passé (\eng{said}) ; pronom \eng{I} $\to$ \eng{she} ; present perfect \eng{have been} $\to$ past perfect \eng{had been} : \eng{Mary said (that) she \textbf{had never been} to Bamenda.} « Mary a dit qu'elle n'était jamais allée à Bamenda. »

3) Question fermée $\to$ \eng{if/whether} ; \eng{will} $\to$ \eng{would} ; \eng{tomorrow} $\to$ \eng{the next day} ; ordre déclaratif : \eng{Peter asked Jane \textbf{if she would help} him \textbf{the next day}.} « Peter a demandé à Jane si elle l'aiderait le lendemain. »
\end{exoresolu}

\begin{exercice}[À toi de jouer]
1) Mets à la voix passive : \eng{The police arrested two suspects last night.} ; \eng{They will build a new radio station in Bamenda.} ; \eng{Someone has stolen my bicycle.}

2) Rapporte au discours indirect : \eng{“I am watching the news now,” John said.} ; \eng{“We won the match yesterday,” the captain said.} ; \eng{“Where did you buy this phone?” she asked me.} ; \eng{“Don't touch the microphone!” the technician told the visitors.}
\end{exercice}

\begin{corrige}[Exercice — Voix passive et discours rapporté]
1) \eng{Two suspects were arrested (by the police) last night.} ; \eng{A new radio station will be built in Bamenda.} ; \eng{My bicycle has been stolen.}

2) \eng{John said (that) he was watching the news then.} ; \eng{The captain said (that) they had won the match the day before.} ; \eng{She asked me where I had bought that phone.} ; \eng{The technician told the visitors not to touch the microphone.}
\end{corrige}

\begin{aretenir}[L'essentiel de la section]
\begin{itemize}
\item Passif : \textbf{be} (au temps voulu) \textbf{+ participe passé} ; seul \eng{be} change.
\item Discours rapporté : verbe introducteur au passé $\to$ recul des temps, adaptation des pronoms et des marqueurs spatio-temporels.
\item \eng{say that...} mais \eng{tell \textbf{me} that...}.
\item Question rapportée : pas d'inversion, pas de \eng{do/does/did}.
\item Ordre rapporté : \eng{tell/ask + quelqu'un + (not) to + infinitif}.
\end{itemize}
\end{aretenir}

\coursec{Conditionnels, souhaits et hypothèses}

Dans la section précédente, nous avons vu comment transformer une phrase active en phrase passive et comment rapporter les paroles de quelqu'un (\eng{reported speech}). Nous quittons maintenant la question du \emph{comment dire} pour celle du \emph{comment imaginer} : comment exprimer une hypothèse, un regret, un souhait ? C'est le domaine des \cle{conditionnels} et de \cle{wish}, un chapitre incontournable de l'épreuve de grammaire du Baccalauréat.

\courssub{Observation : un article de presse}

\begin{textebox}[The Cameroon Weekly — “If only we had acted sooner”]
If the rainy season starts early this year, farmers in the North-West region will plant their maize in March. That is a simple scientific fact: if you plant maize in dry soil, it does not grow. But agriculture is not only about facts. If the government invested more in rural roads, harvests would reach the markets of Bamenda and Buea much faster, and prices would fall for consumers.

An old farmer in Santa told our reporter: “If I were twenty years younger, I would buy a tractor. I wish I had bought one ten years ago, when prices were low. If I had saved more money then, I would have become the richest farmer in the village!” He smiled and added: “Unless the banks help young farmers, nothing will change. Provided that they receive loans, they will modernise their farms. In case you visit us next year, you will see the difference.”
\end{textebox}

\textbf{Glossaire}

\begin{center}
\begin{tabularx}{\linewidth}{|l|l|X|}
\hline
\rowcolor{popblueL} Mot anglais & Nature & Traduction \\
\hline
maize & nom & le maïs \\
\hline
harvest & nom & la récolte \\
\hline
loan & nom & le prêt (bancaire) \\
\hline
to modernise & verbe & moderniser \\
\hline
to postpone & verbe & reporter, remettre à plus tard \\
\hline
\end{tabularx}
\end{center}

\textbf{Questions de compréhension}
\begin{enumerate}
\item What will farmers do if the rainy season starts early?
\item What does the old farmer regret?
\item According to the farmer, what condition would change the situation of young farmers?
\end{enumerate}

\textbf{Réponses attendues :} 1. \eng{They will plant their maize in March.} 2. \eng{He regrets not buying a tractor ten years ago (when prices were low).} 3. \eng{The banks should help them / give them loans.}

Relevons les structures du texte : \eng{If the rainy season starts, farmers will plant} ; \eng{If I were twenty years younger, I would buy} ; \eng{If I had saved more money, I would have become} ; \eng{I wish I had bought one}. Trois constructions différentes, trois rapports différents à la réalité. C'est exactement la logique des trois types de conditionnels.

\courssub{Les quatre types de conditionnels}

\begin{definition}[Conditionnel]
Une phrase conditionnelle comporte deux propositions : la \cle{if-clause} (la condition, introduite par \eng{if}) et la \cle{main clause} (le résultat). L'ordre des deux propositions peut être inversé ; quand la \emph{main clause} vient en premier, on ne met \textbf{pas} de virgule : \eng{If it rains, I will stay at home.} = \eng{I will stay at home if it rains.}
\end{definition}

\begin{propriete}[Type 0 — la vérité générale]
\textbf{Forme :} \eng{If + present simple, present simple.}

\textbf{Emploi :} vérités scientifiques, habitudes, lois générales. \eng{If} signifie ici presque \eng{when} (quand).

\eng{If you heat water to 100 degrees, it boils.} « Si l'on chauffe l'eau à 100 degrés, elle bout. »

\eng{If the Harmattan blows, the air becomes dusty.} « Quand l'harmattan souffle, l'air devient poussiéreux. »
\end{propriete}

\begin{propriete}[Type 1 — l'hypothèse réelle]
\textbf{Forme :} \eng{If + present simple, will + base verbale.}

\textbf{Emploi :} condition \emph{possible} dans le futur ; le locuteur pense que l'événement peut arriver.

\eng{If you revise your lessons, you will pass the Baccalauréat.} « Si tu révises tes leçons, tu réussiras au Bac. »

\eng{If it rains tomorrow, the match will be postponed.} « S'il pleut demain, le match sera reporté. »

\textbf{Variante :} la \emph{main clause} peut contenir un impératif ou un modal : \eng{If you see Paul, tell him to call me.} « Si tu vois Paul, dis-lui de m'appeler. »
\end{propriete}

\begin{propriete}[Type 2 — l'irréel du présent]
\textbf{Forme :} \eng{If + past simple, would + base verbale.}

\textbf{Emploi :} situation imaginaire ou contraire à la réalité présente. Le past simple n'exprime \emph{pas} le passé ici, mais l'éloignement de la réalité.

\eng{If I were rich, I would travel around the world.} « Si j'étais riche, je ferais le tour du monde. » (mais je ne suis pas riche)

\eng{If she lived in Douala, she would visit us every week.} « Si elle habitait à Douala, elle nous rendrait visite chaque semaine. »

\textbf{Exception :} avec le verbe \eng{be}, on utilise \eng{were} pour toutes les personnes à l'écrit soigné : \eng{If I were you}, \eng{if he were taller}.
\end{propriete}

\begin{propriete}[Type 3 — l'irréel du passé]
\textbf{Forme :} \eng{If + past perfect (had + participe passé), would have + participe passé.}

\textbf{Emploi :} regret ou reproche sur le passé ; la condition ne s'est \emph{pas} réalisée et il est trop tard.

\eng{If you had studied, you would have succeeded.} « Si tu avais étudié, tu aurais réussi. » (mais tu n'as pas étudié)

\eng{If they had left earlier, they would not have missed the bus to Yaoundé.} « S'ils étaient partis plus tôt, ils n'auraient pas raté le car pour Yaoundé. »
\end{propriete}

\begin{propriete}[Règle d'or : jamais de will ni de would après if]
Dans la proposition introduite par \eng{if}, on n'emploie \textbf{jamais} \eng{will} ni \eng{would}. Le futur ou le conditionnel se trouve dans la \emph{main clause}, pas dans la \emph{if-clause}.

\eng{If it rains, I will stay at home.} — et non \emph{If it will rain, I will stay at home.}

\eng{If I knew the answer, I would tell you.} — et non \emph{If I would know\dots}
\end{propriete}

\begin{attention}[La double faute du type 3]
La phrase \emph{If I would have known, I would have told you} contient une \textbf{double erreur} : \eng{would} est interdit dans la if-clause, et le type 3 exige le \cle{past perfect}. La forme correcte est :

\eng{If I had known the answer, I would have told you.} « Si j'avais su la réponse, je te l'aurais dite. »

Autre piège : ne pas traduire le « si » français par \eng{if} quand il introduit une interrogation indirecte : « Je ne sais pas s'il viendra » se dit \eng{I don't know whether he will come} (ici \eng{will} est correct, car ce n'est pas une condition !).
\end{attention}

\begin{popfigure}
\begin{tikzpicture}[font=\small]
\node[draw=popblue, fill=popblueL, rounded corners, align=center, minimum width=3.4cm, minimum height=1.1cm] (if1) at (0,0) {\textbf{Type 1}\\ If + present\\ \eng{If you revise}};
\node[draw=popgreen, fill=popgreenL, rounded corners, align=center, minimum width=3.4cm, minimum height=1.1cm] (mc1) at (5.6,0) {\textbf{will + BV}\\ \eng{you will pass}};
\node[draw=poporange, fill=poporangeL, rounded corners, align=center, minimum width=3.4cm, minimum height=1.1cm] (if2) at (0,-1.8) {\textbf{Type 2}\\ If + past simple\\ \eng{If I were rich}};
\node[draw=popgold, fill=popgoldL, rounded corners, align=center, minimum width=3.4cm, minimum height=1.1cm] (mc2) at (5.6,-1.8) {\textbf{would + BV}\\ \eng{I would travel}};
\node[draw=poppurple, fill=poppurpleL, rounded corners, align=center, minimum width=3.4cm, minimum height=1.1cm] (if3) at (0,-3.6) {\textbf{Type 3}\\ If + past perfect\\ \eng{If you had studied}};
\node[draw=poppink, fill=poppinkL, rounded corners, align=center, minimum width=3.4cm, minimum height=1.1cm] (mc3) at (5.6,-3.6) {\textbf{would have + p.p.}\\ \eng{you would have succeeded}};
\draw[-{Stealth}, thick, popblue] (if1) -- (mc1);
\draw[-{Stealth}, thick, poporange] (if2) -- (mc2);
\draw[-{Stealth}, thick, poppurple] (if3) -- (mc3);
\node[align=center, font=\bfseries\small, text=popdark] at (0,1.2) {If-clause};
\node[align=center, font=\bfseries\small, text=popdark] at (5.6,1.2) {Main clause};
\end{tikzpicture}
\legende{Les trois types de conditionnels : chaque couleur associe la forme de la if-clause à celle de la proposition principale.}
\end{popfigure}

\begin{center}
\begin{tabularx}{\linewidth}{|l|X|X|X|}
\hline
\rowcolor{popblueL} Type & If-clause & Main clause & Valeur \\
\hline
0 & present simple & present simple & vérité générale \\
\hline
1 & present simple & will + BV & futur probable \\
\hline
2 & past simple & would + BV & irréel du présent \\
\hline
3 & past perfect & would have + p.p. & irréel du passé \\
\hline
\end{tabularx}
\end{center}

\textbf{Remarque :} les conditionnels \emph{mixtes} existent (une condition passée avec un résultat présent : \eng{If I had learnt English, I would speak it fluently now.} « Si j'avais appris l'anglais, je le parlerais couramment maintenant. »), mais au Baccalauréat on vous demandera surtout de maîtriser parfaitement les types 1, 2 et 3.

\courssub{Wish : exprimer le regret}

\begin{propriete}[Wish + past simple et wish + past perfect]
\textbf{Regret sur le présent :} \eng{wish + past simple} (avec \eng{were} pour \eng{be}).

\eng{I wish I lived near the school.} « J'aimerais habiter près du lycée. » (mais j'habite loin)

\eng{She wishes she were taller.} « Elle aimerait être plus grande. »

\textbf{Regret sur le passé :} \eng{wish + past perfect}.

\eng{I wish I had worked harder last year.} « J'aurais aimé travailler plus dur l'année dernière. »

\eng{We wish we had visited the Waza Park.} « Nous aurions aimé visiter le parc de Waza. »

\textbf{Remarque :} \eng{wish + would} exprime l'irritation ou le souhait que quelqu'un change de comportement : \eng{I wish you would stop making noise!} « J'aimerais que tu arrêtes de faire du bruit ! » On ne peut pas utiliser \eng{wish + would} pour soi-même (\emph{I wish I would\dots} est impossible).
\end{propriete}

\begin{exemplebox}[Exemples traduits]
\eng{If I had known the answer, I would have told you.} « Si j'avais su la réponse, je te l'aurais dite. »

\eng{I wish I had worked harder last year.} « J'aurais aimé travailler plus dur l'année dernière. »

\eng{If only I had listened to my parents!} « Si seulement j'avais écouté mes parents ! » (\eng{if only} renforce le regret et suit les mêmes règles de temps que \eng{wish}.)
\end{exemplebox}

\courssub{Variantes utiles au Bac : unless, provided that, in case}

\begin{definition}[Unless, provided that, in case]
\begin{itemize}
\item \cle{unless} = \eng{if \dots not} : « à moins que ». \eng{Unless you hurry, you will miss the bus.} « À moins que tu ne te dépêches, tu rateras le bus. » (= \eng{If you do not hurry, you will miss the bus.})
\item \cle{provided that} / \cle{as long as} = « à condition que ». \eng{You can go out provided that you finish your homework.} « Tu peux sortir à condition de finir tes devoirs. »
\item \cle{in case} = « au cas où » (précaution). \eng{Take an umbrella in case it rains.} « Prends un parapluie au cas où il pleuvrait. »
\end{itemize}
\end{definition}

\begin{attention}[Unless n'est jamais négatif]
\eng{Unless} contient déjà la négation : on ne dit jamais \emph{unless you don't hurry}. Dites : \eng{Unless you hurry, you will miss the bus.} De même, après \eng{in case}, on utilise le présent, jamais \eng{will} : \eng{in case it rains}, et non \emph{in case it will rain}. Enfin, « à moins que » français est suivi du subjonctif avec \emph{ne} explétif ; l'anglais, lui, garde l'indicatif présent : \eng{unless it rains}.
\end{attention}

\begin{savaistu}[If I were you : le subjonctif anglais]
Pour donner un conseil, l'anglais utilise la forme consacrée \eng{If I were you, I would\dots} (« Si j'étais toi, je\dots »). Ce \eng{were} employé à toutes les personnes est un vestige du \emph{subjonctif} anglais, un mode autrefois aussi riche qu'en français. On le retrouve dans des expressions figées comme \eng{God save the King} ou \eng{be that as it may}. Au Bac, écrire \eng{If I was you} est toléré à l'oral familier, mais \eng{If I were you} reste la norme à l'écrit.
\end{savaistu}

\begin{exoresolu}[Transformation au conditionnel]
\textbf{Énoncé.} Réécrivez chaque phrase avec \eng{if} ou \eng{wish} sans changer le sens.

1. You didn't study, so you failed the test. (type 3)

2. I don't have a car, so I can't drive you to the airport. (type 2)

3. I regret not learning English earlier. (wish)

4. Hurry up, or you will miss the train. (unless)

\tcblower
\textbf{Solution détaillée.}

1. L'événement est passé et non réalisé : type 3 avec past perfect dans la if-clause. \eng{If you had studied, you would have passed the test.} — On remarque la négation qui disparaît : « didn't study » devient la condition positive \eng{had studied}.

2. Situation présente imaginaire : type 2. \eng{If I had a car, I would drive you to the airport.} — Le verbe \eng{have} passe au past simple \eng{had} sans valeur de passé.

3. Regret sur le passé : \eng{wish + past perfect}. \eng{I wish I had learnt English earlier.} — Le participe passé de \eng{learn} est \eng{learnt} (britannique) ou \eng{learned} (américain).

4. \eng{Unless} remplace \eng{if \dots not} : \eng{Unless you hurry up, you will miss the train.} — Pas de négation supplémentaire après \eng{unless}.
\end{exoresolu}

\begin{exercice}[À vous de jouer]
1. Complétez avec la forme correcte du verbe entre parenthèses :
\begin{itemize}
\item If I (be) the President, I (build) more schools in rural areas.
\item If they had left home earlier, they (not/miss) the flight to Lagos.
\item I wish I (know) how to use a computer when I was younger.
\item Provided that you (work) hard, you will succeed.
\end{itemize}

2. Traduisez en anglais : « Si j'avais écouté mes professeurs, j'aurais réussi à l'examen. » et « À moins qu'il ne pleuve, nous jouerons au football. »
\end{exercice}

\begin{corrige}[Exercice « À vous de jouer »]
1. \eng{If I were the President, I would build more schools in rural areas.} (type 2 : \eng{were} à toutes les personnes) ; \eng{If they had left home earlier, they would not have missed the flight to Lagos.} (type 3 : \eng{would not have missed}) ; \eng{I wish I had known how to use a computer when I was younger.} (regret passé : past perfect) ; \eng{Provided that you work hard, you will succeed.} (présent après \eng{provided that}).

2. \eng{If I had listened to my teachers, I would have passed the exam.} ; \eng{Unless it rains, we will play football.}
\end{corrige}

\begin{aretenir}[L'essentiel de la section]
\begin{itemize}
\item Type 0 : \eng{if + present, present} (vérités) ; type 1 : \eng{if + present, will + BV} (futur réel) ; type 2 : \eng{if + past simple, would + BV} (irréel présent) ; type 3 : \eng{if + past perfect, would have + p.p.} (irréel passé).
\item Jamais de \eng{will} ni de \eng{would} dans la if-clause.
\item \eng{Wish + past simple} = regret présent ; \eng{wish + past perfect} = regret passé ; \eng{if only} suit les mêmes règles.
\item \eng{Unless} = \eng{if \dots not} ; \eng{provided that} = à condition que ; \eng{in case} = au cas où.
\end{itemize}
\end{aretenir}

\coursec{Relatifs, modaux et comparaison}

Dans la section précédente, nous avons révisé les conditionnels et les hypothèses (\eng{if I had known…}, \eng{I wish I could…}). Nous abordons maintenant trois piliers de la grammaire anglaise qui tombent chaque année à l'épreuve du Bac : les \cle{pronoms relatifs}, qui permettent de relier deux idées en une seule phrase ; les \cle{auxiliaires modaux}, qui expriment la capacité, l'obligation, le conseil ou la probabilité ; et les \cle{comparatifs et superlatifs}, indispensables pour comparer des personnes, des lieux ou des idées — par exemple deux chaînes de télévision ou deux réseaux sociaux, dans l'esprit de notre chapitre \emph{Media and Communication}.

\courssub{Texte d'observation : une journaliste de la CRTV}

\begin{textebox}[A young reporter in Yaoundé — reading text]
Manka is a young journalist who works for CRTV in Yaoundé. She studied at the school of journalism which is located near the university. Last year, she interviewed a minister whose answers surprised everybody. The report, which was broadcast on the evening news, was watched by thousands of viewers.

Manka believes that a good reporter must check every fact before publishing it. “We mustn't spread false information,” she says, “but we don't have to reveal all our sources.” She thinks young journalists should learn to use social media carefully, because a message which is shared too quickly can become dangerous. When she was a student, she could write three articles a week, but she couldn't film and edit videos. Now she can do both.

Her editor told her: “You should have sent your article earlier yesterday.” Manka replied that she had had to wait for a document that arrived late. She says journalism is more difficult than people imagine, but it is also the most exciting job she knows. “The more you practise,” she smiles, “the better you become.”
\end{textebox}

\textbf{Glossaire :} \eng{to broadcast} (v. irrégulier : \eng{broadcast, broadcast, broadcast}) diffuser ; \eng{viewers} (n.) téléspectateurs ; \eng{to spread} (v. irrégulier : \eng{spread, spread, spread}) répandre ; \eng{sources} (n.) sources ; \eng{editor} (n.) rédacteur en chef ; \eng{to edit} (v.) monter, éditer.

\textbf{Questions de compréhension :}
\begin{enumerate}
\item Where does Manka work and where did she study?
\item What must a good reporter do before publishing, according to Manka?
\item What could Manka do when she was a student, and what couldn't she do?
\item What reproach did her editor make to her?
\end{enumerate}

\begin{corrige}[Questions de compréhension]
\begin{enumerate}
\item She works for CRTV in Yaoundé and she studied at the school of journalism near the university.
\item A good reporter must check every fact before publishing it.
\item She could write three articles a week, but she couldn't film and edit videos.
\item He told her she should have sent her article earlier.
\end{enumerate}
\end{corrige}

\courssub{Les pronoms relatifs}

\begin{definition}[Pronom relatif]
Un \cle{pronom relatif} relie deux phrases en remplaçant un nom déjà mentionné (l'\emph{antécédent}). Il évite la répétition : \eng{The journalist interviewed me. The journalist works for CRTV.} devient \eng{The journalist who interviewed me works for CRTV.} « Le journaliste qui m'a interviewé travaille pour la CRTV. »
\end{definition}

\begin{center}
\begin{tabularx}{\linewidth}{|l|l|X|X|}
\hline
\rowcolor{popblueL} Relatif & Antécédent & Exemple anglais & Traduction \\
\hline
\eng{who} & personnes (sujet ou complément) & \eng{The reporter who called me is kind.} & Le journaliste qui m'a appelé est gentil. \\
\hline
\eng{whom} & personnes (complément, style soutenu) & \eng{The minister whom she interviewed.} & Le ministre qu'elle a interviewé. \\
\hline
\eng{which} & choses, animaux, idées & \eng{The newspaper which I buy is serious.} & Le journal que j'achète est sérieux. \\
\hline
\eng{that} & personnes ou choses (relatives restrictives) & \eng{The news that she read was true.} & Les informations qu'elle a lues étaient vraies. \\
\hline
\eng{whose} & possession (personnes et choses) & \eng{A journalist whose articles I admire.} & Un journaliste dont j'admire les articles. \\
\hline
\eng{where} & lieu & \eng{This is the school where I studied.} & Voici l'école où j'ai étudié. \\
\hline
\eng{when} & temps & \eng{I remember the day when we met.} & Je me souviens du jour où nous nous sommes rencontrés. \\
\hline
\end{tabularx}
\end{center}

\begin{exemplebox}[Exemples traduits]
\begin{itemize}
\item \eng{The journalist who interviewed me works for CRTV.} « Le journaliste qui m'a interviewé travaille pour la CRTV. »
\item \eng{This is the school where I studied.} « Voici l'école où j'ai étudié. »
\item \eng{The phone which she bought in Douala is already broken.} « Le téléphone qu'elle a acheté à Douala est déjà cassé. »
\item \eng{That's the cameraman whose camera was stolen.} « C'est le cameraman dont la caméra a été volée. »
\end{itemize}
\end{exemplebox}

\begin{propriete}[Relatives restrictives et non restrictives]
\textbf{Relative restrictive} (déterminative) : elle est indispensable pour identifier l'antécédent. Elle n'est \emph{jamais} encadrée de virgules et accepte \eng{that}. \eng{Students who work hard pass the exam.} « Les élèves qui travaillent dur réussissent » (pas tous les élèves, seulement ceux-là).

\textbf{Relative non restrictive} (explicative) : elle ajoute une information secondaire. Elle est \emph{toujours encadrée par des virgules} et l'on \emph{ne peut jamais} utiliser \eng{that}. \eng{My father, who is a teacher, lives in Buea.} « Mon père, qui est enseignant, vit à Buea. » (Je n'ai qu'un père ; la relative ajoute un détail.)

\textbf{Préposition + relatif} : en style soutenu, la préposition se place devant \eng{whom} ou \eng{which} (jamais devant \eng{that}) : \eng{The person to whom I spoke was helpful.} « La personne à qui j'ai parlé était serviable. » En langue courante, la préposition va en fin de phrase : \eng{The person (who) I spoke to was helpful.}
\end{propriete}

\begin{attention}[Pièges des francophones]
\begin{itemize}
\item Ne dites jamais \eng{the man which} pour une personne : c'est \eng{the man who/that}.
\item \eng{whose} ne s'emploie pas seulement pour les personnes ; il remplace le français « dont » : \eng{the girl whose brother is famous} « la fille dont le frère est célèbre ».
\item Quand le relatif est complément d'objet, on peut l'omettre : \eng{The film (that) I saw was boring.} « Le film que j'ai vu était ennuyeux. » Mais jamais quand il est sujet : \eng{The man who came…} (on ne peut pas supprimer \eng{who}).
\item Ne confondez pas \eng{what} et \eng{that} : \eng{what} signifie « ce que / ce qui » et n'a \emph{pas d'antécédent} : \eng{I heard what you said.} « J'ai entendu ce que tu as dit. » Mais : \eng{I heard the news that you announced.} « J'ai entendu la nouvelle que tu as annoncée. »
\end{itemize}
\end{attention}

\begin{savaistu}[Le saviez-vous ?]
Dans un style soutenu, \eng{whose} peut aussi s'employer pour des choses, là où le français dirait « dont » : \eng{a country whose economy is growing} « un pays dont l'économie est en croissance ». C'est très apprécié dans les rédactions du Bac !
\end{savaistu}

\courssub{Les auxiliaires modaux}

\begin{definition}[Auxiliaire modal]
Les \cle{modaux} sont de petits auxiliaires (\eng{can, could, must, should, may, might}…) placés devant la base verbale pour exprimer une attitude : capacité, permission, obligation, conseil, probabilité.
\end{definition}

\begin{propriete}[Propriétés communes des modaux]
\begin{enumerate}
\item Ils sont suivis de la \emph{base verbale sans to} : \eng{He must go.} (jamais \eng{He must to go}).
\item Ils ne prennent \emph{jamais de -s} à la 3e personne : \eng{She can swim.} (jamais \eng{She cans swim}).
\item Ils n'ont ni infinitif ni participe : pour les conjuguer à d'autres temps, on utilise des équivalents (\eng{to be able to} pour \eng{can}, \eng{to have to} pour \eng{must}, \eng{to be allowed to} pour la permission) : \eng{She will be able to travel.} « Elle pourra voyager. »
\item Négation et question sans \eng{do} : \eng{You should not smoke.} / \eng{May I come in?}
\item Exceptions : \eng{ought to} et \eng{have to} prennent \eng{to} ; \eng{have to} se conjugue avec \eng{do} : \eng{Does she have to leave?} / \eng{She doesn't have to leave.}
\end{enumerate}
\end{propriete}

\begin{popfigure}
\begin{center}
\begin{tikzpicture}
\draw[fill=popblueL, draw=popblue, thick] (0,0) rectangle (2.6,3.4);
\draw[fill=poporangeL, draw=poporange, thick] (3.2,0) rectangle (6.6,3.4);
\draw[fill=popgreenL, draw=popgreen, thick] (7.2,0) rectangle (11.2,3.4);
\node[align=center, text=popdark, font=\bfseries] at (1.3,3.0) {Modal};
\node[align=center, text=popdark, font=\bfseries] at (4.9,3.0) {Valeur};
\node[align=center, text=popdark, font=\bfseries] at (9.2,3.0) {Exemple};
\node[align=left, text=popdark, font=\small] at (1.3,1.5) {can / could\\must / have to\\should / ought to\\may / might\\mustn't\\don't have to};
\node[align=left, text=popdark, font=\small] at (4.9,1.5) {capacité, permission\\obligation\\conseil\\probabilité\\interdiction\\absence d'obligation};
\node[align=left, text=popdark, font=\small] at (9.2,1.5) {She can type fast.\\You must stop here.\\You should rest.\\It may rain.\\You mustn't smoke.\\You needn't come.};
\end{tikzpicture}
\end{center}
\legende{Tableau TikZ des modaux : modal, valeur, exemple.}
\end{popfigure}

\begin{center}
\begin{tabularx}{\linewidth}{|l|l|X|}
\hline
\rowcolor{popblueL} Modal & Valeur & Exemple traduit \\
\hline
\eng{can / could} & capacité, permission & \eng{She can speak three languages.} « Elle sait parler trois langues. » \\
\hline
\eng{must / have to} & obligation (interne / externe) & \eng{I must finish this article.} « Je dois finir cet article. » \\
\hline
\eng{should / ought to} & conseil & \eng{You should check your sources.} « Tu devrais vérifier tes sources. » \\
\hline
\eng{may / might} & probabilité, permission polie & \eng{It may rain in Buea.} « Il se peut qu'il pleuve à Buea. » \\
\hline
\eng{mustn't} & interdiction & \eng{You mustn't lie.} « Il est interdit de mentir. » \\
\hline
\eng{don't have to} & pas nécessaire & \eng{You don't have to pay.} « Tu n'es pas obligé de payer. » \\
\hline
\end{tabularx}
\end{center}

\begin{attention}[mustn't ou don't have to ?]
Confusion classique des francophones ! \eng{mustn't} exprime une \emph{interdiction} : \eng{You mustn't smoke here.} « Il est interdit de fumer ici. » \eng{don't have to} exprime l'\emph{absence d'obligation} : \eng{You don't have to come.} « Tu n'es pas obligé de venir » (mais tu peux). Ne traduisez donc jamais « tu ne dois pas » automatiquement par \eng{mustn't}.
\end{attention}

\begin{propriete}[Modaux au passé : modal + have + participe passé]
\begin{itemize}
\item \eng{should have + p.p.} : reproche, regret. \eng{You should have revised earlier.} « Tu aurais dû réviser plus tôt. »
\item \eng{could have + p.p.} : possibilité manquée. \eng{We could have won the match.} « Nous aurions pu gagner le match. »
\item \eng{must have + p.p.} : déduction sur le passé (certitude). \eng{She must have missed the bus.} « Elle a dû rater le bus. »
\item \eng{can't have + p.p.} : impossibilité au passé. \eng{He can't have said that!} « Il ne peut pas avoir dit ça ! »
\end{itemize}
\end{propriete}

\courssub{Comparatifs et superlatifs}

\begin{propriete}[Formation]
\begin{itemize}
\item Adjectifs \emph{courts} (une syllabe, ou deux en -y) : \emph{-er} / \emph{the -est}. \eng{tall, taller, the tallest} ; \eng{happy, happier, the happiest}.
\item Adjectifs \emph{longs} (deux syllabes et plus) : \eng{more} / \eng{the most}. \eng{interesting, more interesting, the most interesting}.
\item Irréguliers : \eng{good / better / the best} ; \eng{bad / worse / the worst} ; \eng{far / farther / the farthest} (distance) ou \eng{far / further / the furthest} (distance ou « supplémentaire »).
\item Attention au redoublement de la consonne finale après une voyelle courte : \eng{big, bigger, the biggest} ; \eng{hot, hotter, the hottest}.
\end{itemize}
\end{propriete}

\begin{popfigure}
\begin{center}
\begin{tikzpicture}
\draw[fill=popblueL, draw=popblue, thick] (0,0) rectangle (1.2,1.6);
\node[align=center, text=popdark] at (0.6,1.9) {\eng{tall}};
\draw[fill=poporangeL, draw=poporange, thick] (2.4,0) rectangle (3.6,2.4);
\node[align=center, text=popdark] at (3.0,2.7) {\eng{taller}};
\draw[fill=popgreenL, draw=popgreen, thick] (4.8,0) rectangle (6.0,3.2);
\node[align=center, text=popdark] at (5.4,3.5) {\eng{the tallest}};
\draw[-{Stealth}, thick, popdark] (1.3,0.8) -- (2.3,1.2);
\draw[-{Stealth}, thick, popdark] (3.7,1.2) -- (4.7,1.6);
\node[align=center, text=popmuted, font=\small] at (3.0,-0.5) {gradation : positif, comparatif, superlatif};
\end{tikzpicture}
\end{center}
\legende{Gradation comparative : tall, taller, the tallest.}
\end{popfigure}

\begin{propriete}[Structures de comparaison]
\begin{itemize}
\item Égalité : \eng{as + adjectif + as}. \eng{Radio is as important as television.} « La radio est aussi importante que la télévision. »
\item Infériorité : \eng{not as … as} (ou \eng{less … than}). \eng{My phone is not as fast as yours.} « Mon téléphone est moins rapide que le tien. »
\item Supériorité : \eng{comparatif + than}. \eng{The internet is faster than the post.} « Internet est plus rapide que la poste. »
\item Proportion : \eng{the + comparatif …, the + comparatif …}. \eng{The more you read, the better you write.} « Plus tu lis, mieux tu écris. »
\end{itemize}
\end{propriete}

\begin{attention}[Erreurs fréquentes]
\begin{itemize}
\item Jamais \eng{more taller} : un seul marqueur de comparaison (\eng{taller} ou \eng{more interesting}).
\item Après un comparatif, c'est \eng{than}, pas \eng{that} ni \eng{as} : \eng{better than me}.
\item Le superlatif exige \eng{the} : \eng{the best student}, jamais \eng{a best student}.
\item \eng{good} n'a pas de forme régulière : jamais \eng{gooder} ou \eng{the goodest}.
\item Le superlatif se construit avec \eng{in} pour les lieux et les groupes, pas \eng{of} seul : \eng{the highest mountain in Cameroon} « la plus haute montagne du Cameroun ».
\end{itemize}
\end{attention}

\begin{exoresolu}[Exercice résolu : relatifs, modaux, comparaison]
Énoncé — 1. Combine : \eng{The girl won the prize. She wrote the best article.} 2. Complète avec \eng{mustn't} ou \eng{don't have to} : \eng{You … bring a camera; photos are forbidden.} / \eng{You … bring a dictionary; we have some here.} 3. Mets l'adjectif à la forme correcte : \eng{This newspaper is (good) than that one.} / \eng{Mount Cameroon is (high) mountain in the country.}
\tcblower
Solution — 1. \eng{The girl who wrote the best article won the prize.} (relative restrictive, antécédent humain : \eng{who}). 2. \eng{You mustn't bring a camera; photos are forbidden} (interdiction) ; \eng{You don't have to bring a dictionary} (absence d'obligation). 3. \eng{This newspaper is better than that one} (irrégulier : \eng{good, better}) ; \eng{Mount Cameroon is the highest mountain in the country} (superlatif avec \eng{the}).
\end{exoresolu}

\begin{exercice}[À vous de jouer]
\begin{enumerate}
\item Reliez : \eng{This is the village. I was born there.} / \eng{The reporter interviewed me. He works in Bamenda.}
\item Choisissez : \eng{You (should / should have) studied yesterday; now it is too late.}
\item Traduisez : « Plus tu écoutes la radio, mieux tu comprends l'anglais. »
\item Corrigez : \eng{He musts to go now.} / \eng{My sister is more tall than me.}
\item Complétez avec le relatif qui convient : \eng{The man … car was stolen reported it to the police.} / \eng{The article … I read yesterday was about fake news.}
\end{enumerate}
\end{exercice}

\begin{corrige}[Exercice « À vous de jouer »]
\begin{enumerate}
\item \eng{This is the village where I was born.} / \eng{The reporter who interviewed me works in Bamenda.}
\item \eng{You should have studied yesterday; now it is too late.} (reproche sur le passé).
\item \eng{The more you listen to the radio, the better you understand English.}
\item \eng{He must go now.} (pas de -s, pas de \eng{to}) ; \eng{My sister is taller than me.} (adjectif court : \emph{-er}, jamais \eng{more tall}).
\item \eng{The man whose car was stolen reported it to the police.} (possession : \eng{whose}) ; \eng{The article (which/that) I read yesterday was about fake news.} (relatif complément d'objet : omissible).
\end{enumerate}
\end{corrige}

\begin{aretenir}[L'essentiel de la section]
\begin{itemize}
\item \eng{who} = personnes, \eng{which} = choses, \eng{that} = restrictif seulement, \eng{whose} = possession, \eng{where} = lieu, \eng{when} = temps ; \eng{what} = « ce que/ce qui », sans antécédent.
\item Relative non restrictive : virgules obligatoires, jamais \eng{that}.
\item Modaux : base verbale sans \eng{to}, pas de -s ; \eng{mustn't} = interdit, \eng{don't have to} = pas nécessaire ; \eng{should/could/must have + p.p.} pour parler du passé.
\item Comparatif : \emph{-er}/\eng{more} + \eng{than} ; superlatif : \eng{the} + \emph{-est}/\eng{most} ; irréguliers : \eng{better/best}, \eng{worse/worst}, \eng{farther/farthest}, \eng{further/furthest}.
\end{itemize}
\end{aretenir}

\coursec{Gérondif/infinitif et pièges des francophones}

Dans la section précédente, nous avons révisé les propositions relatives, les verbes modaux et la comparaison. Nous abordons maintenant un point où les francophones trébuchent presque toujours : que mettre après un verbe anglais, un \cle{gérondif} en \eng{-ing} ou un \cle{infinitif} avec \eng{to} ? Nous terminerons par les pièges de vocabulaire les plus fréquents : faux amis et mots indénombrables.

\courssub{Observation : un texte pour découvrir}

\begin{textebox}[A letter from a penfriend]
Dear Marie-Claire,

Thank you for your last letter. I really enjoy reading about life in Yaoundé. I have decided to visit Cameroon next year, and I hope to spend at least two weeks there. I can't afford to stay in expensive hotels, but I don't mind sleeping in simple guesthouses. I am very interested in learning about your culture, and I am quite good at cooking, so I promise to prepare a French meal for your family!

Last month, I stopped eating meat because I want to be healthier. My mother suggested trying your famous ndolé, so please send me the recipe. I remember tasting it when your cousin visited us, and I always remember to buy plantains when I go to the African market here.

I keep practising my French every day, but I must avoid making too many mistakes when I speak. I look forward to meeting you soon. By the way, the news is good: my brother has finished writing his book!

Your friend,
Emily
\end{textebox}

\textbf{Glossaire :} \emph{penfriend} (n., correspondant(e)) ; \emph{guesthouse} (n., maison d'hôtes) ; \emph{recipe} (n., recette) ; \emph{plantain} (n., plantain) ; \emph{by the way} (loc. adv., à propos).

\textbf{Questions de compréhension :}
\begin{enumerate}
\item What has Emily decided to do next year?
\item Why did she stop eating meat?
\item What does she look forward to?
\end{enumerate}

Relevez dans le texte : \eng{enjoy reading}, \eng{decided to visit}, \eng{hope to spend}, \eng{can't afford to stay}, \eng{mind sleeping}, \eng{interested in learning}, \eng{good at cooking}, \eng{stopped eating}, \eng{suggested trying}, \eng{keep practising}, \eng{avoid making}, \eng{look forward to meeting}. Deux constructions apparaissent : verbe + \cle{-ing} et verbe + \cle{to + base verbale}. Comment savoir laquelle choisir ?

\courssub{Verbes suivis du gérondif (-ing)}

\begin{propriete}[Règle : verbe + -ing]
Certains verbes anglais sont \textbf{toujours} suivis du gérondif (base verbale + \eng{-ing}), jamais de l'infinitif. Les plus fréquents au programme de Terminale :
\begin{itemize}
\item \eng{enjoy} (apprécier), \eng{avoid} (éviter), \eng{mind} (déranger, gêner), \eng{suggest} (suggérer), \eng{finish} (finir), \eng{practise} (s'entraîner), \eng{keep} (continuer à), \eng{admit} (admettre), \eng{consider} (envisager), \eng{imagine} (imaginer), \eng{miss} (rater, regretter), \eng{risk} (risquer).
\end{itemize}
Forme : \textbf{verbe + -ing}. Exemple : \eng{I enjoy reading novels.} « J'aime lire des romans. »
\end{propriete}

\begin{exemplebox}[Exemples traduits]
\eng{She avoids travelling at night.} « Elle évite de voyager la nuit. »

\eng{Would you mind opening the window?} « Cela vous dérangerait-il d'ouvrir la fenêtre ? »

\eng{He suggested visiting the Limbe Wildlife Centre.} « Il a suggéré de visiter le centre de faune de Limbé. »

\eng{Have you finished doing your homework?} « As-tu fini de faire tes devoirs ? »

\eng{They keep making the same mistakes.} « Ils continuent à faire les mêmes erreurs. »
\end{exemplebox}

\begin{attention}[Piège : suggest]
\eng{Suggest} n'est jamais suivi de \eng{to + infinitif} ni d'un objet + infinitif. On ne dit pas \emph{He suggested me to go}. Dites : \eng{He suggested going.} « Il a suggéré d'y aller. » ou \eng{He suggested that I (should) go.} « Il a suggéré que j'y aille. »
\end{attention}

\courssub{Verbes suivis de l'infinitif (to + base verbale)}

\begin{propriete}[Règle : verbe + to + BV]
D'autres verbes sont \textbf{toujours} suivis de l'infinitif avec \eng{to} :
\begin{itemize}
\item \eng{want} (vouloir), \eng{decide} (décider), \eng{hope} (espérer), \eng{promise} (promettre), \eng{refuse} (refuser), \eng{afford} (avoir les moyens de), \eng{plan} (prévoir), \eng{agree} (accepter), \eng{learn} (apprendre à), \eng{offer} (proposer), \eng{expect} (s'attendre à), \eng{fail} (ne pas réussir à).
\end{itemize}
Forme : \textbf{verbe + to + base verbale}. Exemple : \eng{We decided to leave early.} « Nous avons décidé de partir tôt. »
\end{propriete}

\begin{exemplebox}[Exemples traduits]
\eng{I can't afford to buy a new phone.} « Je n'ai pas les moyens d'acheter un nouveau téléphone. »

\eng{She promised to help me with my essay.} « Elle a promis de m'aider avec ma rédaction. »

\eng{They refused to answer the journalist's questions.} « Ils ont refusé de répondre aux questions du journaliste. »

\eng{We plan to spend the holidays in Buea.} « Nous prévoyons de passer les vacances à Buéa. »
\end{exemplebox}

\begin{center}
\begin{tabularx}{\linewidth}{|X|X|X|}
\hline
\rowcolor{popblueL} Verb + -ing & Verb + to-infinitive & Remarque \\
\hline
enjoy reading & want to read & aucune logique française : mémoriser les listes \\
\hline
avoid making & decide to make & \eng{avoid} jamais suivi de \eng{to} \\
\hline
mind waiting & hope to wait & \eng{hope} jamais suivi de \eng{-ing} \\
\hline
suggest going & promise to go & \eng{suggest} + that + sujet + BV aussi possible \\
\hline
finish writing & refuse to write & — \\
\hline
practise speaking & plan to speak & — \\
\hline
keep trying & afford to try & \eng{afford} presque toujours avec \eng{can/could} \\
\hline
\end{tabularx}
\end{center}

\begin{popfigure}
\begin{tikzpicture}[
  boxA/.style={rectangle, rounded corners, draw=popblue, fill=popblueL, align=center, text width=5.2cm, inner sep=6pt},
  boxB/.style={rectangle, rounded corners, draw=poporange, fill=poporangeL, align=center, text width=5.2cm, inner sep=6pt},
  boxC/.style={rectangle, rounded corners, draw=popgreen, fill=popgreenL, align=center, text width=11.5cm, inner sep=6pt},
  titre/.style={font=\bfseries}]
\node[titre, text=popblue] (ta) at (-3.2,0) {Verb + -ing};
\node[boxA, below=2pt of ta, align=center] (a){enjoy, avoid, mind, suggest, finish, practise, keep\\[3pt]
\emph{I enjoy reading.}\\ \emph{She finished cooking.}};
\node[titre, text=poporange] (tb) at (3.2,0) {Verb + to-infinitive};
\node[boxB, below=2pt of tb, align=center] (b){want, decide, hope, promise, refuse, afford, plan\\[3pt]
\emph{We decided to leave.}\\ \emph{He promised to call.}};
\node[titre, text=popgreen] (tc) at (0,-4.6) {Uncountable nouns : no plural, no \emph{a/an}};
\node[boxC, below=2pt of tc, align=center] (c){information, advice, news, homework, luggage, furniture, money, bread\\
\emph{a piece of information} --- \emph{some advice} --- \emph{The news is good.}};
\end{tikzpicture}
\legende{Carte mentale : gérondif, infinitif et mots indénombrables.}
\end{popfigure}

\courssub{Verbes à double construction : le sens change}

\begin{propriete}[stop, remember : deux constructions, deux sens]
Certains verbes acceptent les deux constructions, mais avec un \textbf{changement de sens} :
\begin{itemize}
\item \eng{stop + -ing} = cesser une activité. \eng{He stopped smoking.} « Il a arrêté de fumer. »
\item \eng{stop + to-infinitive} = s'arrêter \textbf{pour} faire autre chose. \eng{He stopped to smoke.} « Il s'est arrêté pour fumer. »
\item \eng{remember + -ing} = se souvenir d'un fait passé. \eng{I remember meeting her in Douala.} « Je me souviens de l'avoir rencontrée à Douala. »
\item \eng{remember + to-infinitive} = penser à faire quelque chose (ne pas oublier). \eng{Remember to lock the door.} « N'oublie pas de fermer la porte. »
\end{itemize}
Même logique avec \eng{forget} : \eng{I'll never forget seeing Mount Cameroon.} « Je n'oublierai jamais d'avoir vu le mont Cameroun. » / \eng{Don't forget to bring your dictionary.} « N'oublie pas d'apporter ton dictionnaire. »
\end{propriete}

\begin{propriete}[try : deux constructions, deux nuances]
\begin{itemize}
\item \eng{try + to-infinitive} = essayer de faire quelque chose (tenter). \eng{Try to open the window.} « Essaie d'ouvrir la fenêtre. »
\item \eng{try + -ing} = essayer quelque chose comme expérience, pour voir le résultat. \eng{Try adding more salt.} « Essaie d'ajouter du sel (pour voir). »
\end{itemize}
Enfin, \eng{like}, \eng{love} et \eng{hate} acceptent les deux constructions sans grande différence : \eng{I like reading.} = \eng{I like to read.} « J'aime lire. »
\end{propriete}

\begin{exemplebox}[La paire classique à connaître par cœur]
\eng{He stopped smoking last year.} « Il a arrêté de fumer l'année dernière. »

\eng{He stopped to smoke a cigarette.} « Il s'est arrêté pour fumer une cigarette. »
\end{exemplebox}

\courssub{Préposition + gérondif}

\begin{propriete}[Après toute préposition, le verbe est en -ing]
Règle absolue : en anglais, tout verbe placé \textbf{après une préposition} (\eng{in, at, of, for, without, before, after, by}...) prend la forme en \eng{-ing}.
\begin{itemize}
\item \eng{I am interested in learning English.} « Je suis intéressé par l'apprentissage de l'anglais. »
\item \eng{She is good at speaking in public.} « Elle est douée pour parler en public. »
\item \eng{He left without saying goodbye.} « Il est parti sans dire au revoir. »
\item \eng{Thank you for coming.} « Merci d'être venu. »
\end{itemize}
\end{propriete}

\begin{attention}[Le piège n° 1 du Bac : look forward to]
Dans \eng{look forward to} (avoir hâte de), le mot \eng{to} est une \textbf{préposition}, pas la marque de l'infinitif. Le verbe qui suit est donc en \eng{-ing} :

\eng{I look forward to meeting you.} « J'ai hâte de vous rencontrer. »

\emph{I look forward to meet you} est \textbf{faux}. Même piège avec \eng{to be used to + -ing} (être habitué à) : \eng{I am used to getting up early.} « J'ai l'habitude de me lever tôt. »
\end{attention}

\courssub{Faux amis et calques fréquents}

\begin{attention}[Faux amis : attention aux traductions automatiques]
\begin{itemize}
\item \eng{actually} = \textbf{en fait, en réalité} (et non « actuellement »). \eng{Actually, I live in Bamenda.} « En fait, j'habite à Bamenda. » Pour « actuellement », dites \eng{currently, at present}.
\item \eng{to assist} = \textbf{aider} (et non « assister à » au sens d'être présent). \eng{She assisted the old man.} « Elle a aidé le vieil homme. » Pour « assister à une réunion », dites \eng{to attend a meeting}.
\item \eng{sensible} = \textbf{raisonnable, sensé} (et non « sensible »). \eng{That is a sensible decision.} « C'est une décision raisonnable. » Pour « sensible », dites \eng{sensitive}.
\item \eng{to demand} = exiger ; pour « demander », dites \eng{to ask}.
\item \eng{a library} = une bibliothèque ; une « librairie » se dit \eng{a bookshop}.
\end{itemize}
\end{attention}

\begin{attention}[Erreurs typiques des francophones : mots indénombrables]
\eng{information}, \eng{advice}, \eng{news}, \eng{homework}, \eng{luggage} sont \textbf{indénombrables} : jamais de \eng{-s}, jamais de \eng{a/an}.
\begin{itemize}
\item \emph{an information}, \emph{informations} : impossible. Dites \eng{a piece of information} ou \eng{some information}.
\item \emph{advices} : impossible. Dites \eng{some advice} ou \eng{a piece of advice}.
\item \emph{to be agree with} : impossible. Dites \eng{to agree with}. \eng{I agree with you.} « Je suis d'accord avec toi. »
\item \emph{married with} : impossible. Dites \eng{married to}. \eng{She is married to a teacher.} « Elle est mariée à un enseignant. »
\item \eng{news} se termine par \eng{-s} mais est singulier : \eng{The news is good.} « Les nouvelles sont bonnes. »
\end{itemize}
\end{attention}

\begin{savaistu}[Pourquoi « news » a-t-il un -s ?]
Le mot \eng{news} vient de l'anglais tardif du Moyen Âge \emph{newes}, forme plurielle de \emph{new} (« chose nouvelle ») : les « nouvelles » étaient une collection de faits nouveaux. Aujourd'hui, le mot a perdu son sens pluriel et s'accorde au singulier : \eng{The news is good today.} « Les nouvelles sont bonnes aujourd'hui. » Même phénomène avec \eng{mathematics} et \eng{politics} : \eng{Mathematics is difficult.} « Les mathématiques sont difficiles. »
\end{savaistu}

\begin{methode}[Comment éviter les calques du français]
\begin{enumerate}
\item Traduisez mentalement votre phrase française en structure anglaise \textbf{apprise en classe}, pas mot à mot.
\item Demandez-vous : « Ai-je déjà vu cette structure dans un texte anglais authentique ? » Si non, méfiez-vous.
\item Vérifiez les points chauds : verbe + \eng{-ing} ou \eng{to} ? préposition correcte (\eng{married to}, \eng{agree with}, \eng{interested in}) ? mot dénombrable ou non (\eng{information}, \eng{advice}) ?
\item En cas de doute, privilégiez la \textbf{structure simple} : une phrase correcte et simple vaut mieux qu'une phrase compliquée et fautive.
\item Tenez un carnet personnel de vos erreurs fréquentes et relisez-le avant chaque épreuve.
\end{enumerate}
\end{methode}

\begin{exoresolu}[Exercice résolu : choisir la bonne forme]
Complétez avec la forme correcte du verbe entre parenthèses (gérondif ou infinitif) :
\begin{enumerate}
\item I enjoy (listen) to the radio every morning.
\item We decided (visit) our grandmother in Bafoussam.
\item She suggested (watch) the football match together.
\item He stopped (smoke) two years ago.
\item I look forward to (hear) from you soon.
\item They can't afford (buy) a car.
\end{enumerate}
\tcblower
\textbf{Solution détaillée :}
\begin{enumerate}
\item \eng{I enjoy listening to the radio every morning.} — \eng{enjoy} est toujours suivi du gérondif.
\item \eng{We decided to visit our grandmother in Bafoussam.} — \eng{decide} est toujours suivi de l'infinitif.
\item \eng{She suggested watching the football match together.} — \eng{suggest} + gérondif (jamais \eng{to}).
\item \eng{He stopped smoking two years ago.} — il a cessé l'habitude : gérondif. (\eng{stopped to smoke} signifierait qu'il s'est arrêté \emph{pour} fumer.)
\item \eng{I look forward to hearing from you soon.} — après \eng{look forward to}, préposition + \eng{-ing}.
\item \eng{They can't afford to buy a car.} — \eng{afford} + infinitif.
\end{enumerate}
\end{exoresolu}

\begin{exercice}[À vous de jouer]
A. Corrigez les erreurs dans les phrases suivantes :
\begin{enumerate}
\item She gave me some useful advices for the exam.
\item I am agree with your opinion about social media.
\item My sister is married with a doctor from Douala.
\item Actually, more students own smartphones than before. (l'élève veut dire « actuellement »)
\item I look forward to meet the new English teacher.
\end{enumerate}
B. Traduisez en anglais :
\begin{enumerate}
\item J'ai décidé d'apprendre l'anglais parce que j'aime voyager.
\item Il a arrêté de regarder la télévision pour finir ses devoirs.
\item N'oublie pas d'envoyer la lettre ; je me souviens de l'avoir écrite hier.
\end{enumerate}
\end{exercice}

\begin{corrige}[Exercice « À vous de jouer »]
A. 1. \eng{She gave me some useful advice for the exam.} (\eng{advice} indénombrable). 2. \eng{I agree with your opinion about social media.} (\eng{agree} est un verbe, pas un adjectif). 3. \eng{My sister is married to a doctor from Douala.} 4. \eng{Currently, more students own smartphones than before.} (\eng{actually} = en fait). 5. \eng{I look forward to meeting the new English teacher.}

B. 1. \eng{I decided to learn English because I enjoy travelling.} 2. \eng{He stopped watching television to finish his homework.} (il a cessé de regarder : \eng{stopped watching} ; afin de finir : \eng{to finish}). 3. \eng{Don't forget to send the letter; I remember writing it yesterday.}
\end{corrige}

\begin{aretenir}[L'essentiel de la section]
\begin{itemize}
\item Après \eng{enjoy, avoid, mind, suggest, finish, practise, keep} : gérondif (\eng{-ing}).
\item Après \eng{want, decide, hope, promise, refuse, afford, plan} : infinitif (\eng{to} + BV).
\item \eng{stop/remember + -ing} = action passée ou cessée ; \eng{stop/remember + to} = but ou action future.
\item Après toute préposition, y compris le \eng{to} de \eng{look forward to} : \eng{-ing}.
\item \eng{information, advice, news, homework, luggage} : indénombrables, jamais de \eng{-s} ni de \eng{a}.
\item Faux amis : \eng{actually} = en fait ; \eng{to assist} = aider ; \eng{sensible} = raisonnable.
\end{itemize}
\end{aretenir}

\coursec{Diagnostic et carte mentale des structures}

Avant de plonger dans les exercices de type Bac, il est indispensable de faire le point sur les structures grammaticales révisées tout au long de l'année. Cette carte mentale synthétise les \cle{neuf piliers} de la grammaire de Terminale A. Utilisez-la comme checklist mentale : si une structure vous semble floue, reportez-vous à la leçon correspondante (EN49 à EN57) avant de continuer.

\begin{popfigure}
\begin{tikzpicture}[
  mindmap,
  concept color=popblue,
  level 1/.style={level distance=3.5cm, sibling angle=40},
  level 2/.style={level distance=2.8cm, sibling angle=45},
  every node/.style={concept, execute at begin node=\hskip0pt},
  root concept/.append style={concept color=popblue, font=\large\bfseries, text width=3.5cm, align=center},
  level 1 concept/.append style={font=\normalsize\bfseries, text width=2.8cm, align=center},
  level 2 concept/.append style={font=\small, text width=2.2cm, align=center}
]
\node[root concept, align=center]{Révision\\ Générale\\ (Terminale A)}
  child[concept color=popgreen] {node {Temps verbaux}
    child {node {Présent / Prétérit}}
    child {node {Perfect (Have + V-en)}}
    child {node {Progressif (Be + -ing)}}
    child {node {Futur (Will / Be going to)}}
    child {node {Futur antérieur}}}
  child[concept color=poporange] {node {Voix passive}
    child {node {Be + Participe passé}}
    child {node {Agent (by) / Sans agent}}
    child {node {Tous les temps}}}
  child[concept color=poppurple] {node {Discours rapporté}
    child {node {Backshift (Will$\to$Would...)}}
    child {node {Pronoms / Adverbes}}
    child {node {Questions indirectes (If/Whether)}}
    child {node {Ordres (Tell sb to / Not to)}}}
  child[concept color=poppink] {node {Conditionnels \& Hypothèses}
    child {node {Type 0 (Réalité)}}
    child {node {Type 1 (Probable)}}
    child {node {Type 2 (Irréel présent)}}
    child {node {Type 3 (Irréel passé)}}
    child {node {Wish / If only / As if}}}
  child[concept color=popgold] {node {Propositions relatives}
    child {node {Who / Which / That}}
    child {node {Whose (Possession)}}
    child {node {Whom (Objet formel)}}
    child {node {Ø (Omission objet)}}
    child {node {Préposition + Whom/Which}}}
  child[concept color=popblue] {node {Verbes modaux}
    child {node {Obligation (Must/Have to)}}
    child {node {Capacité (Can/Could)}}
    child {node {Probabilité (May/Might)}}
    child {node {Conseil (Should/Ought to)}}
    child {node {Modaux au passé (Have + V-en)}}}
  child[concept color=popdark] {node {Comparaison}
    child {node {Supériorité / Infériorité}}
    child {node {Égalité (As...as)}}
    child {node {Superlatif (The most / -est)}}
    child {node {Irréguliers (Better, Worse, Further)}}}
  child[concept color=popmuted] {node {Gérondif / Infinitif}
    child {node {Verbes + -ing (Enjoy, Avoid)}}
    child {node {Verbes + To (Decide, Want)}}
    child {node {Verbes + Obj + To (Advise)}}
    child {node {Used to vs Be used to}}
    child {node {Préposition + -ing}}};
\end{tikzpicture}
\legende{Carte mentale des 9 structures clés du programme de Terminale A.}
\end{popfigure}

\coursec{Entraînement type Bac et production guidée}

Après avoir révisé les structures essentielles — temps verbaux, passif, discours rapporté, conditionnels, relatifs, modaux, comparaison, gérondif et infinitif — nous mettons maintenant ces connaissances à l'épreuve dans des exercices de type baccalauréat. L'objectif est double : acquérir la rapidité et la précision nécessaires à l'examen, et savoir mobiliser ces structures dans une production écrite personnelle et cohérente.

\courssub{Exercice de transformation de phrases}

\begin{methode}[Réussir les transformations au Bac]
\begin{enumerate}
\item \textbf{Lisez la consigne} : identifiez la structure cible (passif, discours rapporté, conditionnel, relatif, etc.).
\item \textbf{Repérez le temps} du verbe principal dans la phrase de départ.
\item \textbf{Transformez} en appliquant la règle correspondante sans changer le sens.
\item \textbf{Relisez} en vérifiant : accord sujet-verbe, orthographe, ordre des mots, conservation du sens.
\end{enumerate}
\end{methode}

\begin{exercice}[Transformations — 6 phrases]
Transformez chaque phrase selon l'indication entre parenthèses. Conservez le sens exact.
\begin{enumerate}
\item \eng{They broadcast the news in three languages.} (Passif, même temps)
\item \eng{``I will write an article about social media,'' she said.} (Discours rapporté)
\item \eng{If you don't check your sources, you will spread fake news.} (Conditionnel premier type $\to$ deuxième type)
\item \eng{The journalist who interviewed the minister is my cousin.} (Relatif : remplacez par \eng{whom} si possible, sinon conservez)
\item \eng{Although he was tired, he finished editing the video.} (Concession : commencez par \eng{Despite})
\item \eng{``Don't share unverified information,'' the teacher told us.} (Discours rapporté : ordre / infinitif)
\end{enumerate}
\end{exercice}

\begin{exoresolu}[Corrigé détaillé]
\textbf{Reportez-vous à l'exercice ci-dessus.}
\tcblower
\begin{enumerate}
\item \eng{The news is broadcast in three languages.} (Passif présent simple : \eng{is} + participe passé ; \eng{broadcast} est invariable : \eng{broadcast - broadcast - broadcast}).
\item \eng{She said (that) she would write an article about social media.} (\eng{will} $\to$ \eng{would} = \cle{backshift} ; pronom \eng{I} $\to$ \eng{she} ; suppression des guillemets).
\item \eng{If you didn't check your sources, you would spread fake news.} (Type 2 : \eng{if} + prétérit \eng{didn't check}, \eng{would} + base verbale \eng{spread} — hypothèse présente/future irréelle).
\item \eng{The journalist who interviewed the minister is my cousin.} (\textbf{Attention} : \eng{who} est \textbf{sujet} de \eng{interviewed} (le journaliste a interviewé le ministre). \eng{Whom} est impossible car c'est un pronom objet. On conserve \eng{who} ou \eng{that}).
\item \eng{Despite being tired, he finished editing the video.} (\eng{Despite} + gérondif \eng{being} ; transformation \eng{Although} + clause $\to$ \eng{Despite} + \eng{-ing}).
\item \eng{The teacher told us not to share unverified information.} (\eng{told} + objet \eng{us} + \eng{not to} + infinitif ; ordre négatif rapporté).
\end{enumerate}
\end{exoresolu}

\begin{attention}[Pièges fréquents en transformation]
\begin{itemize}
\item \textbf{Passif} : oublie l'auxiliaire \eng{be} au bon temps (\eng{The news broadcast} \texttimes $\to$ \eng{The news \textbf{is} broadcast} \checkmark).
\item \textbf{Discours rapporté} : oublie le \cle{backshift} des temps (\eng{will} $\to$ \eng{would}, \eng{can} $\to$ \eng{could}, présent $\to$ prétérit) ou conserve les guillemets.
\item \textbf{Conditionnel} : écrit \eng{If I would have...} au lieu de \eng{If I had...} (calque du français \og si j'aurais \fg).
\item \textbf{Relatifs} : utilise \eng{who} pour un objet (\eng{The man who I saw} \texttimes $\to$ \eng{The man \textbf{whom/that/Ø} I saw} \checkmark) ou \eng{whom} pour un sujet (\eng{The man whom called} \texttimes).
\item \textbf{Orthographe} : \eng{broadcasted} \texttimes (verbe irrégulier : \eng{broadcast - broadcast - broadcast}) ; \eng{writen} \texttimes $\to$ \eng{written}.
\end{itemize}
\end{attention}

\courssub{Exercice de correction d'erreurs}

\begin{exercice}[Correction — 8 phrases]
Chaque phrase contient \textbf{une seule erreur} typique de francophone. Corrigez-la et expliquez la règle.
\begin{enumerate}
\item \eng{If I would know the answer, I will tell you.}
\item \eng{The article has been wrote by a famous journalist.}
\item \eng{She asked me where did I buy this newspaper.}
\item \eng{He is the only one who can speaks English fluently.}
\item \eng{Despite of the rain, the match continued.}
\item \eng{I am used to read the news on my phone every morning.}
\item \eng{The informations given by the radio were very useful.}
\item \eng{He suggested me to apply for the internship.}
\end{enumerate}
\end{exercice}

\begin{exoresolu}[Corrigé et explications]
\textbf{Reportez-vous à l'exercice ci-dessus.}
\tcblower
\begin{enumerate}
\item \eng{If I \textbf{knew} the answer, I \textbf{would} tell you.} (Type 2 : \eng{if} + prétérit, \eng{would} + BV ; \textbf{pas de \eng{would} après \eng{if}} — calque du français \og si j'aurais \fg).
\item \eng{The article has been \textbf{written} by a famous journalist.} (Participe passé irrégulier : \eng{write - wrote - \textbf{written}}).
\item \eng{She asked me \textbf{where I bought} this newspaper.} (Discours rapporté : ordre sujet-verbe affirmatif, pas d'auxiliaire \eng{did} ; \eng{where I bought}).
\item \eng{He is the only one who can \textbf{speak} English fluently.} (Modaux : \eng{can} + base verbale \textbf{sans \eng{to}}, sans \eng{-s} à la 3e pers. du singulier).
\item \eng{\textbf{Despite} the rain, the match continued.} (\eng{Despite} préposition directe, sans \eng{of} ; \eng{In spite \textbf{of}} prend \eng{of}, pas \eng{Despite}).
\item \eng{I am used to \textbf{reading} the news on my phone every morning.} (\eng{be used to} + gérondif \eng{-ing} = habitude actuelle ; \eng{used to} + BV = habitude passée révolue).
\item \eng{The \textbf{information} given by the radio \textbf{was} very useful.} (\eng{Information} = indénombrable / \cle{uncountable}, jamais de \eng{-s} ; accord verbe au singulier \eng{was}).
\item \eng{He \textbf{suggested (that) I apply} for the internship.} (\eng{Suggest} + \eng{that}-clause au \cle{subjonctif mandatif} (base verbale \eng{apply}) ou \eng{should apply} ; \textbf{jamais} \eng{suggest someone to do}).
\end{enumerate}
\end{exoresolu}

\begin{center}
\begin{tabularx}{\linewidth}{|X|X|X|}
\hline
\rowcolor{popblueL} \textbf{Erreur francophone} & \textbf{Forme correcte} & \textbf{Règle / Explication} \\
\hline
\eng{If I would know...} & \eng{If I knew...} & Pas de \eng{would} dans la proposition \eng{if} (type 2) \\
\hline
\eng{has been wrote} & \eng{has been written} & Participe passé irrégulier (\eng{write/wrote/written}) \\
\hline
\eng{where did I buy} & \eng{where I bought} & Discours rapporté : ordre affirmatif, pas d'auxiliaire \eng{do} \\
\hline
\eng{can speaks} & \eng{can speak} & Modal + base verbale sans \eng{-s} \\
\hline
\eng{Despite of} & \eng{Despite} / \eng{In spite of} & \eng{Despite} préposition directe ; \eng{In spite of} = trois mots \\
\hline
\eng{used to read} & \eng{used to reading} & \eng{be used to} + \eng{-ing} (habitude) vs \eng{used to} + BV (passé) \\
\hline
\eng{informations} & \eng{information} & Nom indénombrable (\eng{uncountable}) \\
\hline
\eng{suggested me to apply} & \eng{suggested (that) I apply} & \eng{Suggest} + \eng{that}-clause (subjonctif) ; pas d'infinitif \\
\hline
\end{tabularx}
\end{center}

\courssub{Exercice à trous : les médias au Cameroun}

\begin{textebox}[Texte à compléter — Médias et langues locales]
Complétez le texte en conjuguant les verbes entre parenthèses au temps approprié (actif ou passif). Attention aux indices temporels et au sens.

Since the new community radio station \_\_\_\_\_\_\_\_\_\_ (open) in Buea last year, many programmes \_\_\_\_\_\_\_\_\_\_ (broadcast) in local languages such as Bakweri and Pidgin. Young presenters \_\_\_\_\_\_\_\_\_\_ (train) by experienced journalists before they \_\_\_\_\_\_\_\_\_\_ (go) on air. If more young people \_\_\_\_\_\_\_\_\_\_ (encourage) to study journalism, the quality of local information \_\_\_\_\_\_\_\_\_\_ (improve) significantly. The station manager said that the next project \_\_\_\_\_\_\_\_\_\_ (be) a TV channel for youth. He added that they \_\_\_\_\_\_\_\_\_\_ (look) for sponsors at the moment. By the end of next year, the station \_\_\_\_\_\_\_\_\_\_ (reach) over 50,000 listeners in the region.
\end{textebox}

\begin{exoresolu}[Corrigé du texte à trous]
\textbf{Reportez-vous au texte ci-dessus.}
\tcblower
\begin{enumerate}
\item \eng{opened} (prétérit : \eng{last year} $\to$ action passée datée).
\item \eng{have been broadcast} / \eng{are broadcast} (présent perfect passif / présent passif : conséquence présente d'une action passée / habitude actuelle).
\item \eng{are trained} / \eng{were trained} (passif présent / prétérit : processus actuel ou passé selon perspective).
\item \eng{go} (présent simple après \eng{before} + présent pour futur ; ou \eng{went} si récit au passé — ici \eng{are trained} présent $\to$ \eng{go}).
\item \eng{were encouraged} (conditionnel type 2 : \eng{if} + prétérit passif $\to$ \eng{would improve}).
\item \eng{would improve} (conséquence type 2 : \eng{would} + base verbale).
\item \eng{would be} (discours rapporté : \eng{will} $\to$ \eng{would} = \cle{backshift}).
\item \eng{were looking} (discours rapporté : \eng{are looking} $\to$ \eng{were looking} ; \eng{at the moment} $\to$ \eng{then} mais temps changé).
\item \eng{will have reached} (futur antérieur : \eng{By the end of next year} $\to$ action achevée à une date future).
\end{enumerate}
\end{exoresolu}

\begin{attention}[Vigilance à l'examen]
Une transformation grammaticalement juste mais entachée d'une \textbf{faute d'orthographe} (\eng{broadcasted}, \eng{informations}, \eng{writen}) ou d'un \textbf{défaut d'accord} (\eng{The news are}, \eng{He suggest}) perd des points. \textbf{Relisez-vous systématiquement} : accords, orthographe des irréguliers, majuscules (\eng{I}, noms propres, jours, mois, nationalités).
\end{attention}

\courssub{QCM de grammaire — Entraînement rapidité}

\begin{exercice}[Choisissez la bonne réponse (A, B, C ou D)]
\begin{enumerate}
\item The reporter \_\_\_\_\_\_\_\_ the interview when the power went out.
  \begin{itemize}
  \item[A)] was doing \item[B)] did \item[C)] has done \item[D)] has been doing
  \end{itemize}
\item If the government \_\_\_\_\_\_\_\_ press freedom, journalists would work more freely.
  \begin{itemize}
  \item[A)] respects \item[B)] respected \item[C)] will respect \item[D)] would respect
  \end{itemize}
\item The article \_\_\_\_\_\_\_\_ by the editor before publication yesterday.
  \begin{itemize}
  \item[A)] was checked \item[B)] has been checked \item[C)] is checked \item[D)] checks
  \end{itemize}
\item She asked me \_\_\_\_\_\_\_\_ I had seen the documentary.
  \begin{itemize}
  \item[A)] that \item[B)] if \item[C)] whether \item[D)] B et C sont corrects
  \end{itemize}
\item \_\_\_\_\_\_\_\_ his lack of experience, he got the job as a news anchor.
  \begin{itemize}
  \item[A)] Despite \item[B)] Although \item[C)] In spite \item[D)] Even though
  \end{itemize}
\item I look forward to \_\_\_\_\_\_\_\_ your article in the next issue.
  \begin{itemize}
  \item[A)] read \item[B)] reading \item[C)] to read \item[D)] be reading
  \end{itemize}
\item The journalist \_\_\_\_\_\_\_\_ article won the prize works for a daily newspaper.
  \begin{itemize}
  \item[A)] who \item[B)] which \item[C)] whose \item[D)] whom
  \end{itemize}
\item He denied \_\_\_\_\_\_\_\_ the confidential document.
  \begin{itemize}
  \item[A)] to leak \item[B)] leaking \item[C)] leak \item[D)] have leaked
  \end{itemize}
\end{enumerate}
\end{exercice}

\begin{corrige}[QCM]
\begin{enumerate}
\item \textbf{A} : \eng{was doing} (passé progressif = action en cours interrompue par \eng{when} + prétérit).
\item \textbf{B} : \eng{respected} (conditionnel type 2 : \eng{if} + prétérit).
\item \textbf{A} : \eng{was checked} (passif prétérit : \eng{yesterday} $\to$ passé daté).
\item \textbf{D} : \eng{if} ou \eng{whether} pour question indirecte fermée (oui/non) ; \eng{that} introduit une déclaration, pas une question.
\item \textbf{A} : \eng{Despite} + nom/GN (\eng{his lack of experience}) ; \eng{Although} + clause (sujet + verbe).
\item \textbf{B} : \eng{look forward to} + gérondif (\eng{to} = préposition ici, pas marqueur d'infinitif).
\item \textbf{C} : \eng{whose} (possessif relatif, antécédent personne, \eng{article} = possédé).
\item \textbf{B} : \eng{deny} + gérondif (\eng{deny doing sth}).
\end{enumerate}
\end{corrige}

\courssub{Production guidée : \eng{The role of media in my country}}

\begin{methode}[Rédiger un paragraphe mobilisant 4 structures révisées]
\begin{enumerate}
\item \textbf{Planifiez} : idée principale + 2-3 arguments + exemples camerounais concrets (radio communautaire, CRTV, réseaux sociaux, langues locales, \eng{fake news}).
\item \textbf{Cochez les structures obligatoires} (au minimum 4) :
  \begin{itemize}
  \item Un passif (ex. \eng{Programmes are broadcast in local languages.})
  \item Un conditionnel (type 2 ou 3 : \eng{If more journalists were trained...})
  \item Un relatif (ex. \eng{Journalists who investigate corruption...})
  \item Un discours rapporté ou un modal de nuance (ex. \eng{The minister stated that...} / \eng{Media must verify...})
  \end{itemize}
\item \textbf{Rédigez} 8-10 lignes (environ 100-120 mots) en anglais correct, phrases complètes, connecteurs logiques (\eng{Furthermore, However, For instance, Consequently}).
\item \textbf{Relisez} avec la grille d'auto-évaluation ci-dessous.
\end{enumerate}
\end{methode}

\begin{exemplebox}[Modèle de paragraphe (≈ 110 mots)]
\eng{The media play a crucial role in Cameroon by informing citizens and promoting cultural diversity. Programmes \textbf{are broadcast} in over 200 local languages, which allows rural communities to access vital information. \textbf{If more young journalists were trained} in investigative reporting, the quality of news \textbf{would improve} significantly. Journalists \textbf{who expose corruption} often face pressure, yet their work is essential for democracy. The Minister of Communication recently \textbf{stated that} new laws \textbf{would protect} press freedom. \textbf{However}, social media \textbf{must be used} responsibly to avoid the spread of fake news. Ultimately, a free and professional media landscape benefits the whole nation.}
\end{exemplebox}

\begin{center}
\begin{tabularx}{\linewidth}{|X|X|}
\hline
\rowcolor{popblueL} \textbf{Structure mobilisée} & \textbf{Exemple dans le modèle} \\
\hline
Passif (présent) & \eng{Programmes are broadcast} \\
\hline
Conditionnel type 2 & \eng{If more young journalists were trained... would improve} \\
\hline
Relatif (sujet) & \eng{Journalists who expose corruption} \\
\hline
Discours rapporté & \eng{stated that new laws would protect} \\
\hline
Modal + passif & \eng{social media must be used} \\
\hline
\end{tabularx}
\end{center}

\courssub{Auto-évaluation — Grille de vérification finale}

Avant de rendre votre copie, cochez chaque critère :

\begin{center}
\begin{tabular}{|l|c|}
\hline
\rowcolor{popblueL} \textbf{Critère} & \textbf{Validé ?} \\
\hline
Tous les verbes sont au temps correct (indices temporels respectés) & \square \\
\hline
Accords sujet-verbe corrects (3e pers. sing. \eng{-s}, passif \eng{be} + PP) & \square \\
\hline
Au moins \textbf{un passif} correctement formé & \square \\
\hline
Au moins \textbf{un conditionnel} (type 2 ou 3) correct & \square \\
\hline
Au moins \textbf{un relatif} (\eng{who/which/whose/whom/that/Ø}) correct & \square \\
\hline
Aucun calque du français (\eng{If I would}, \eng{Despite of}, \eng{suggest me to}, \eng{informations}, \eng{can speaks}...) & \square \\
\hline
Orthographe des irréguliers vérifiée (\eng{broadcast, written, spoken, chosen}...) & \square \\
\hline
Majuscules : \eng{I}, noms propres, jours, mois, nationalités & \square \\
\hline
Connecteurs logiques variés (\eng{Furthermore, However, For instance, Consequently}) & \square \\
\hline
\end{tabular}
\end{center}

\begin{savaistu}[Ce que valorisent les correcteurs du Bac]
Les correcteurs du baccalauréat camerounais (et des examens internationaux type GCE A Level) accordent une note significative à la \textbf{variété et à la justesse des structures} dans la production écrite (Composition / Guided Writing). Une copie qui n'utilise que le présent simple et des phrases courtes (\eng{Subject + Verb + Object}) ne peut pas obtenir la note maximale, même si le contenu est pertinent. \textbf{Montrez que vous maîtrisez} : passif, conditionnels, relatifs, discours rapporté, modaux, gérondif/infinitif, connecteurs. La qualité prime sur la quantité : 10 lignes bien structurées valent mieux qu'une page truffée d'erreurs de base.
\end{savaistu}

\begin{popfigure}
\begin{tikzpicture}[
  node distance=1.2cm and 1.5cm,
  decision/.style={diamond, draw=popblue, thick, fill=popblueL, text width=3.2cm, align=center, inner sep=4pt},
  process/.style={rectangle, draw=popgreen, thick, fill=popgreenL, text width=3.2cm, align=center, inner sep=4pt, rounded corners=4pt},
  start/.style={ellipse, draw=poporange, thick, fill=poporangeL, text width=2.5cm, align=center, inner sep=4pt},
  end/.style={ellipse, draw=poppink, thick, fill=poppinkL, text width=2.5cm, align=center, inner sep=4pt},
  arrow/.style={-Stealth, thick, popdark}
]
\node[start] (start) {Ma phrase est-elle correcte ?};
\node[decision, below=of start] (temps) {Le temps est-il justifié par le contexte ?};
\node[decision, below=of temps] (accord) {Accord sujet-verbe OK ? (3e pers. sg, passif be+PP)};
\node[decision, below=of accord] (structure) {Structure demandée respectée ? (passif, conditionnel, relatif, etc.)};
\node[decision, below=of structure] (sens) {Le sens est-il inchangé (transformation) / clair (production) ?};
\node[end, below=of sens] (ok) {\textbf{VALIDÉE} \checkmark};
\node[process, right=3.5cm of temps] (notemps) {Corrigez le temps};
\node[process, right=3.5cm of accord] (noaccord) {Corrigez l'accord};
\node[process, right=3.5cm of structure] (nostruct) {Revoir la règle};
\node[process, right=3.5cm of sens] (nosens) {Reformulez};

\draw[arrow] (start) -- (temps);
\draw[arrow] (temps) -- node[left, font=\small, text=popgreen] {OUI} (accord);
\draw[arrow] (accord) -- node[left, font=\small, text=popgreen] {OUI} (structure);
\draw[arrow] (structure) -- node[left, font=\small, text=popgreen] {OUI} (sens);
\draw[arrow] (sens) -- node[left, font=\small, text=popgreen] {OUI} (ok);

\draw[arrow] (temps) -- node[above, font=\small, text=poporange] {NON} (notemps);
\draw[arrow] (accord) -- node[above, font=\small, text=poporange] {NON} (noaccord);
\draw[arrow] (structure) -- node[above, font=\small, text=poporange] {NON} (nostruct);
\draw[arrow] (sens) -- node[above, font=\small, text=poporange] {NON} (nosens);

\draw[arrow, popmuted] (notemps) |- (temps);
\draw[arrow, popmuted] (noaccord) |- (accord);
\draw[arrow, popmuted] (nostruct) |- (structure);
\draw[arrow, popmuted] (nosens) |- (sens);
\end{tikzpicture}
\legende{Organigramme de vérification finale : appliquez-le à chaque phrase avant de rendre votre copie.}
\end{popfigure}

\begin{exercice}[Entraînement final — Mini-Bac blanc (20 min)]
\begin{enumerate}
\item \textbf{Transformation (4 pts)} : Transformez sans changer le sens.
  \begin{enumerate}
  \item \eng{They have translated the report into Fulfulde.} (Passif)
  \item \eng{``We are launching a new app tomorrow,'' the director announced.} (Discours rapporté)
  \item \eng{Unless the signal improves, we won't hear the news.} (Commencez par \eng{If})
  \item \eng{The editor who rejected my article is very strict.} (Relatif : \eng{whom} si possible)
  \end{enumerate}
\item \textbf{Correction (4 pts)} : Corrigez l'unique erreur.
  \begin{enumerate}
  \item \eng{If I would be you, I will check the source.}
  \item \eng{The news were broadcasted at 8 pm.}
  \item \eng{He advised me to not believe everything on Facebook.}
  \item \eng{She is used to write articles for the school magazine.}
  \end{enumerate}
\item \textbf{Production (12 pts)} : Rédigez 8-10 lignes sur \eng{The impact of social media on youth in Cameroon} en utilisant \textbf{au moins 4 structures} révisées (cochez-les dans la marge).
\end{enumerate}
\end{exercice}

\begin{corrige}[Mini-Bac blanc]
\begin{enumerate}
\item \textbf{Transformation}
  \begin{enumerate}
  \item \eng{The report has been translated into Fulfulde.}
  \item \eng{The director announced (that) they were launching a new app the next day / the following day.}
  \item \eng{If the signal doesn't improve, we won't hear the news.} (ou type 2 : \eng{If the signal didn't improve, we wouldn't hear the news.})
  \item \eng{The editor by whom my article was rejected is very strict.} / \eng{The editor whom my article was rejected by is very strict.} (\textbf{Note} : \eng{reject} prend un objet direct ; au passif \eng{my article was rejected by the editor}, le relatif sur \eng{editor} exige la préposition \eng{by} $\to$ \eng{by whom} ou \eng{whom ... by}).
  \end{enumerate}
\item \textbf{Correction}
  \begin{enumerate}
  \item \eng{If I \textbf{were} you, I \textbf{would} check the source.} (Type 2 : \eng{were} + \eng{would}).
  \item \eng{The news \textbf{was} broadcast at 8 pm.} (\eng{news} singulier ; \eng{broadcast} invariable).
  \item \eng{He advised me \textbf{not to} believe everything on Facebook.} (\eng{advise} + objet + \eng{not to} + BV ; négation de l'infinitif = \eng{not to}).
  \item \eng{She is used to \textbf{writing} articles for the school magazine.} (\eng{be used to} + \eng{-ing} = habitude).
  \end{enumerate}
\item \textbf{Production} : Voir grille d'auto-évaluation. Exemple de structures attendues : passif (\eng{are influenced}), conditionnel (\eng{would be}), relatif (\eng{who spend}), modal passif (\eng{must be regulated}), gérondif (\eng{by spending}), connecteurs (\eng{However, Consequently}).
\end{enumerate}
\end{corrige}

\newpage

\coursec{Activité d'intégration}

\courssub{Objectifs de l'activité}
\begin{itemize}
    \item Mobiliser l'ensemble des structures grammaticales révisées (temps, voix passive, discours rapporté, conditionnels, relatifs, modaux, gérondif/infinitif) dans une production orale et écrite cohérente.
    \item Travailler en collaboration pour concevoir, rédiger et présenter une émission de radio authentique sur le thème des médias au Cameroun.
    \item Développer l'écoute active et l'analyse grammaticale \emph{via} une grille d'évaluation par les pairs.
    \item Renforcer l'autocorrection en identifiant et en expliquant les erreurs fréquentes des francophones (\og erreur de la semaine \fg).
\end{itemize}

\courssub{Situation}
\begin{textebox}[Contexte : Le lancement de \eng{Grammar Clinic} sur \eng{Campus FM}]
\eng{The administration of Government Bilingual High School (GBHS) Molyko, Buea, has just launched \textbf{Campus FM}, a new school radio station managed entirely by students. To celebrate the launch, the English Club organises a special weekly programme called \textbf{``Grammar Clinic''}. Each episode is prepared by a team of four students (Form 5 and Upper Sixth) and broadcast live every Friday at 1:00 p.m. during the lunch break. The target audience is the whole student body, from Form 1 to Upper Sixth, and the language of the show is exclusively English.}

\eng{Your group has been selected to produce the very first episode. The theme is \textbf{``Media and Communication in Cameroon: Challenges and Opportunities''}. The episode must last exactly five minutes and follow a strict format designed to showcase your mastery of advanced grammar structures. The best episode, chosen by the students' vote and validated by the English teacher, will become the pilot for the whole school year.}
\end{textebox}

\begin{textebox}[Script de l'annonce officielle (Teacher's brief)]
\eng{``Good morning, everyone. Welcome to the first production meeting of \textbf{Grammar Clinic}. As you know, \textbf{Campus FM} is now on air. Your mission, should you accept it, is to produce a five-minute magazine show. Here is the mandatory running order:}

\eng{\textbf{1. The Headlines (Direct $\to$ Reported Speech):} One anchor reads three news headlines in direct speech. A ``correspondent'' then reports them using appropriate reporting verbs and backshifting.}

\eng{\textbf{2. The Interview (Reported Questions):} The host interviews a ``special guest'' (a student role-playing a journalist, a blogger, or a media expert). The host must ask at least two questions. After the interview, the host summarises the guest's answers using \emph{reported questions} (e.g., \emph{He asked me whether... / She wanted to know why...}).}

\eng{\textbf{3. The Editorial (Passive Voice + Second Conditional):} One member delivers a short opinion piece (4--5 sentences) on the state of school media. It must contain at least \textbf{three passive forms} (in different tenses) and \textbf{one Second Conditional sentence} starting with \emph{If our school had a radio station, we would...} (or a similar hypothesis).}

\eng{\textbf{4. Error of the Week:} The team presents one classic mistake made by French speakers (e.g., confusion between \emph{make} and \emph{do}, \emph{since} and \emph{for}, a wrong preposition after an adjective, or a false friend). They give the wrong sentence, the correct version, and a 30-second explanation in English.}

\eng{During the other groups' broadcasts, you will fill in a \textbf{Listening Grid} to identify the structures heard. Let's make it professional, accurate, and fun!''}
\end{textebox}

\courssub{Consignes}
\begin{enumerate}
    \item \textbf{Formation des équipes et répartition des rôles (10 min).} Constituez des groupes de 4. Désignez : un \emph{Producer/Timekeeper} (veille au timing et à la cohérence), un \emph{Main Anchor} (titres, transitions), un \emph{Correspondent/Reporter} (discours rapporté, interview), un \emph{Editor/Grammar Expert} (éditorial, erreur de la semaine). Tous participent à l'écriture.
    
    \item \textbf{Rédaction du script complet (35 min).} Rédigez le script \emph{mot pour mot} sur papier ou tablette. Chaque segment doit être numéroté. Vérifiez collectivement chaque phrase selon la \textbf{Grille d'auto-évaluation} fournie dans les Ressources. \textbf{Contrainte :} pas de lecture improvisée ; le script final est remis au professeur avant le passage.
    
    \item \textbf{Répétition et ajustement (15 min).} Répétez à voix basse. Chronométrez : 5 minutes exactement ($\pm$ 15 secondes). Travaillez l'intonation (montante pour les questions fermées directes, descendante pour les questions en \eng{wh-} et les questions rapportées), le rythme et la prononciation des formes faibles : \eng{was} se dit « woz », \eng{had} se dit « hed », \eng{would} se dit « woud » dans le flux de la parole.
    
    \item \textbf{Enregistrement / Diffusion en direct (20 min par groupe).} Les groupes passent tour à tour. Les auditeurs remplissent la \textbf{Grille d'écoute} (voir Ressources). Le professeur note la justesse grammaticale (critère principal), la fluidité, la prononciation et le respect du format.
    
    \item \textbf{Dépouillement et vote (10 min).} Échange des grilles. Discussion collective : quel groupe a respecté toutes les contraintes structurelles ? Quel éditorial était le plus convaincant ? Vote à main levée. Le professeur annonce le groupe gagnant et commente les points forts et les faiblesses globaux.
    
    \item \textbf{Extension écrite (Devoir / 30 min).} Individuellement, réécrivez l'éditorial de votre groupe en développant l'argumentation (150 mots minimum). Intégrez au moins un conditionnel de type 3 (\eng{If we had invested earlier, we would have attracted...}) et deux propositions relatives (dont une avec préposition antéposée : \eng{the studio \textbf{in which} we record...}). Remettez la copie lors du prochain cours.
\end{enumerate}

\courssub{Ressources à mobiliser}
\begin{propriete}[Rappel : Discours rapporté (Reported Speech) -- Backshifting principal]
\begin{center}
\begin{tabularx}{\linewidth}{|X|X|X|}
\hline
\rowcolor{popblueL} \textbf{Direct Speech} & \textbf{Reported Speech} & \textbf{Exemple} \\
\hline
Present Simple & Past Simple & \eng{``I \textbf{work} here.''} $\to$ He said he \textbf{worked} there. \\
Present Continuous & Past Continuous & \eng{``I \textbf{am writing}.''} $\to$ She said she \textbf{was writing}. \\
Present Perfect / Past Simple & Past Perfect & \eng{``I \textbf{have finished}.''} $\to$ He said he \textbf{had finished}. \\
Will & Would & \eng{``I \textbf{will} call.''} $\to$ She said she \textbf{would} call. \\
Can & Could & \eng{``I \textbf{can} speak.''} $\to$ He said he \textbf{could} speak. \\
\hline
\end{tabularx}
\end{center}
\textbf{Questions rapportées :} pas d'auxiliaire \eng{do/does/did} ; ordre sujet + verbe. \eng{``Why \textbf{did} you \textbf{choose} this?''} $\to$ He asked \textbf{why I had chosen} that. \eng{``\textbf{Do} you \textbf{like} it?''} $\to$ She asked \textbf{if/whether I liked} it.

\textbf{Changements de repères :} \eng{now $\to$ then, today $\to$ that day, here $\to$ there, this $\to$ that, tomorrow $\to$ the next/following day, yesterday $\to$ the day before, ago $\to$ before/earlier}.

\textbf{Exception :} pas de backshifting si le verbe de parole est au présent (\eng{He says he \textbf{is} tired}) ou si l'on rapporte une vérité générale encore valable (\eng{The teacher said the earth \textbf{is} round}).
\end{propriete}

\begin{propriete}[Voix passive (Passive Voice) -- Formes attendues dans l'éditorial]
Forme : \eng{be (conjugué au temps voulu) + Past Participle (V3)}.
\begin{center}
\begin{tabular}{|l|l|l|}
\hline
\textbf{Temps} & \textbf{Forme} & \textbf{Exemple contextuel} \\
\hline
Present Simple & \eng{is/are + V3} & \eng{The news \textbf{is broadcast} daily.} \\
Past Simple & \eng{was/were + V3} & \eng{The studio \textbf{was built} last year.} \\
Present Perfect & \eng{has/have been + V3} & \eng{Equipment \textbf{has been donated} by former students.} \\
Modal + \eng{be} + V3 & \eng{must/should/can be + V3} & \eng{Programmes \textbf{must be approved} by the teacher in charge.} \\
Future (\eng{will}) & \eng{will be + V3} & \eng{A new show \textbf{will be launched} soon.} \\
\hline
\end{tabular}
\end{center}
\emph{Attention :} le passif évite de dire \og qui \fg{} fait l'action (agent inconnu, évident ou secondaire). N'utilisez \eng{by + agent} que si l'agent apporte une information réelle. Certains participes sont irréguliers : \eng{broadcast--broadcast--broadcast}, \eng{build--built--built}, \eng{write--wrote--written}, \eng{speak--spoke--spoken}.
\end{propriete}

\begin{propriete}[Conditionnel Type 2 (Second Conditional) -- Hypothèse présente ou future irréelle]
\eng{If + Past Simple (ou Past Continuous), ... would/could/might + Base Verbale (BV)}.

\eng{If our school \textbf{had} a radio station, we \textbf{would produce} great shows.} « Si notre école avait une station de radio, nous produirions de grandes émissions. »

\eng{If I \textbf{were} you, I \textbf{would join} the club.} « Si j'étais toi, je rejoindrais le club. » (\eng{were} pour toutes les personnes, forme soutenue ; \eng{was} est toléré à l'oral pour \eng{I/he/she}.)

\emph{Piège :} ne jamais mettre \eng{would} dans la proposition en \eng{if} : \eng{*If we would have...} est \textbf{FAUX}. Rappel du type 3 (hypothèse passée irréelle) : \eng{If + Past Perfect, ... would have + V3} : \eng{If we had invested earlier, we would have attracted more listeners.}
\end{propriete}

\begin{propriete}[Gérondif vs Infinitif -- Verbes fréquents dans le thème des médias]
\begin{center}
\begin{tabularx}{\linewidth}{|X|X|X|}
\hline
\rowcolor{popblueL} \textbf{Verbe + Gérondif (-ing)} & \textbf{Verbe + Infinitif (to + BV)} & \textbf{Verbe + Objet + Infinitif} \\
\hline
\eng{enjoy, avoid, suggest, consider, keep (on), risk, practise} & \eng{decide, hope, plan, promise, agree, refuse, manage} & \eng{ask, tell, encourage, warn, advise, allow, enable} \\
\eng{The host \textbf{avoided answering}.} & \eng{The guest \textbf{agreed to speak}.} & \eng{The teacher \textbf{encouraged us to participate}.} \\
\hline
\end{tabularx}
\end{center}
\emph{Verbes à double sens :} \eng{stop smoking} « arrêter de fumer » vs \eng{stop to smoke} « s'arrêter pour fumer » ; \eng{remember posting} « se souvenir d'avoir posté » vs \eng{remember to post} « ne pas oublier de poster ». Après une préposition, toujours le gérondif : \eng{interested \textbf{in joining}}, \eng{good \textbf{at speaking}}, \eng{keen \textbf{on broadcasting}}.
\end{propriete}

\begin{attention}[Top 5 des erreurs de francophones à traquer (\og Erreur de la semaine \fg)]
\begin{enumerate}
    \item \textbf{Make vs Do :} \eng{*do a mistake} (FAUX) $\to$ \eng{\textbf{make} a mistake} ; \eng{*do a decision} $\to$ \eng{\textbf{make} a decision}. Mais : \eng{\textbf{do} an interview, \textbf{do} homework, \textbf{do} the dishes, \textbf{do} one's best}.
    \item \textbf{Since vs For :} \eng{*I listen since 2 hours} $\to$ \eng{I \textbf{have listened} \textbf{for} 2 hours} (durée) / \eng{I \textbf{have listened} \textbf{since} 10 a.m.} (point de départ). Avec \eng{since/for}, l'anglais exige le Present Perfect, pas le présent.
    \item \textbf{Prépositions après adjectifs :} \eng{interested \textbf{in}, good \textbf{at}, keen \textbf{on}, aware \textbf{of}, responsible \textbf{for}, satisfied \textbf{with}}. \eng{*interested to} et \eng{*good in} sont \textbf{FAUX}.
    \item \textbf{Faux amis :} \eng{actually} = « en fait » (« actuellement » se dit \eng{currently}) ; \eng{eventually} = « finalement » (« éventuellement » se dit \eng{possibly}) ; \eng{a journalist} = « un journaliste » ; \eng{a library} = « une bibliothèque » (« une librairie » se dit \eng{a bookshop}).
    \item \textbf{Concordance des temps (Sequence of Tenses) :} \eng{*He said he \textbf{will} come} $\to$ \eng{He said he \textbf{would} come}. Ne gardez pas \eng{will} après un verbe de parole au passé.
\end{enumerate}
\end{attention}

\begin{textebox}[Grille d'écoute \eng{(Listener's Grid)} -- À distribuer aux auditeurs]
\eng{\textbf{Group Name: \underline{\hspace{5cm}} \quad Date: \underline{\hspace{3cm}}}}

\eng{\textbf{Segment 1: Headlines (Direct $\to$ Reported)}}

\eng{Reporting verbs used: \underline{\hspace{6cm}} (e.g., \emph{stated, announced, revealed, warned})}

\eng{Backshifting correct? Yes / No \quad Examples heard: \underline{\hspace{5cm}}}

\eng{Time/place markers changed? Yes / No \quad (\emph{now $\to$ then, today $\to$ that day})}

\eng{\textbf{Segment 2: Interview (Reported Questions)}}

\eng{Question 1 (Direct): \underline{\hspace{7cm}}}

\eng{Reported version: \underline{\hspace{7cm}} \quad [Check: no \emph{do/does/did}, subject--verb order]}

\eng{Question 2 (Direct): \underline{\hspace{7cm}}}

\eng{Reported version: \underline{\hspace{7cm}}}

\eng{\textbf{Segment 3: Editorial (Passive + Conditional)}}

\eng{Passive 1 (Tense: \underline{\hspace{2cm}}): \underline{\hspace{7cm}}}

\eng{Passive 2 (Tense: \underline{\hspace{2cm}}): \underline{\hspace{7cm}}}

\eng{Passive 3 (Tense: \underline{\hspace{2cm}}): \underline{\hspace{7cm}}}

\eng{Second Conditional sentence heard: \underline{\hspace{7cm}}}

\eng{\textbf{Segment 4: Error of the Week}}

\eng{Error presented: \underline{\hspace{7cm}}}

\eng{Correction given: \underline{\hspace{7cm}}}

\eng{Explanation clear? Yes / No}

\eng{\textbf{Global Score (Grammar Accuracy /10): \underline{\hspace{1cm}}} \quad \textbf{Best moment:} \underline{\hspace{5cm}}}
\end{textebox}

\courssub{Proposition de production et corrigé commenté}
\begin{exoresolu}[Script modèle -- Groupe \eng{``The Broadcasters''} (GBHS Molyko)]
\textbf{SCRIPT COMPLET (5 min)}

\eng{\textbf{[JINGLE: Upbeat Afrobeat music, 5 sec. Fade under.]}}

\eng{\textbf{ANCHOR (Direct Speech -- Headlines):} ``Good afternoon, Buea! This is \textbf{Campus FM}, your school voice. I'm \textbf{Ngwa}, your anchor. Here are the top stories at 1:00 p.m.}

\eng{\textbf{One:} The Minister of Communication announced yesterday: `We \textbf{will launch} a national media literacy campaign next month.'}

\eng{\textbf{Two:} The Principal stated: `Our new studio \textbf{is equipped} with professional microphones.'}

\eng{\textbf{Three:} A Form 5 student asked: `\textbf{Can} we \textbf{broadcast} in Pidgin English too?'''}

\eng{\textbf{CORRESPONDENT (Reported Speech):} ``Thank you, Ngwa. This is \textbf{Tako}, your correspondent. The Minister announced that they \textbf{would launch} a national media literacy campaign \textbf{the following month}. The Principal stated that the new studio \textbf{was equipped} with professional microphones. A Form 5 student asked \textbf{if they could broadcast} in Pidgin English too. Back to you, Ngwa.''}

\eng{\textbf{[JINGLE: Short transition sting.]}}

\eng{\textbf{HOST (Interview -- Direct Questions):} ``Now, our special segment. Today we welcome \textbf{Mrs Fotso}, a freelance journalist and blogger based in Douala. Welcome, Mrs Fotso!''}

\eng{\textbf{GUEST:} ``Thank you for having me.''}

\eng{\textbf{HOST:} ``My first question is: \textbf{how has social media changed} journalism in Cameroon?''}

\eng{\textbf{GUEST:} ``It has democratised information. Anyone with a smartphone \textbf{can report} live. But verification is harder.''}

\eng{\textbf{HOST:} ``And my second question: \textbf{do you think} school radios \textbf{are} important?''}

\eng{\textbf{GUEST:} ``Absolutely. They train future professionals and give students a voice.''}

\eng{\textbf{HOST (Reported Questions Summary):} ``Thank you, Mrs Fotso. Listeners, let me recap. I asked Mrs Fotso \textbf{how social media had changed} journalism in Cameroon. She explained that it \textbf{had democratised} information but added that verification \textbf{was} harder. Then I asked her \textbf{whether she thought} school radios \textbf{were} important. She confirmed that they \textbf{were} essential for training and student expression.''}

\eng{\textbf{[JINGLE: Solemn, piano background.]}}

\eng{\textbf{EDITOR (Editorial -- Passive + Conditional):} ``And now, \textbf{The Editor's Desk}. Campus FM \textbf{has been awaited} for years by the English-speaking community here in Molyko. The studios \textbf{were renovated} last summer thanks to a partnership with the British Council. Programmes \textbf{must be submitted} for approval 48 hours before airing, which ensures quality. If our school \textbf{had not invested} in this equipment, we \textbf{would not have} this platform today. If we \textbf{had} a dedicated newsroom, we \textbf{would produce} daily bulletins. Let's make every broadcast count.''}

\eng{\textbf{[JINGLE: Quirky ``wrong answer'' buzzer sound, then cheerful music.]}}

\eng{\textbf{GRAMMAR EXPERT (Error of the Week):} ``It's time for \textbf{Error of the Week}! Listen to this sentence, heard in a Form 3 essay last week:}

\eng{`*The journalist \textbf{has interviewed} the minister \textbf{since} two hours.'}

\eng{\textbf{Stop! Wrong.} \emph{Since} indicates a \textbf{starting point} (e.g., \emph{since 10 a.m., since Monday}). For a \textbf{duration}, we use \emph{\textbf{for}} with the Present Perfect.}

\eng{\textbf{Correct version:} `The journalist \textbf{has interviewed} the minister \textbf{for} two hours.' Or: `The journalist \textbf{has been interviewing} the minister \textbf{for} two hours.'}

\eng{\textbf{Remember:} \emph{for} = duration (\emph{for 5 minutes, for 3 years}); \emph{since} = starting point (\emph{since 2020, since breakfast}). Don't mix them up!''}

\eng{\textbf{[JINGLE: Closing theme. Fade out.]}}

\eng{\textbf{ANCHOR:} ``That's all for today's \textbf{Grammar Clinic}. We are \textbf{The Broadcasters}. Stay tuned, stay grammatical. Goodbye!''}

\tcblower

\textbf{COMMENTAIRE DES CHOIX DE LANGUE (Corrigé commenté)}

\begin{enumerate}
    \item \textbf{Headlines $\to$ Reported Speech :}
    \begin{itemize}
        \item \eng{``We will launch'' $\to$ ``they would launch''} : backshifting \eng{will $\to$ would} + changement de pronom \eng{we $\to$ they}. \eng{Next month $\to$ the following month}.
        \item \eng{``is equipped'' $\to$ ``was equipped''} : passif au Present Simple $\to$ passif au Past Simple (backshifting de l'auxiliaire \eng{be}).
        \item \eng{``Can we broadcast'' $\to$ ``if they could broadcast''} : question fermée $\to$ \eng{if/whether} + sujet \eng{they} + \eng{could} (backshifting de \eng{can}). Pas d'inversion, pas de \eng{do}.
    \end{itemize}
    
    \item \textbf{Interview $\to$ Reported Questions :}
    \begin{itemize}
        \item \eng{``How has social media changed...?'' $\to$ ``I asked... how social media had changed...''} : question en \eng{wh-}. Backshifting Present Perfect $\to$ Past Perfect (\eng{had changed}). Ordre affirmatif (sujet + verbe), sans point d'interrogation.
        \item \eng{``Do you think...?'' $\to$ ``I asked her whether she thought...''} : question fermée $\to$ \eng{whether} (plus soutenu que \eng{if}). Backshifting Present Simple $\to$ Past Simple (\eng{thought}). \eng{Are} devient \eng{were} dans la subordonnée (\eng{whether they were important}).
    \end{itemize}
    
    \item \textbf{Éditorial (Passif + Conditionnel) :}
    \begin{itemize}
        \item \eng{has been awaited} (Present Perfect passif) : insiste sur le résultat présent d'une longue attente.
        \item \eng{were renovated} (Past Simple passif) : action datée (\eng{last summer}).
        \item \eng{must be submitted} (passif avec modal) : obligation institutionnelle, agent omis.
        \item \eng{If our school had not invested..., we would not have this platform today} : \textbf{conditionnel mixte} (condition au type 3 : \eng{If + Past Perfect} ; résultat au type 2 : \eng{would + BV}, car la conséquence est présente). Un type 3 complet serait : \eng{we would not have had this platform}.
        \item \eng{If we had a dedicated newsroom, we would produce...} : \textbf{conditionnel type 2} (obligatoire). Hypothèse présente irréelle (nous n'avons pas de salle de rédaction actuellement). \eng{Had} (et non \eng{*had had}), \eng{would produce}.
        \item \eng{which ensures quality} : proposition relative non restrictive avec \eng{which} (bonus du programme).
    \end{itemize}
    
    \item \textbf{Erreur de la semaine :} cible l'erreur n°2 de la fiche \textbf{Attention} (\eng{since/for} + choix du temps). Explication métalinguistique en anglais simple, exemples contrastifs. Prononciation à faire remarquer : \eng{for} se dit « fo » et \eng{since} se dit « sins ».
    
    \item \textbf{Cohésion / Lexique des médias :} \eng{anchor, correspondent, freelance journalist, blogger, media literacy, verification, platform, bulletin, submit for approval, partnership, renovate}. Niveau B2/C1 adapté à la Terminale A.
\end{enumerate}
\end{exoresolu}

\courssub{Bilan de l'activité}
\begin{aretenir}[Ce que cette activité a permis de valider]
\begin{itemize}
    \item \textbf{Compétence grammaticale :} l'élève est capable de \textbf{basculer} du discours direct au discours rapporté (backshifting, pronoms, repères spatio-temporels), de \textbf{manipuler la voix passive} à divers temps et avec des modaux, de \textbf{construire des conditionnels de type 2 (et 3)} correctement, et d'\textbf{identifier et corriger} une erreur fossilisée (\eng{since/for}, \eng{make/do}, prépositions).
    \item \textbf{Compétence discursive :} production d'un texte oral structuré (genre \emph{radio magazine}), respect des conventions de prise de parole (accroche, transitions, conclusion), adaptation du registre (formel pour l'éditorial, plus direct pour l'interview).
    \item \textbf{Compétence stratégique :} planification collaborative, auto-évaluation \emph{via} une grille de relecture, gestion du temps (5 min strictes), écoute sélective pour repérer des structures cibles chez les pairs.
    \item \textbf{Compétence culturelle :} ancrage dans le réel camerounais (GBHS Molyko, Buea, Douala, Pidgin English, British Council), valorisation de l'anglais comme langue de communication scolaire et médiatique.
\end{itemize}
\end{aretenir}

\begin{methode}[Pour réinvestir seul(e) : « Mini-Grammar Clinic » en autonomie]
\begin{enumerate}
    \item Choisis un sujet d'actualité camerounaise (ex. \eng{the Africa Cup of Nations, cocoa prices, the new school curriculum}).
    \item Écris 3 titres au discours direct. Transforme-les au discours rapporté (à l'écrit).
    \item Invente 2 questions pour un expert. Écris le résumé de ses réponses au style indirect.
    \item Rédige un micro-éditorial (5 lignes) : 3 passifs à des temps différents + 1 conditionnel de type 2.
    \item Piège-toi : écris une phrase contenant ton erreur favorite. Corrige-la et explique la règle en anglais.
    \item Enregistre-toi (téléphone). Réécoute-toi avec la grille d'écoute. Note tes progrès.
\end{enumerate}
\end{methode}

\begin{savaistu}[Le saviez-vous ? -- La radio au Cameroun]
La radiodiffusion au Cameroun remonte à l'époque coloniale, avec des émissions depuis Douala dès les années 1940. Après l'indépendance, la radio d'État (aujourd'hui la \eng{CRTV}, \eng{Cameroon Radio Television}) a longtemps dominé le paysage audiovisuel, avant la libéralisation des ondes autour de l'an 2000. Depuis, de nombreuses \textbf{stations privées} (\eng{Radio Balafon, Magic FM, Sweet FM}...) et des \textbf{radios communautaires, universitaires et scolaires} se sont développées, comme les radios des campus de Buea et de Yaoundé qui ont inspiré notre \eng{Campus FM} fictive. Le \eng{Cameroon Pidgin English}, langue véhiculaire majeure, y occupe une place croissante, aux côtés de l'anglais standard et du français, les deux langues officielles. Comme on dit à l'antenne : \eng{``Make we go!''} (« Allons-y ! »).
\end{savaistu}

\newpage

\coursec{Méthodes et exercices résolus}

\courssub{Lesson 58 — Grammar: General Revision}

\courssub{Diagnostic et carte mentale des structures}

\begin{methode}[Faire le point sur ses acquis : checklist de révision]
Avant d'entrer dans le détail, cochez les cases pour identifier vos zones de fragilité. Chaque item renvoie à une section de cette leçon.
\begin{enumerate}
\item \textbf{Temps verbaux :} Je distingue \cle{Present Simple/Continuous}, \cle{Past Simple/Continuous}, \cle{Present Perfect Simple/Continuous}, \cle{Past Perfect}, \cle{Future (will/going to)}. Je repère les marqueurs temporels. (\S 2)
\item \textbf{Voix passive :} Je transforme Actif $\rightarrow$ Passif à tous les temps (\eng{be + V-en}). Je sais omettre l'agent. (\S 3)
\item \textbf{Discours rapporté :} J'applique le \cle{backshift} (recul des temps), j'adapte pronoms et adverbes (\eng{now $\rightarrow$ then}), je respecte l'ordre affirmatif dans la subordonnée. (\S 3)
\item \textbf{Conditionnels \& Wish :} Je construis les 4 types (\cle{Type 0, 1, 2, 3, Mixed}) et les formes de \eng{Wish/If only} (regret présent/passé, volonté). (\S 4)
\item \textbf{Relatifs :} Je choisis \eng{who/which/whose/where/when}, je distingue définissante (pas de virgule, \eng{that} possible) et non-définissante (virgules, \eng{that} interdit). (\S 5)
\item \textbf{Modaux :} Je maîtrise les valeurs (capacité, obligation, déduction, conseil) au présent et au passé (\eng{modal + have + V-en}). (\S 5)
\item \textbf{Comparaison :} Je forme les comparatifs/superlatifs (courts/longs/irréguliers) et j'évite \eng{more better}, \eng{prefer ... than}. (\S 5)
\item \textbf{Gérondif / Infinitif :} Je classe les verbes (\eng{-ing} vs \eng{to}), je gère les changements de sens (\eng{stop, remember, forget, regret, try, mean}) et les prépositions (\eng{look forward to + -ing}). (\S 6)
\item \textbf{Pièges francophones :} Je traque les 10 erreurs « tueuses » (\eng{used to + -ing}, \eng{suggest + -ing}, \eng{interested}, \eng{since/for}, \eng{must/to}, \eng{tell/say}, \eng{depend on}, \eng{enough}, \eng{look forward to}, \eng{actually}). (\S 6)
\end{enumerate}
\end{methode}

\begin{popfigure}
\begin{tikzpicture}[mindmap, grow cyclic, every node/.style=concept, concept color=popblue,
  level 1/.append style={level distance=4cm, sibling angle=72},
  level 2/.append style={level distance=3cm, sibling angle=45},
  concept/.append style={font=\small\bfseries, text=white, execute at begin node=\hskip0pt}]
\node[align=center]{Grammaire\\Révision Générale}
  child[concept color=poporange] { node {Temps Verbaux}
    child { node {Marqueurs temporels} }
    child { node {Aspect : Simple/Cont./Perf.} }
    child { node {Concordance} }
  }
  child[concept color=popgreen] { node {Passif \& Discours Rapporté}
    child { node {Be + V-en (tous temps)} }
    child { node {Backshift + Pronoms/Adverbes} }
    child { node {Ordres/Questions rapportées} }
  }
  child[concept color=poppurple] { node {Conditionnels \& Wish}
    child { node {Type 0, 1, 2, 3, Mixed} }
    child { node {Wish/If only + Past/Perf./Would} }
    child { node {Unless, Provided that} }
  }
  child[concept color=poppink] { node {Relatifs, Modaux, Compar.}
    child { node {Who/Which/Whose/Where} }
    child { node {Defining vs Non-defining} }
    child { node {Modaux + Have V-en} }
    child { node {Er/Est, More/Most, Irrég.} }
  }
  child[concept color=popgold] { node {Gérondif/Infinitif \& Pièges}
    child { node {Verbes + -ing / + to} }
    child { node {Stop/Remember/Forget/Try} }
    child { node {Faux-amis (Actually...)} }
    child { node {Prépositions verbes/adj.} }
  };
\end{tikzpicture}
\legende{Carte mentale des 5 piliers de la révision grammaticale (Terminale A — Chapter 5)}
\end{popfigure}

\courssub{Révision des temps verbaux}

\begin{methode}[Réviser les temps verbaux : identifier, choisir, conjuguer]
\textbf{Étape 1 : Repérer les marqueurs temporels.} Soulignez les indicateurs de temps (\eng{now, at the moment, currently} $\rightarrow$ \cle{present continuous} ; \eng{yesterday, in 2010, ago, last week} $\rightarrow$ \cle{past simple} ; \eng{since, for, just, already, yet, ever, never} $\rightarrow$ \cle{present perfect} ; \eng{by the time, before, after} $\rightarrow$ \cle{past perfect} ; \eng{tomorrow, next week, in two days} $\rightarrow$ \cle{future}).
\textbf{Étape 2 : Déterminer l'aspect.} Action en cours (\cle{be + V-ing}) ? Action achevée avec résultat présent (\cle{have + V-en}) ? Antériorité (\cle{had + V-en}) ? Habitude ou vérité générale (\cle{present simple}) ? Projet/Prédiction (\cle{will/going to}) ?
\textbf{Étape 3 : Vérifier la concordance des temps.} Dans une phrase complexe, le temps de la subordonnée dépend souvent de la principale (ex. \eng{He said he \textbf{would} come} ; \eng{When I arrived, he \textbf{had already left}}).
\textbf{Étape 4 : Conjuguer sans erreur.} Attention aux verbes irréguliers (\eng{go/went/gone}, \eng{write/wrote/written}, \eng{take/took/taken}, \eng{see/saw/seen}) et à la 3\textsuperscript{e} personne du singulier au \cle{present simple} (\eng{he/she/it + s} : \eng{he watches, she goes, it does}).
\textbf{Structure à retenir :} \eng{Subject + Auxiliary (tense) + Verb (form) + Complements.}
\end{methode}

\begin{exoresolu}[Application : Texte à trous « Une matinée à Yaoundé »]
Complétez le texte ci-dessous en conjuguant les verbes entre parenthèses au temps qui convient. Justifiez chaque choix.

\textbf{Text:} ``Every morning, Pierre \textbf{\_\_\_\_\_} (wake up) at 5:30 a.m. He \textbf{\_\_\_\_\_} (prepare) his lessons when his phone \textbf{\_\_\_\_\_} (ring) yesterday. It \textbf{\_\_\_\_\_} (be) his friend Paul. ``I \textbf{\_\_\_\_\_} (wait) for you at the bendskin station for twenty minutes!'', he \textbf{\_\_\_\_\_} (say). Pierre \textbf{\_\_\_\_\_} (realise) he \textbf{\_\_\_\_\_} (forget) the appointment. He \textbf{\_\_\_\_\_} (rush) out immediately. By the time he \textbf{\_\_\_\_\_} (arrive), Paul \textbf{\_\_\_\_\_} (already / leave).''
\tcblower
\textbf{Solution détaillée et justifications :}
\begin{enumerate}
\item \textbf{wakes up} : \cle{Present simple} (habitude) + marqueur \eng{Every morning}. 3\textsuperscript{e} pers. sg $\rightarrow$ \eng{s}.
\item \textbf{was preparing} : \cle{Past continuous} (action en cours dans le passé) interrompue par \eng{when his phone rang}. \eng{was + V-ing}.
\item \textbf{rang} : \cle{Past simple} (action courte, ponctuelle, passée). Verbe irrégulier \eng{ring/rang/rung}.
\item \textbf{was} : \cle{Past simple} de \eng{be} (sujet \eng{It}, 3\textsuperscript{e} pers. sg).
\item \textbf{have been waiting} : \cle{Present perfect continuous} (action commencée dans le passé \eng{for twenty minutes} et qui continue encore au moment de parler). \eng{have been + V-ing}.
\item \textbf{said} : \cle{Past simple} (récit passé). Verbe irrégulier \eng{say/said/said}.
\item \textbf{realised} : \cle{Past simple} (action suivante dans la narration). Orthographe britannique \eng{-ise}.
\item \textbf{had forgotten} : \cle{Past perfect} (antériorité par rapport à \eng{realised} : l'oubli est antérieur à la prise de conscience). \eng{had + V-en}.
\item \textbf{rushed} : \cle{Past simple} (action principale suivante).
\item \textbf{arrived} : \cle{Past simple} (marqueur \eng{By the time} introduit une action de référence au passé).
\item \textbf{had already left} : \cle{Past perfect} (antériorité : le départ de Paul est antérieur à l'arrivée de Pierre). \eng{already} se place entre auxiliaire et participe.
\end{enumerate}
\textbf{Conclusion :} La maîtrise des marqueurs temporels et de la chronologie (antériorité, simultanéité, postériorité) est la clé pour ne pas mélanger \cle{past simple}, \cle{past continuous} et \cle{past perfect}.
\end{exoresolu}

\courssub{Voix passive et discours rapporté}

\begin{methode}[Maîtriser la voix passive et le discours rapporté]
\textbf{Voix passive (Passive Voice) :}
\textbf{1.} Identifier l'objet du complément d'objet direct (COD) de la phrase active $\rightarrow$ devient \textbf{sujet} de la passive.
\textbf{2.} Mettre l'auxiliaire \eng{be} au \textbf{même temps} que le verbe actif.
\textbf{3.} Mettre le verbe principal au \textbf{participe passé} (V-en).
\textbf{4.} L'ancien sujet devient \textbf{complément d'agent} (introduit par \eng{by}) ou est omis s'il est inconnu/général (\eng{someone, people, they}).
\textbf{Attention :} Seuls les verbes transitifs directs (suivis de COD) peuvent se mettre au passif. \eng{Happen, sleep, arrive, seem, lack} ne s'utilisent pas au passif.

\textbf{Discours rapporté (Reported Speech) :}
\textbf{1.} Identifier le verbe introducteur (\eng{say, tell, ask, wonder, advise, order, promise}). S'il est au \textbf{passé} (\eng{said, told, asked}) $\rightarrow$ \textbf{backshift} (recul d'un cran des temps).
\textbf{2.} Appliquer la \textbf{concordance des temps} :
\begin{center}
\begin{tabular}{|l|l|}\hline
\textbf{Direct} & \textbf{Rapporté (Backshift)} \\ \hline
Present Simple & Past Simple \\ \hline
Present Continuous & Past Continuous \\ \hline
Present Perfect / Past Simple & Past Perfect \\ \hline
Will & Would \\ \hline
Can & Could \\ \hline
Must (obligation) & Had to \\ \hline
Must (déduction) & Must (inchangé) \\ \hline
May / Might & Might \\ \hline
\end{tabular}
\end{center}
\textbf{3.} Adapter \textbf{pronoms} (I $\rightarrow$ he/she), \textbf{adverbes temps/lieu} (now $\rightarrow$ then, today $\rightarrow$ that day, tomorrow $\rightarrow$ the next day, here $\rightarrow$ there, this $\rightarrow$ that).
\textbf{4.} Structure :
\begin{itemize}
\item \textbf{Statement :} \eng{Subject + verb + (that) + clause.}
\item \textbf{Question (Wh-) :} \eng{Subject + verb + wh-word + clause (ordre affirmatif : sujet + verbe).}
\item \textbf{Question (Yes/No) :} \eng{Subject + verb + if/whether + clause (ordre affirmatif).}
\item \textbf{Imperative :} \eng{Subject + verb + object + to-infinitive.} (Négatif : \eng{not to + infinitive}).
\end{itemize}
\end{methode}

\begin{exoresolu}[Transformation : Passif \& Discours rapporté (Type Bac)]
\textbf{A. Passive Voice.} Transformez les phrases actives en passives. Conservez le temps.
1. \eng{The students are cleaning the classroom now.}
2. \eng{Someone stole my phone yesterday.}
3. \eng{They will have finished the project by Friday.}
4. \eng{You must submit the application before Friday.}

\textbf{B. Reported Speech.} Rapportez les propos suivants (verbe introducteur au passé).
5. \eng{Paul: ``I am reading a fascinating article about AI.''}
6. \eng{The teacher: ``Did you do your homework?''}
7. \eng{Mother: ``Don't forget to buy bread!''}
8. \eng{The coach: ``Run faster!''}
\tcblower
\textbf{Solution détaillée :}
\textbf{A. Passive Voice}
1. \eng{The classroom \textbf{is being cleaned} by the students now.} (Present continuous : \eng{be} $\rightarrow$ \eng{is being} + V-en).
2. \eng{My phone \textbf{was stolen} yesterday.} (Past simple : \eng{be} $\rightarrow$ \eng{was} + V-en. Agent \eng{by someone} omis).
3. \eng{The project \textbf{will have been finished} by Friday.} (Future perfect : \eng{will have been} + V-en).
4. \eng{The application \textbf{must be submitted} before Friday.} (Modal \eng{must} : \eng{modal + be + V-en}).

\textbf{B. Reported Speech (Backshift appliqué)}
5. \eng{Paul said (that) \textbf{he was reading} a fascinating article about AI \textbf{the previous day / the day before}.}
\begin{itemize}
\item \eng{am reading} (pres. cont.) $\rightarrow$ \eng{was reading} (past cont.).
\item \eng{I} $\rightarrow$ \eng{he}.
\item \eng{today/now} (implicite) $\rightarrow$ \eng{the previous day}.
\end{itemize}
6. \eng{The teacher asked \textbf{if/whether I had done} my homework.}
\begin{itemize}
\item Question fermée $\rightarrow$ \eng{if/whether}.
\item Ordre affirmatif (sujet + verbe) : \eng{I had done} (pas \eng{did I do}).
\item \eng{did ... do} (past simple) $\rightarrow$ \eng{had done} (past perfect).
\item \eng{your} $\rightarrow$ \eng{my}.
\end{itemize}
7. \eng{Mother told/reminded me \textbf{not to forget} to buy bread.}
\begin{itemize}
\item Impératif négatif $\rightarrow$ \eng{not to + infinitif}.
\item Verbe introducteur : \eng{tell/remind/order + objet + to-inf}.
\end{itemize}
8. \eng{The coach ordered/advised me \textbf{to run} faster.}
\begin{itemize}
\item Impératif affirmatif $\rightarrow$ \eng{to + infinitif}.
\end{itemize}
\textbf{Conclusion :} La transformation mécanique (backshift + pronoms + adverbes + ordre des mots) doit devenir un réflexe automatique. Vérifiez toujours l'ordre des mots dans la subordonnée rapportée (sujet avant verbe).
\end{exoresolu}

\courssub{Conditionnels, souhaits et hypothèses}

\begin{methode}[Structurer les conditionnels, souhaits et hypothèses]
\textbf{Conditionnels (If-clauses) :} Identifier le \textbf{type} selon la probabilité/réalité.
\begin{itemize}
\item \textbf{Type 0 (Réalité générale / Loi scientifique) :} \eng{If + Present Simple, ... Present Simple.} \eng{If you heat ice, it melts.}
\item \textbf{Type 1 (Possible / Futur probable) :} \eng{If + Present Simple, ... Will/Can/May/Might/Should + Inf.} \eng{If it rains, I will stay home.}
\item \textbf{Type 2 (Imaginaire / Présent irréel / Conseil) :} \eng{If + Past Simple, ... Would/Could/Might + Inf.} \eng{If I knew, I would tell you.} $\rightarrow$ \eng{were} pour \textbf{toutes} les personnes (\eng{If I \textbf{were} you...}).
\item \textbf{Type 3 (Regret passé / Impossible / Contre-factuel) :} \eng{If + Past Perfect, ... Would/Could/Might + Have + V-en.} \eng{If I had studied, I would have passed.}
\item \textbf{Mixed Conditionals :}
\begin{itemize}
\item Passé $\rightarrow$ Présent : \eng{If + Past Perfect, ... Would + Inf.} (\eng{If I had taken the medicine, I would feel better now.})
\item Présent $\rightarrow$ Passé : \eng{If + Past Simple, ... Would + Have + V-en.} (\eng{If he were faster, he would have won the race.})
\end{itemize}
\end{itemize}
\textbf{Mots de liaison :} \eng{Unless = If not} (\eng{Unless you study, you will fail}). \eng{Provided/Providing (that) / As long as / On condition that / Supposing} + \textbf{present} pour valeur future.

\textbf{Wish / If only :}
\begin{itemize}
\item \textbf{Regret présent/futur (irréel présent) :} \eng{Wish/If only + Past Simple.} \eng{I wish I \textbf{knew} the answer.}
\item \textbf{Regret passé (irréel passé) :} \eng{Wish/If only + Past Perfect.} \eng{I wish I \textbf{had come} earlier.}
\item \textbf{Ennui / Volonté de changement (autre sujet) :} \eng{Wish + Would + Inf.} \eng{I wish it \textbf{would stop} raining.} \eng{I wish he \textbf{would listen}.}
\end{itemize}
\textbf{Attention :} \eng{Wish + Would} \textbf{jamais} avec le même sujet (\eng{*I wish I would go} $\rightarrow$ \eng{I want to go} ou \eng{I wish I could go}).
\end{methode}

\begin{exoresolu}[Conditionnels, Wish \& Mixte : Complétion et Réécriture]
\textbf{1. Complétez avec la forme correcte du verbe.}
a) If you \textbf{\_\_\_\_\_} (heat) water to 100°C, it \textbf{\_\_\_\_\_} (boil).
b) If I \textbf{\_\_\_\_\_} (be) you, I \textbf{\_\_\_\_\_} (apply) for that scholarship.
c) If they \textbf{\_\_\_\_\_} (not / miss) the bus, they \textbf{\_\_\_\_\_} (arrive) on time yesterday.
d) I wish I \textbf{\_\_\_\_\_} (can) speak English fluently. (Regret présent)
e) She wishes she \textbf{\_\_\_\_\_} (not / spend) all her money last week.

\textbf{2. Réécrivez les phrases en commençant par \eng{I wish / If only}.}
f) \eng{I don't know the answer.} $\rightarrow$ \eng{I wish...}
g) \eng{He didn't wake up early.} $\rightarrow$ \eng{If only...}
h) \eng{It is raining now (and I hate it).} $\rightarrow$ \eng{I wish...}
\tcblower
\textbf{Solution détaillée :}
\textbf{1. Complétion}
a) \eng{heat / boils} : \textbf{Type 0} (vérité scientifique). Deux \cle{present simple}.
b) \eng{were / would apply} : \textbf{Type 2} (hypothèse imaginaire présente). \eng{If I \textbf{were} you} (subjonctif/forme \eng{were} pour toutes personnes). \eng{would + inf}.
c) \eng{hadn't missed / would have arrived} : \textbf{Type 3} (regret passé, fait contraire à la réalité). \eng{If + past perfect}, \eng{would have + V-en}.
d) \eng{could} : \textbf{Wish + regret présent} $\rightarrow$ \eng{past simple} (ici modal \eng{can} $\rightarrow$ \eng{could}). \eng{I wish I \textbf{could} speak...}
e) \eng{hadn't spent} : \textbf{Wish + regret passé} $\rightarrow$ \eng{past perfect}. \eng{She wishes she \textbf{hadn't spent}...}

\textbf{2. Réécriture (Wish / If only)}
f) \eng{I wish I \textbf{knew} the answer.} (Regret présent : \eng{don't know} $\rightarrow$ \eng{knew} / past simple).
g) \eng{If only he \textbf{had woken} up early.} (Regret passé : \eng{didn't wake} $\rightarrow$ \eng{had woken} / past perfect).
h) \eng{I wish it \textbf{would stop} raining.} (Volonté/Ennui présent : sujet différent \eng{it} + \eng{would + inf}). \textbf{Pas} \eng{I wish it stopped} (ce qui exprimerait un regret statique, pas une demande de changement).
\textbf{Conclusion :} La distinction Type 2 / Type 3 et Wish Present / Wish Past est l'une des plus testées au Bac. Mémorisez les formules « \eng{If + Past, Would + Inf} » et « \eng{If + Had + V-en, Would + Have + V-en} ».
\end{exoresolu}

\courssub{Relatifs, modaux et comparaison}

\begin{methode}[Relatifs, Modaux et Comparaison : Précision et nuances]
\textbf{Propositions relatives (Relative Clauses) :}
\textbf{1. Choix du pronom :}
\begin{center}
\begin{tabular}{|l|l|l|}\hline
\textbf{Antécédent / Fonction} & \textbf{Pronom} & \textbf{Exemple} \\ \hline
Personne (Sujet) & \eng{who / that} & \eng{The man \textbf{who/that} called} \\ \hline
Personne (Objet) & \eng{who / whom / that / Ø} & \eng{The man \textbf{(who/that)} I saw} \\ \hline
Chose (Sujet/Objet) & \eng{which / that} & \eng{The book \textbf{which/that} fell / \textbf{(which/that)} I read} \\ \hline
Possesseur (Pers./Chose) & \eng{whose} & \eng{The man \textbf{whose} car was stolen} \\ \hline
Lieu & \eng{where} & \eng{The school \textbf{where} I studied} \\ \hline
Temps & \eng{when} & \eng{The day \textbf{when} we met} \\ \hline
Raison & \eng{why} & \eng{The reason \textbf{why} he left} \\ \hline
\end{tabular}
\end{center}
\textbf{2. Définissante (Defining / Restrictive) vs Non-définissante (Non-defining / Non-restrictive) :}
\begin{itemize}
\item \textbf{Définissante :} Information essentielle pour identifier l'antécédent. \textbf{Pas de virgule}. \eng{That} possible. Pronom objet omissible. \eng{The student who/that won the prize is my friend.}
\item \textbf{Non-définissante :} Information supplémentaire, déjà identifié. \textbf{Virgules obligatoires}. \eng{That} \textbf{interdit}. Pronom \textbf{non omissible}. \eng{My brother, who lives in Douala, is visiting.}
\end{itemize}
\textbf{3. Préposition en fin :} \eng{The man \textbf{(who/that)} I spoke \textbf{to}} (standard) vs \eng{The man \textbf{to whom} I spoke} (formel). \eng{The chair \textbf{(which/that)} I sat \textbf{on}}.

\textbf{Verbes modaux (Modals) — Valeurs et formes passées :}
\begin{center}
\begin{tabularx}{\linewidth}{|X|X|X|}
\hline \rowcolor{popblueL} \textbf{Fonction} & \textbf{Present / Future} & \textbf{Past / Perfect Infinitive} \\ \hline
Capacité & \eng{can / be able to} & \eng{could} (générale) / \eng{was able to} (action précise réussie) \\ \hline
Obligation (extérieure) & \eng{must / have to} & \eng{had to} \\ \hline
Interdiction & \eng{mustn't / can't} & \eng{couldn't / wasn't allowed to} \\ \hline
Absence d'obligation & \eng{don't have to / needn't} & \eng{didn't have to / didn't need to} \\ \hline
Conseil / Opinion & \eng{should / ought to} & \eng{should / ought to + have + V-en} (\eng{You should have come}) \\ \hline
Déduction forte (Certitude) & \eng{must / can't} & \eng{must / can't + have + V-en} \\ \hline
Possibilité / Probabilité & \eng{may / might / could} & \eng{may / might / could + have + V-en} \\ \hline
\end{tabularx}
\end{center}
\textbf{Rappel :} Pas de \eng{to} après modaux (sauf \eng{ought to, have to, be able to, used to}). Pas de \eng{s} à la 3\textsuperscript{e} pers. sg. Pas de \eng{do/does/did} pour question/négation.

\textbf{Comparaison (Comparison) :}
\textbf{1. Formes :}
\begin{itemize}
\item Court adjectif (1 syllabe, qq 2 syll. en \eng{-y, -er, -ow}) : \eng{-er / -est} (\eng{fast, faster, fastest} ; \eng{happy, happier, happiest}).
\item Long adjectif (2+ syllabes) : \eng{more / most} (\eng{beautiful, more beautiful, most beautiful}).
\item Irréguliers : \eng{good/well $\rightarrow$ better $\rightarrow$ best} ; \eng{bad/badly $\rightarrow$ worse $\rightarrow$ worst} ; \eng{far $\rightarrow$ further/farther $\rightarrow$ furthest/farthest} ; \eng{much/many $\rightarrow$ more $\rightarrow$ most} ; \eng{little $\rightarrow$ less $\rightarrow$ least}.
\end{itemize}
\textbf{2. Structures :}
\begin{itemize}
\item Égalité : \eng{as + adj/adv + as} (\eng{as tall as}).
\item Supériorité : \eng{comparative + than} (\eng{taller than, more interesting than}).
\item Infériorité : \eng{less + adj/adv + than} (\eng{less expensive than}).
\item Corrélation : \eng{The + comparative ..., the + comparative} (\eng{The more you read, the more you learn}).
\end{itemize}
\textbf{3. Pièges majeurs :}
\begin{itemize}
\item \eng{Prefer + noun/gerund + to + noun/gerund} (\eng{I prefer coffee to tea} / \eng{I prefer reading to watching TV}). \textbf{Pas} \eng{than}.
\item \eng{Would rather + bare infinitive + than + bare infinitive} (\eng{I would rather go than stay}).
\item \eng{Similar to / Different from} (\textbf{pas} \eng{than}).
\item Double comparatif interdit : \eng{more better, more easier}.
\end{itemize}
\end{methode}

\begin{exoresolu}[Synthèse : Relatifs, Modaux, Comparaison]
\textbf{1. Reliez les phrases par un pronom relatif (indiquez si virgules nécessaires).}
a) \eng{The journalist interviewed the minister. The minister arrived late.}
b) \eng{My uncle lives in Buea. He is a doctor.}
c) \eng{I bought a new phone. Its battery lasts two days.}

\textbf{2. Choisissez le modal approprié (déduction/conseil/obligation).}
d) \eng{You \_\_\_\_\_ drive fast near schools. It is forbidden.} (Interdiction)
e) \eng{Look at those clouds! It \_\_\_\_\_ rain soon.} (Certitude forte future basée sur indices)
f) \eng{He \_\_\_\_\_ have taken the wrong train; he is in Douala, not Yaoundé.} (Déduction passée forte)

\textbf{3. Corrigez les erreurs de comparaison.}
g) \eng{This exercise is more easier than the last one.}
h) \eng{She speaks English more good than her brother.}
i) \eng{He is the most tall student in the class.}
j) \eng{I prefer coffee than tea.}
\tcblower
\textbf{Solution détaillée :}
\textbf{1. Relatifs}
a) \eng{The minister \textbf{who/that} the journalist interviewed arrived late.} (Définissante : pas de virgules. \eng{The minister} = objet de \eng{interviewed} $\rightarrow$ \eng{who/that} ou omis. Mais ici sujet de \eng{arrived}, on garde le pronom).
b) \eng{My uncle, \textbf{who} is a doctor, lives in Buea.} (Non-définissante : info supplémentaire sur l'oncle $\rightarrow$ \textbf{virgules obligatoires}, \eng{who} pas \eng{that}).
c) \eng{I bought a new phone \textbf{whose} battery lasts two days.} (\eng{Whose} = possessif pour choses aussi. Pas \eng{which battery}).

\textbf{2. Modaux}
d) \eng{\textbf{mustn't}} (Interdiction formelle/règle). \eng{Can't} possible mais \eng{mustn't} plus fort pour règle écrite.
e) \eng{\textbf{must}} (Déduction forte présente : preuves visibles \eng{clouds}). \eng{Will} = prédiction pure, pas déduction.
f) \eng{\textbf{must}} (Déduction forte passée : \eng{must have + V-en}). \eng{Can't have taken} = impossibilité logique.

\textbf{3. Comparaison (Erreurs classiques francophones)}
g) \eng{This exercise is \textbf{easier} than the last one.} (\textbf{Erreur :} \eng{more easier} = double comparatif interdit. Court adj $\rightarrow$ \eng{-er}).
h) \eng{She speaks English \textbf{better} than her brother.} (\textbf{Erreur :} \eng{more good} n'existe pas. Irrégulier \eng{good/well} $\rightarrow$ \eng{better}).
i) \eng{He is the \textbf{tallest} student in the class.} (\textbf{Erreur :} \eng{most tall}. Court adj $\rightarrow$ \eng{-est}).
j) \eng{I prefer coffee \textbf{to} tea.} (\textbf{Erreur :} \eng{prefer ... than}. Structure correcte : \eng{prefer + noun/gerund + to + noun/gerund}). Ou \eng{I \textbf{would rather} drink coffee \textbf{than} tea.}
\textbf{Conclusion :} Les modaux au passé (\eng{have + V-en}) et les irréguliers de comparaison (\eng{good/better/best, bad/worse/worst, far/further/furthest, much/many/more/most, little/less/least}) sont des points de bonus faciles s'ils sont maîtrisés.
\end{exoresolu}

\courssub{Gérondif / Infinitif et pièges des francophones}

\begin{methode}[Gérondif vs Infinitif \& Pièges des francophones : Le dernier kilomètre]
\textbf{Verbes suivis du gérondif (\eng{V-ing}) :}
\eng{enjoy, avoid, mind, deny, risk, suggest, imagine, finish, stop (arrêter de), regret (regretter d'avoir fait), remember (se souvenir d'avoir fait), go (activités : go swimming), admit, consider, delay, dislike, fancy, involve, keep (continuer), miss, postpone, practise, quit, recommend, can't help, can't stand, look forward to, be used to, get used to.}
\textbf{+ Prépositions :} \eng{interested in, good at, bad at, keen on, fond of, used to + -ing, look forward to + -ing, succeed in, insist on, apologise for, accuse of, prevent from.}

\textbf{Verbes suivis de l'infinitif (\eng{to + V}) :}
\eng{want, decide, hope, plan, promise, refuse, agree, manage, afford, learn, choose, fail, offer, prepare, threaten, pretend, claim, deserve, expect, need, would like/love/prefer/hate.}
\textbf{+ Adjectifs :} \eng{happy to, sorry to, surprised to, difficult to, easy to, ready to, likely to.}
\textbf{+ Too / Enough :} \eng{too hot to eat, not old enough to vote, enough money to buy}. Ordre : \eng{Adj/Adv + enough} mais \eng{enough + Noun}.

\textbf{Verbes + \eng{-ing} OU \eng{to} (CHANGEMENT DE SENS) :}
\begin{center}
\begin{tabularx}{\linewidth}{|X|X|X|}
\hline \rowcolor{popblueL} \textbf{Verbe} & \textbf{+ V-ing} & \textbf{+ to Inf} \\ \hline
\eng{Stop} & Arrêter une action en cours (\eng{Stop smoking}) & S'arrêter POUR faire qqch (\eng{Stop to smoke}) \\ \hline
\eng{Remember} & Se souvenir d'une action PASSÉE (\eng{Remember locking the door}) & Se souvenir DE FAIRE qqch FUTUR (\eng{Remember to lock the door}) \\ \hline
\eng{Forget} & Oublier une action PASSÉE (\eng{Forget meeting him}) & Oublier DE FAIRE qqch FUTUR (\eng{Forget to meet him}) \\ \hline
\eng{Regret} & Regretter d'AVOIR FAIT (\eng{Regret telling him}) & Regretter de DEVOIR FAIRE (formel : \eng{Regret to inform you}) \\ \hline
\eng{Try} & Essayer de / par curiosité (\eng{Try opening the window}) & Faire un effort pour / tenter de (\eng{Try to open the window}) \\ \hline
\eng{Mean} & Signifier / Impliquer (\eng{This means war}) & Avoir l'intention de (\eng{I meant to call you}) \\ \hline
\eng{Go on} & Continuer la même action (\eng{Go on reading}) & Passer à l'action suivante (\eng{Go on to read}) \\ \hline
\end{tabularx}
\end{center}

\textbf{Top 10 Pièges Francophones (Faux-amis / Calques / Structures) :}
\begin{enumerate}
\item \eng{Actually} = \textbf{En fait / En réalité} (pas \eng{actuellement} $\rightarrow$ \eng{currently, at present, nowadays}).
\item \eng{Eventually} = \textbf{Finalement / Au bout du compte} (pas \eng{éventuellement} $\rightarrow$ \eng{possibly, potentially}).
\item \eng{Sensible} = \textbf{Raisonnable / Sensé} (pas \eng{sensible} $\rightarrow$ \eng{sensitive}).
\item \eng{Library} = \textbf{Bibliothèque} (pas \eng{librairie} $\rightarrow$ \eng{bookshop / bookstore}).
\item \eng{Opportunity} = \textbf{Occasion / Chance} (pas \eng{opportunité} $\rightarrow$ \eng{appropriateness, timeliness}).
\item \eng{Make a mistake / a decision / an effort / a choice / money / progress} (pas \eng{do}).
\item \eng{Do homework / housework / the shopping / business / an exam / a favour} (pas \eng{make}).
\item \eng{Since} (point de départ : \eng{since Monday, since 2010}) vs \eng{For} (durée : \eng{for two hours, for ages}).
\item \eng{Look forward to + V-ing} (pas \eng{to + inf}). \eng{To} = préposition.
\item \eng{Depend \textbf{on}} (pas \eng{of}) ; \eng{Listen \textbf{to}} ; \eng{Wait \textbf{for}} ; \eng{Interested \textbf{in}} ; \eng{Married \textbf{to}} ; \eng{Similar \textbf{to}} ; \eng{Good \textbf{at}}.
\end{enumerate}
\end{methode}

\begin{exoresolu}[Correction d'erreurs « Spécial Francophones » (10 phrases)]
Corrigez les erreurs (grammaire, vocabulaire, structure) dans les phrases suivantes. Expliquez la règle.
1. \eng{I am used to wake up early.}
2. \eng{He suggested me to go to the cinema.}
3. \eng{The match was very interested.}
4. \eng{She has been waiting for you since two hours.}
5. \eng{I must to finish my homework.}
6. \eng{He told that he was tired.}
7. \eng{It depends of the weather.}
8. \eng{He is not enough tall to play basketball.}
9. \eng{I look forward to see you.}
10. \eng{Actually, I live in Douala.}
\tcblower
\textbf{Solution détaillée (Correction + Justification) :}
1. \eng{I am used to \textbf{waking} up early.} (\eng{Used to + V-ing} = avoir l'habitude DE faire. \eng{Used to + inf} = avait l'habitude de (passé révolu). \eng{To} = préposition ici $\rightarrow$ gérondif).
2. \eng{He suggested \textbf{going} to the cinema.} OU \eng{He suggested \textbf{that I (should) go} to the cinema.} (\textbf{Erreur :} \eng{suggest} ne prend JAMAIS \eng{to-inf} ni COI direct \eng{me}. Structure : \eng{suggest + V-ing} ou \eng{suggest + that-clause}).
3. \eng{The match was very \textbf{interesting}.} (\textbf{Participe adjectif :} \eng{-ing} = cause du sentiment (le match \textbf{intéresse}) ; \eng{-ed} = ressenti du sujet (j'étais \textbf{intéressé}). \eng{An interested match} = un match qui a des intérêts financiers).
4. \eng{She has been waiting for you \textbf{for} two hours.} (\textbf{Since} = point de départ \eng{since 2010, since Monday, since 10am} ; \textbf{For} = durée \eng{for two hours, for ages, for a long time}).
5. \eng{I \textbf{must} finish my homework.} (\textbf{Modal :} \eng{must, can, could, should, might, shall, will, would, may} $\rightarrow$ \textbf{jamais de \eng{to}} après. \eng{Have to} $\rightarrow$ \eng{to}).
6. \eng{He \textbf{told me} that he was tired.} OU \eng{He \textbf{said} that he was tired.} (\textbf{Transitivité :} \eng{Tell} = dire \textbf{à qqn} $\rightarrow$ COI obligatoire \eng{tell me/him/us/them}. \eng{Say} = dire \textbf{qqch} $\rightarrow$ pas de COI).
7. \eng{It depends \textbf{on} the weather.} (\textbf{Verbe prépositionnel :} \eng{Depend ON}. Calque français « dépendre \textbf{de} »).
8. \eng{He is \textbf{not tall enough} to play basketball.} (\textbf{Ordre :} \eng{Enough} suit l'adjectif/adverbe (\eng{tall enough, quickly enough}), mais précède le nom (\eng{enough time, enough money}). \eng{Not ... enough} = pas assez).
9. \eng{I look forward to \textbf{seeing} you.} (\textbf{Structure figée :} \eng{Look forward to + V-ing}. \eng{To} = préposition, pas marqueur d'infinitif. Pareil pour \eng{be used to + -ing, get used to + -ing}).
10. \eng{\textbf{Currently} / \textbf{At present} / \textbf{Nowadays}, I live in Douala.} (\textbf{Faux ami :} \eng{Actually} = \eng{En fait / En réalité} (correction/contraste : \eng{Actually, I don't like it}). \eng{Actuellement} = \eng{Currently / At the moment / Nowadays}).
\textbf{Conclusion :} Ces 10 erreurs couvrent 80\% des points perdus en grammaire au Bac. Les apprendre par cœur (liste « \eng{Stop/Remember/Forget/Regret/Try/Mean/Go on} », \eng{Suggest/Tell/Say}, \eng{Used to}, \eng{Look forward to}, \eng{Enough}, \eng{Since/For}, Faux-amis \eng{Actually/Eventually/Sensible/Library/Opportunity}, \eng{Make/Do}, Prépositions) garantit un gain de 5 à 10 points.
\end{exoresolu}

\courssub{Entraînement type Bac et production guidée}

\begin{exercice}[Grammaire : QCM Type Bac (20 items)]
Choisissez la bonne réponse (a, b, c ou d).
\begin{enumerate}
\item By the time the police arrived, the thieves \_\_\_\_\_.
  a) escaped \quad b) have escaped \quad c) had escaped \quad d) escape
\item I wish I \_\_\_\_\_ harder for the exam. I failed.
  a) studied \quad b) had studied \quad c) would study \quad d) study
\item The bridge \_\_\_\_\_ next year according to the plan.
  a) will build \quad b) will be built \quad c) is building \quad d) builds
\item She asked me where \_\_\_\_\_.
  a) I live \quad b) do I live \quad c) did I live \quad d) I lived
\item If only he \_\_\_\_\_ so much noise! I can't concentrate.
  a) didn't make \quad b) wouldn't make \quad c) doesn't make \quad d) hadn't made
\item This is the best movie \_\_\_\_\_ I have ever seen.
  a) who \quad b) which \quad c) that \quad d) Ø (rien)
\item You \_\_\_\_\_ have seen him; he is in London.
  a) must \quad b) can't \quad c) might \quad d) should
\item I'm tired. I'd rather \_\_\_\_\_ at home tonight.
  a) stay \quad b) staying \quad c) to stay \quad d) stayed
\item He denied \_\_\_\_\_ the money.
  a) to take \quad b) taking \quad c) take \quad d) took
\item The teacher made us \_\_\_\_\_ the exercise again.
  a) to do \quad b) doing \quad c) do \quad d) did
\item She speaks English \_\_\_\_\_ than her sister.
  a) more good \quad b) better \quad c) weller \quad d) more well
\item I'm not used to \_\_\_\_\_ so early.
  a) get up \quad b) getting up \quad c) got up \quad d) have got up
\item The report \_\_\_\_\_ by the committee tomorrow.
  a) will discuss \quad b) will be discussed \quad c) discusses \quad d) is discussed
\item He \_\_\_\_\_ be at home; his car is in the driveway.
  a) must \quad b) can't \quad c) might \quad d) should
\item I regret \_\_\_\_\_ you that your application was rejected.
  a) to inform \quad b) informing \quad c) inform \quad d) informed
\item The man \_\_\_\_\_ car was stolen has gone to the police.
  a) who \quad b) which \quad c) whose \quad d) whom
\item She \_\_\_\_\_ me not to open the window.
  a) said \quad b) told \quad c) asked \quad d) suggested
\item This exercise is \_\_\_\_\_ difficult for me.
  a) too \quad b) enough \quad c) very \quad d) so
\item I \_\_\_\_\_ him since we left school.
  a) didn't see \quad b) haven't seen \quad c) don't see \quad d) hadn't seen
\item \_\_\_\_\_ you work hard, you won't pass.
  a) If \quad b) Unless \quad c) Provided \quad d) As long as
\end{enumerate}
\end{exercice}

\begin{exercice}[Transformation de phrases (Type Bac — 10 items)]
Réécrivez chaque phrase en commençant par les mots donnés, sans en changer le sens.
\begin{enumerate}
\item ``I will call you tomorrow,'' he said. $\rightarrow$ He said that he \_\_\_\_\_.
\item They are building a new stadium in Buea. $\rightarrow$ A new stadium \_\_\_\_\_.
\item ``Don't touch the wire!'' the electrician said to me. $\rightarrow$ The electrician warned me \_\_\_\_\_.
\item I didn't buy the book because I didn't have enough money. $\rightarrow$ If I \_\_\_\_\_, I \_\_\_\_\_.
\item It is forbidden to smoke here. $\rightarrow$ You \_\_\_\_\_.
\item ``Have you finished your work?'' the boss asked. $\rightarrow$ The boss asked \_\_\_\_\_.
\item He is too young to drive. $\rightarrow$ He is not \_\_\_\_\_.
\item I regret not learning English earlier. $\rightarrow$ I wish \_\_\_\_\_.
\item The last time I saw him was in 2015. $\rightarrow$ I \_\_\_\_\_ him \_\_\_\_\_ 2015.
\item ``Let's go to the cinema,'' she said. $\rightarrow$ She suggested \_\_\_\_\_.
\end{enumerate}
\end{exercice}

\begin{experience}[Débat oral : « Social Media — A Blessing or a Curse? »]
\textbf{Contexte :} Dans le cadre du Chapter 5 (Media and Communication), organisez un débat structuré.
\textbf{Consignes :}
\begin{enumerate}
\item \textbf{Préparation (10 min) :} En groupes de 4, listez 3 arguments \textbf{Pour} (connectivité, information, business, éducation) et 3 arguments \textbf{Contre} (addiction, cyberharcèlement, fake news, isolement, vie privée). Utilisez le vocabulaire du chapitre (\eng{social media, platform, user, content, algorithm, privacy, cyberbullying, misinformation, influencer, viral, screen time, digital detox}).
\item \textbf{Structure du débat (20 min) :}
  \begin{itemize}
  \item \textbf{Opening statements} (1 min par camp) : \eng{``In my opinion... / It is undeniable that... / One major advantage/drawback is...''}
  \item \textbf{Rebuttals / Répliques} (tour à tour) : \eng{``That's true, but... / I see your point, however... / On the contrary...''}
  \item \textbf{Conclusion} (30 sec par camp) : \eng{``To sum up... / All things considered...''}
  \end{itemize}
\item \textbf{Contraintes linguistiques :} Utilisez au moins 3 structures de la révision : 1 conditionnel (\eng{If we regulate..., it would...}), 1 passif (\eng{Data is collected...}), 1 discours rapporté (\eng{Experts say that...}), 1 comparatif (\eng{More harmful than...}), 1 gérondif (\eng{Stopping scrolling is hard}).
\end{enumerate}
\textbf{Évaluation :} Fluidité, justesse grammaticale (ciblée sur les points révisés), richesse lexicale, interaction.
\end{experience}

\begin{methode}[Production écrite guidée : Article d'opinion (Composition Type Bac)]
\textbf{Sujet :} ``Social media has done more harm than good to today's youth.'' Discuss. (250--300 mots)
\textbf{Plan type (4 paragraphes) :}
\textbf{1. Introduction (40--50 mots) :}
\begin{itemize}
\item Accroche : \eng{Nowadays, social media platforms like Facebook, WhatsApp and TikTok are ubiquitous in young people's lives.}
\item Présentation du débat : \eng{While they offer undeniable benefits, a growing number of voices warn against their negative impact.}
\item Thèse (votre opinion) : \eng{This essay will argue that, despite their utility, the dangers outweigh the advantages if usage is not regulated.}
\end{itemize}
\textbf{2. Arguments « Pour » (Avantages) (70--80 mots) :}
\begin{itemize}
\item \eng{Firstly, social media facilitates communication and networking.} \eng{Students can join educational groups, access tutorials, and collaborate on projects.}
\item \eng{Secondly, they provide a platform for creativity and entrepreneurship.} \eng{Many young Cameroonians showcase their talents or sell products online.}
\item \eng{Furthermore, they enable rapid access to information and global awareness.}
\end{itemize}
\textbf{3. Arguments « Contre » (Inconvénients) (70--80 mots) :}
\begin{itemize}
\item \eng{However, the drawbacks are alarming.} \eng{Addiction is rampant; many youths spend hours scrolling instead of studying or sleeping.}
\item \eng{Cyberbullying and online harassment cause severe psychological trauma.}
\item \eng{Moreover, the spread of fake news and misinformation manipulates vulnerable minds.}
\item \eng{Privacy is often compromised as personal data is collected and sold.}
\end{itemize}
\textbf{4. Conclusion (40--50 mots) :}
\begin{itemize}
\item \eng{In conclusion, social media is a double-edged sword.} \eng{If used wisely, it empowers; if abused, it destroys.}
\item \eng{Parents, schools and governments must educate the youth on digital literacy and responsible use.} \eng{Only then can we harness the benefits without falling prey to the pitfalls.}
\end{itemize}
\textbf{Check-list grammaticale à réviser avant de rendre :} Passif (\eng{data is collected}), Conditionnel (\eng{If used wisely...}), Relatifs (\eng{youths who spend...}), Modaux (\eng{must educate}), Comparatif (\eng{more harm than good}), Gérondif (\eng{spending hours scrolling}), Vocabulaire précis (\eng{ubiquitous, rampant, harness, prey to, double-edged sword}).
\end{methode}

\courssub{Points de vigilance pour le Bac}

\begin{attention}[Les 5 erreurs qui coûtent des points au Bac]
\begin{enumerate}
\item \textbf{Concordance des temps oubliée (Reported Speech / Conditionnel) :} Écrire \eng{He said he \textbf{will} come} au lieu de \eng{he \textbf{would} come} ; ou \eng{If I \textbf{would} know, I \textbf{would} tell you} (double \eng{would} interdit en anglais standard dans la proposition \eng{if}). \textit{Réflexe :} Verbe introducteur passé $\rightarrow$ tout recule d'un cran. \eng{If} $\rightarrow$ jamais \eng{would} dans la proposition \eng{if}.
\item \textbf{Passif mal formé (Auxiliaire \eng{be} + Participe Passé) :} Oublier \eng{be} (\eng{The book written by...}) ou mettre le verbe au passé simple au lieu du participe (\eng{The car was \textbf{stole}}). \textit{Réflexe :} \eng{be (temps du verbe actif) + V-en (colonne 3 des irréguliers)}. Vérifier : \eng{written, done, seen, taken, chosen, spoken, broken}.
\item \textbf{Faux-amis « Transparents » :} \eng{Actually} pour « actuellement », \eng{Eventually} pour « éventuellement », \eng{Sensible} pour « sensible », \eng{Library} pour « librairie », \eng{Opportunity} pour « opportunité ». \textit{Réflexe :} Liste des 20 faux-amis « tueurs » à connaître par cœur (voir fiche vocabulaire Chapter 5).
\item \textbf{Prépositions verbes/adjectifs (Calque du français) :} \eng{Depend \textbf{of}}, \eng{Listen \textbf{---}}, \eng{Wait \textbf{---}}, \eng{Interested \textbf{by}}, \eng{Married \textbf{with}}, \eng{Similar \textbf{than}}, \eng{Good \textbf{in}}, \eng{Arrive \textbf{to}}. \textit{Réflexe :} Apprendre le verbe/adjectif \textbf{avec} sa préposition (\eng{Depend \textbf{on}}, \eng{Listen \textbf{to}}, \eng{Wait \textbf{for}}, \eng{Interested \textbf{in}}, \eng{Married \textbf{to}}, \eng{Similar \textbf{to}}, \eng{Good \textbf{at}}, \eng{Arrive \textbf{at/in}}).
\item \textbf{Ordre des mots (Négation / Question / Assez / Comparatif) :} \eng{He \textbf{not} likes} (pas d'auxiliaire), \eng{Why \textbf{he is} late?} (inversion manquante), \eng{He is \textbf{enough tall}} (ordre), \eng{More \textbf{better}} (double comparatif), \eng{Prefer \textbf{than}} (préposition). \textit{Réflexe :} Structure rigide \eng{Sujet + Aux + Not + V} ; \eng{Wh + Aux + Sujet + V} ; \eng{Adj/Adv + Enough} ; \eng{Comparatif court = -er / long = more} ; \eng{Prefer ... to}.
\end{enumerate}
\textbf{Conseil final :} Relisez votre copie \textbf{spécifiquement} pour ces 5 points (passage « chasse aux erreurs » de 5 min). C'est souvent la différence entre 12/20 et 16/20.
\end{attention}

\newpage

\coursec{Exercices}

\courssub{Je vérifie mes connaissances}

\begin{exercice}[QCM : la bonne option]
Entoure la lettre correspondant à la bonne réponse.
\begin{enumerate}
\item Unless you \_\_\_ harder, you will fail the exam.\\
(a) work \quad (b) will work \quad (c) worked \quad (d) would work
\item The news \_\_\_ broadcast every evening at 8 p.m.\\
(a) is \quad (b) are \quad (c) have been \quad (d) were
\item She asked me where \_\_\_ the previous night.\\
(a) was I \quad (b) I was \quad (c) did I be \quad (d) I am
\item I am looking forward to \_\_\_ you in Douala.\\
(a) see \quad (b) seeing \quad (c) saw \quad (d) to see
\item This is the journalist \_\_\_ article won the first prize.\\
(a) who \quad (b) which \quad (c) whose \quad (d) where
\item By the time we arrived, the press conference \_\_\_.\\
(a) has started \quad (b) started \quad (c) had started \quad (d) starts
\end{enumerate}
\end{exercice}

\begin{exercice}[Vrai ou faux — justifie ta réponse]
Indique si chaque phrase est correcte (True) ou incorrecte (False) et justifie brièvement en citant la règle.
\begin{enumerate}
\item \eng{I have been living in Yaoundé since 2020.}
\item \eng{He said me that he was tired.}
\item \eng{If I were rich, I would travel around the world.}
\item \eng{She gave me many advices about the interview.}
\item \eng{The match was watched by millions of viewers.}
\item \eng{You must to switch off your phone during the news.}
\end{enumerate}
\end{exercice}

\begin{exercice}[Vocabulaire des médias]
Complète chaque phrase avec le mot anglais approprié choisi dans la liste : \eng{headline — broadcast — reporter — advertisement — editor — audience}.
\begin{enumerate}
\item The \_\_\_ interviewed the minister after the press conference in Yaoundé.
\item The main \_\_\_ of today's newspaper is about the new media law.
\item Equinoxe TV will \_\_\_ the football match live on Saturday.
\item The \_\_\_ checks every article before the newspaper is printed.
\item The programme attracts a large \_\_\_ of young viewers.
\item There was an \_\_\_ for a new mobile phone between the two news bulletins.
\end{enumerate}
\end{exercice}

\begin{exercice}[Conjugaison à compléter]
Mets le verbe entre parenthèses au temps et à la forme qui conviennent.
\begin{enumerate}
\item Right now, the presenter \_\_\_ (read) the news on CRTV.
\item Yesterday, while I \_\_\_ (listen) to the radio, the electricity went off.
\item We \_\_\_ (already / finish) our composition when the bell rang.
\item Next year, our school \_\_\_ (organise) a media workshop.
\item Newspapers \_\_\_ (sell) at every street corner in Douala.
\item If the journalist \_\_\_ (not hurry), he will miss the press conference.
\end{enumerate}
\end{exercice}

\courssub{Je m'entraîne}

\begin{exercice}[Traduction]
Traduis les phrases suivantes en anglais.
\begin{enumerate}
\item Le gouvernement a lancé une nouvelle campagne médiatique la semaine dernière.
\item On m'a dit que le journal serait publié le lendemain.
\item Si j'avais regardé le débat, je t'aurais raconté ce qui s'est passé.
\item Voici la radio que mon grand-père écoute tous les matins.
\item Il est interdit de fumer dans les studios de télévision.
\end{enumerate}
\end{exercice}

\begin{exercice}[Transformation de phrases]
Transforme chaque phrase selon la consigne entre parenthèses, sans changer le sens.
\begin{enumerate}
\item \eng{The government has launched a new media campaign.} (Turn into the passive voice.)
\item \eng{“I will watch the debate tonight,” said Mary.} (Report what Mary said.)
\item \eng{The reporter interviewed me. He works for Equinoxe TV.} (Combine the two sentences with a relative pronoun.)
\item \eng{Work hard, or you will fail the exam.} (Rewrite beginning with \eng{Unless...})
\item \eng{I don't have a smartphone, so I can't follow the news online.} (Rewrite beginning with \eng{If I had...})
\item \eng{People speak English and French in Cameroon.} (Turn into the passive voice.)
\end{enumerate}
\end{exercice}

\begin{exercice}[Texte à trous : la presse camerounaise]
Complète l'article suivant en conjuguant les verbes entre parenthèses au temps approprié (actif ou passif).
\begin{textebox}[The press in Cameroon]
Every morning, thousands of newspapers \_\_\_ (sell) in the streets of Yaoundé and Douala. The Cameroonian press \_\_\_ (exist) for more than a century, and it \_\_\_ (play) an important role in public life ever since. Last month, a new media law \_\_\_ (discuss) in Parliament, and many journalists \_\_\_ (hope) that it would protect their freedom. “We \_\_\_ (work) in difficult conditions for years,” said one reporter, “but things \_\_\_ (improve) slowly.” If the public \_\_\_ (support) independent newspapers, quality journalism will survive. At the moment, more and more young people \_\_\_ (read) the news on their phones instead of buying printed papers. Experts believe that by 2030, most information \_\_\_ (share) online.
\end{textebox}
\end{exercice}

\begin{exercice}[Appariement : corrige les erreurs]
Chaque phrase de la colonne A contient une erreur typique de francophone. Associe-la à sa correction dans la colonne B, puis explique la règle pour deux d'entre elles.
\begin{center}
\begin{tabularx}{\linewidth}{|X|X|}
\hline
\rowcolor{popblueL} \textbf{Colonne A (phrases fautives)} & \textbf{Colonne B (corrections)} \\
\hline
1. \eng{I am living in Yaoundé since 2020.} & a. \eng{She gave me a lot of advice.} \\
\hline
2. \eng{He said me that he was tired.} & b. \eng{I have been living in Yaoundé since 2020.} \\
\hline
3. \eng{If I would be rich, I would travel.} & c. \eng{He told me that he was tired.} \\
\hline
4. \eng{She gave me many advices.} & d. \eng{If I were rich, I would travel.} \\
\hline
5. \eng{I want that you come to the studio.} & e. \eng{The informations are interesting.} \\
\hline
6. \eng{The information are interesting.} & f. \eng{I want you to come to the studio.} \\
\hline
\end{tabularx}
\end{center}
\end{exercice}

\courssub{Je me prépare au Bac}

\begin{exercice}[Compréhension et langue — type Bac]
Lis attentivement le texte suivant, puis réponds aux questions en anglais, avec tes propres mots.
\begin{textebox}[Social media and young Cameroonians]
In recent years, social media have transformed the way young Cameroonians communicate. In Yaoundé, Douala and Bamenda, teenagers spend several hours a day on their phones, sharing photos, watching videos and commenting on the news. “Ten years ago, we used to wait for the evening news on television,” explains Mr Ndi, a teacher in Buea. “Today, my students know about events before I do.”

However, this revolution has a dark side. False information is spread very quickly online, and many young people believe everything they read. Last year, a rumour about a school in Douala was shared by thousands of users before anyone checked whether it was true. Journalists insist that media education should be taught in every secondary school. “If young people learned to verify sources, they would not be manipulated so easily,” says a reporter from a private television channel. Some schools have already started media clubs where students are taught how to analyse articles and identify fake news. If these initiatives were supported by the government, the situation would certainly improve.
\end{textebox}
\textbf{A. Comprehension questions}
\begin{enumerate}
\item According to the text, how do young Cameroonians use social media? (Give two activities.)
\item What does Mr Ndi mean when he says his students know about events before he does?
\item What happened in Douala last year? Answer in one sentence.
\item What solution do journalists propose, and why?
\end{enumerate}
\textbf{B. Language work}
\begin{enumerate}
\item Find in the text: (a) one sentence in the passive voice; (b) one reported statement; (c) one type 2 conditional.
\item Turn into the passive voice: \eng{Thousands of users shared the rumour.}
\item Report the sentence: \eng{“Media education should be taught in every school,” the reporter said.}
\item Complete with the correct form: \eng{If young people \_\_\_ (verify) sources, they would not be manipulated.}
\end{enumerate}
\end{exercice}

\begin{exercice}[Traduction — type Bac]
Traduis le passage suivant en anglais.
\begin{textebox}[Texte à traduire]
Hier soir, ma sœur m'a dit qu'elle avait regardé un reportage très intéressant à la télévision. Si j'avais été à la maison, je l'aurais regardé avec elle. Le reportage, qui a été réalisé par un journaliste camerounais, montrait comment les fausses informations sont fabriquées et diffusées sur les réseaux sociaux. On devrait apprendre aux élèves à vérifier ce qu'ils lisent en ligne.
\end{textebox}
\end{exercice}

\begin{exercice}[Composition — type Bac]
Choisis l'un des deux sujets suivants et rédige ta production en anglais.

\textbf{Sujet 1 (rédaction guidée).} Write a paragraph of 80–100 words entitled \eng{“If I were a journalist...”}. In your paragraph, you must use at least: one type 2 conditional, one passive form, one reported statement and one relative clause. Underline each of these four structures.

\textbf{Sujet 2 (essai).} \eng{“Social media do more harm than good to young people in Cameroon.”} Do you agree? Write an essay of 150–180 words. Organise your work in three parts: an introduction presenting the debate, a development with at least two arguments illustrated by examples, and a conclusion giving your personal opinion. Pay attention to tenses, linking words and spelling.
\end{exercice}

\newpage

\coursec{Corrigés des exercices}

\coursec{Corrigés détaillés — Révision générale (Lesson 58)}

\begin{corrige}[Exercice 1 — QCM : la bonne option]
\begin{enumerate}
\item \textbf{(a) work} — Après \cle{unless} (« à moins que »), on emploie le présent simple, jamais \eng{will} : \eng{Unless you work harder, you will fail the exam.} (« À moins que tu ne travailles plus dur, tu échoueras à l'examen. »)
\item \textbf{(a) is} — \cle{news} est un indénombrable à l'aspect pluriel : il prend toujours un verbe au singulier. \eng{The news is broadcast every evening at 8 p.m.} (« Les informations sont diffusées chaque soir à 20 h. »)
\item \textbf{(b) I was} — Dans une question rapportée, on retrouve l'ordre sujet + verbe (pas d'inversion) : \eng{She asked me where I was the previous night.} (« Elle m'a demandé où j'étais la veille au soir. »)
\item \textbf{(b) seeing} — Dans \cle{look forward to}, \eng{to} est une préposition, donc le verbe suivant est au gérondif : \eng{I am looking forward to seeing you in Douala.} (« J'ai hâte de te voir à Douala. »)
\item \textbf{(c) whose} — Il s'agit d'une possession (l'article du journaliste) : \eng{This is the journalist whose article won the first prize.} (« Voici le journaliste dont l'article a remporté le premier prix. »)
\item \textbf{(c) had started} — Avec \cle{by the time} + prétérit, l'action antérieure est au past perfect : \eng{By the time we arrived, the press conference had started.} (« Quand nous sommes arrivés, la conférence de presse avait déjà commencé. »)
\end{enumerate}
\end{corrige}

\begin{corrige}[Exercice 2 — Vrai ou faux]
\begin{enumerate}
\item \textbf{True.} \eng{I have been living in Yaoundé since 2020.} — Action commencée dans le passé et qui continue : present perfect continuous + \cle{since} + point de départ. Correct.
\item \textbf{False.} Correction : \eng{He told me that he was tired.} — \cle{say} ne prend pas de complément de personne directement ; on dit \eng{say something (to somebody)} ou \eng{tell somebody something}.
\item \textbf{True.} \eng{If I were rich, I would travel around the world.} — Conditionnel de type 2 (hypothèse irréelle au présent) : \eng{if} + prétérit, \eng{would} + base verbale ; \eng{were} pour toutes les personnes. Correct.
\item \textbf{False.} Correction : \eng{She gave me a lot of advice about the interview.} — \cle{advice} est indénombrable : jamais de \eng{-s}, jamais \eng{many} ; on dit \eng{some advice}, \eng{a piece of advice}.
\item \textbf{True.} \eng{The match was watched by millions of viewers.} — Passif au prétérit correctement formé : \eng{was} + participe passé \eng{watched}. Correct.
\item \textbf{False.} Correction : \eng{You must switch off your phone during the news.} — Les auxiliaires modaux (\cle{must}, \cle{can}, \cle{should}...) sont suivis directement de la base verbale, sans \eng{to}.
\end{enumerate}
\end{corrige}

\begin{corrige}[Exercice 3 — Vocabulaire des médias]
\begin{enumerate}
\item \eng{The \textbf{reporter} interviewed the minister after the press conference in Yaoundé.} (« Le journaliste a interviewé le ministre après la conférence de presse à Yaoundé. »)
\item \eng{The main \textbf{headline} of today's newspaper is about the new media law.} (« Le grand titre du journal d'aujourd'hui concerne la nouvelle loi sur les médias. »)
\item \eng{Equinoxe TV will \textbf{broadcast} the football match live on Saturday.} (« Equinoxe TV diffusera le match de football en direct samedi. »)
\item \eng{The \textbf{editor} checks every article before the newspaper is printed.} (« Le rédacteur en chef vérifie chaque article avant que le journal ne soit imprimé. »)
\item \eng{The programme attracts a large \textbf{audience} of young viewers.} (« L'émission attire un large public de jeunes téléspectateurs. »)
\item \eng{There was an \textbf{advertisement} for a new mobile phone between the two news bulletins.} (« Il y avait une publicité pour un nouveau téléphone portable entre les deux journaux télévisés. »)
\end{enumerate}
\textbf{Point de vigilance :} ne pas confondre \eng{advertisement} (« annonce publicitaire », prononcé « ad-vèr-tiss-ment ») avec \eng{advertising} (la publicité comme activité, indénombrable).
\end{corrige}

\begin{corrige}[Exercice 4 — Conjugaison à compléter]
\begin{enumerate}
\item \eng{Right now, the presenter \textbf{is reading} the news on CRTV.} — Présent continu : action en cours au moment où l'on parle (\eng{right now}).
\item \eng{Yesterday, while I \textbf{was listening} to the radio, the electricity went off.} — Prétérit continu après \cle{while} : action longue interrompue par une action brève au prétérit simple.
\item \eng{We \textbf{had already finished} our composition when the bell rang.} — Past perfect : action antérieure à une autre action passée (\eng{when the bell rang}).
\item \eng{Next year, our school \textbf{will organise} a media workshop.} — Futur simple avec \eng{next year} (projet prévu ; \eng{is going to organise} est aussi accepté).
\item \eng{Newspapers \textbf{are sold} at every street corner in Douala.} — Passif au présent simple : vérité générale ; le sujet subit l'action.
\item \eng{If the journalist \textbf{does not hurry}, he will miss the press conference.} — Conditionnel de type 1 : \eng{if} + présent simple, futur dans la principale. Jamais \eng{will} après \eng{if}.
\end{enumerate}
\end{corrige}

\begin{corrige}[Exercice 5 — Traduction]
\begin{enumerate}
\item \eng{The government launched a new media campaign last week.} — Prétérit simple avec \eng{last week} ; attention : pas de \eng{has launched} avec un repère passé terminé.
\item \eng{I was told that the newspaper would be published the next day.} — Passif (\eng{I was told}) + discours rapporté : \eng{will} devient \eng{would}, \eng{tomorrow} devient \eng{the next day}.
\item \eng{If I had watched the debate, I would have told you what happened.} — Conditionnel de type 3 : \eng{if} + past perfect, \eng{would have} + participe passé.
\item \eng{This is the radio (that) my grandfather listens to every morning.} — Relatif \eng{that} (facultatif) ; ne pas oublier la préposition finale \eng{to} (\eng{listen to}).
\item \eng{Smoking is not allowed in television studios.} ou \eng{It is forbidden to smoke in television studios.} — « Il est interdit de » se rend par \eng{it is forbidden to} ou \eng{... is not allowed}, jamais par \eng{it is interdicted} (faux ami).
\end{enumerate}
\end{corrige}

\begin{corrige}[Exercice 6 — Transformation de phrases]
\begin{enumerate}
\item \eng{A new media campaign has been launched by the government.} — Passif au present perfect : \eng{has been} + participe passé ; l'objet devient sujet.
\item \eng{Mary said (that) she would watch the debate that night.} — Recul des temps : \eng{will} $\rightarrow$ \eng{would} ; \eng{tonight} $\rightarrow$ \eng{that night}.
\item \eng{The reporter who interviewed me works for Equinoxe TV.} (ou \eng{The reporter who works for Equinoxe TV interviewed me.}) — Relatif \cle{who} pour une personne ; pas de virgule car la relative est déterminative.
\item \eng{Unless you work hard, you will fail the exam.} — \cle{unless} = \eng{if... not} ; présent simple après \eng{unless}.
\item \eng{If I had a smartphone, I could follow the news online.} — Conditionnel de type 2 ; \eng{can} devient \eng{could}.
\item \eng{English and French are spoken in Cameroon.} — Passif au présent simple ; deux langues = sujet pluriel : \eng{are spoken}.
\end{enumerate}
\end{corrige}

\begin{corrige}[Exercice 7 — Texte à trous : la presse camerounaise]
\begin{enumerate}
\item \eng{are sold} — passif, présent simple (\eng{every morning} : habitude).
\item \eng{has existed} — present perfect avec \eng{for more than a century} (durée jusqu'à aujourd'hui).
\item \eng{has played} — present perfect avec \eng{ever since}.
\item \eng{was discussed} — passif au prétérit (\eng{last month}).
\item \eng{hoped} — prétérit simple, en cohérence avec le récit au passé.
\item \eng{have been working} ou \eng{have worked} — present perfect (continuous) avec \eng{for years} : situation qui dure encore.
\item \eng{are improving} — présent continu : changement en cours (\eng{slowly}, aujourd'hui).
\item \eng{supports} — présent simple après \eng{if} (conditionnel de type 1).
\item \eng{are reading} — présent continu avec \eng{at the moment}.
\item \eng{will be shared} — futur passif avec \eng{by 2030} (\eng{will have been shared} est également correct si l'on insiste sur l'antériorité).
\end{enumerate}
\textbf{Texte complet :}
\begin{textebox}[Texte reconstitué]
Every morning, thousands of newspapers are sold in the streets of Yaoundé and Douala. The Cameroonian press has existed for more than a century, and it has played an important role in public life ever since. Last month, a new media law was discussed in Parliament, and many journalists hoped that it would protect their freedom. “We have been working in difficult conditions for years,” said one reporter, “but things are improving slowly.” If the public supports independent newspapers, quality journalism will survive. At the moment, more and more young people are reading the news on their phones instead of buying printed papers. Experts believe that by 2030, most information will be shared online.
\end{textebox}
\end{corrige}

\begin{corrige}[Exercice 8 — Appariement : corrige les erreurs]
\textbf{Appariements :} 1 $\rightarrow$ b ; 2 $\rightarrow$ c ; 3 $\rightarrow$ d ; 4 $\rightarrow$ a ; 5 $\rightarrow$ f ; 6 $\rightarrow$ e.

\textbf{Explications demandées (deux au choix) :}
\begin{itemize}
\item \textbf{1 $\rightarrow$ b :} \eng{I have been living in Yaoundé since 2020.} — Avec \cle{since} + point de départ et une action qui continue, l'anglais exige le present perfect (continuous), alors que le français utilise le présent (« j'habite à Yaoundé depuis 2020 »). Le présent simple anglais exprime une habitude générale, pas une durée inachevée.
\item \textbf{2 $\rightarrow$ c :} \eng{He told me that he was tired.} — \cle{say} ne peut pas être suivi directement d'un complément de personne : on dit \eng{he said to me} (rare) ou \eng{he told me}. C'est le verbe \cle{tell} qui prend une personne pour complément.
\item \textbf{3 $\rightarrow$ d :} \eng{If I were rich, I would travel.} — Dans le conditionnel de type 2, \eng{would} ne se met jamais dans la subordonnée introduite par \eng{if} : \eng{if} + prétérit, \eng{would} dans la principale.
\item \textbf{4 $\rightarrow$ a :} \eng{She gave me a lot of advice.} — \cle{advice} est indénombrable : pas de pluriel en \eng{-s}, pas de \eng{many} ; on utilise \eng{some}, \eng{a lot of}, \eng{a piece of advice}.
\item \textbf{5 $\rightarrow$ f :} \eng{I want you to come to the studio.} — Après \cle{want}, l'anglais n'emploie pas de subordonnée en \eng{that} ; structure : \eng{want} + complément + \eng{to} + base verbale.
\item \textbf{6 $\rightarrow$ e :} \eng{The information is interesting.} — \cle{information} est indénombrable : singulier obligatoire (\eng{is}), jamais \eng{informations} comme en français.
\end{itemize}
\end{corrige}

\begin{corrige}[Exercice 9 — Compréhension et langue — type Bac]
\textbf{Barème indicatif (sur 20) :} A. Compréhension : 4 x 2 pts = 8 pts ; B. Langue : 4 x 2 pts = 8 pts ; qualité globale de l'expression (correction grammaticale, propres mots) : 4 pts.

\textbf{A. Comprehension questions (réponses modèles)}
\begin{enumerate}
\item \eng{They use social media to share photos and watch videos. They also comment on the news.} (Deux activités suffisent ; toute formulation correcte avec ses propres mots est acceptée.)
\item \eng{He means that his students learn about events on social media very quickly, even before their teacher hears about them.} (Les élèves sont informés avant lui grâce aux réseaux sociaux.)
\item \eng{Last year, a false rumour about a school in Douala was shared by thousands of users before anyone checked it.}
\item \eng{Journalists propose that media education should be taught in every secondary school, so that young people can learn to verify sources and avoid being manipulated.}
\end{enumerate}

\textbf{B. Language work}
\begin{enumerate}
\item (a) Passif : \eng{False information is spread very quickly online.} (ou \eng{...a rumour about a school in Douala was shared by thousands of users...} / \eng{...students are taught how to analyse articles...}) — (b) Discours rapporté : \eng{Journalists insist that media education should be taught in every secondary school.} — (c) Conditionnel de type 2 : \eng{If young people learned to verify sources, they would not be manipulated so easily.} (ou la dernière phrase du texte).
\item \eng{The rumour was shared by thousands of users.} — Passif au prétérit : \eng{was} + participe passé.
\item \eng{The reporter said (that) media education should be taught in every school.} — \eng{should} est invariable au discours rapporté ; la phrase reste au passif.
\item \eng{If young people \textbf{verified} sources, they would not be manipulated.} — Prétérit simple après \eng{if} (type 2).
\end{enumerate}
\textbf{Points de vigilance :} répondre par des phrases complètes ; ne pas recopier le texte mot à mot ; vérifier l'accord sujet-verbe et l'emploi du prétérit pour les événements passés.
\end{corrige}

\begin{corrige}[Exercice 10 — Traduction — type Bac]
\textbf{Barème indicatif (sur 10) :} 2 pts par phrase (1 pt pour le sens, 1 pt pour la correction grammaticale).

\textbf{Traduction modèle :}
\begin{textebox}[Modèle]
Last night, my sister told me (that) she had watched a very interesting documentary on television. If I had been at home, I would have watched it with her. The documentary, which was made by a Cameroonian journalist, showed how false information is created and spread on social media. Students should be taught to check what they read online.
\end{textebox}

\textbf{Points de vigilance :}
\begin{itemize}
\item « m'a dit qu'elle avait regardé » : \eng{told me she had watched} — \cle{tell} + personne ; past perfect pour l'antériorité au discours rapporté.
\item « Si j'avais été... je l'aurais regardé » : conditionnel de type 3 — \eng{If I had been..., I would have watched...} ; \eng{it} pour reprendre \eng{documentary} (pas \eng{him}).
\item « qui a été réalisé » : relative non déterminative entre virgules avec \cle{which} + passif au prétérit (\eng{was made}).
\item « sont fabriquées et diffusées » : passif au présent, sujet singulier \eng{information} $\rightarrow$ \eng{is created and spread}.
\item « On devrait apprendre aux élèves » : passif avec \eng{should be taught}, plus naturel que \eng{we should teach} ici.
\end{itemize}
\end{corrige}

\begin{corrige}[Exercice 11 — Composition — type Bac]
\textbf{Barème indicatif (sur 20) :} respect de la consigne et du nombre de mots : 4 pts ; richesse et correction grammaticale (structures exigées) : 6 pts ; vocabulaire approprié : 4 pts ; organisation et connecteurs : 3 pts ; orthographe et ponctuation : 3 pts.

\textbf{Sujet 1 — Production modèle :}
\begin{textebox}[Sujet 1 : If I were a journalist...]
If I were a journalist, I would work for a television channel in Douala. I would investigate stories that are hidden from the public, because I believe the truth should be told to everyone. My editor once said that courage was the most important quality for a reporter, and I completely agree. I would also interview ordinary people whose voices are never heard, such as market women and bendskin drivers. Being a journalist is a job which requires honesty and hard work, and I would be proud to serve my country in this way.
\end{textebox}
(96 mots ; structures à souligner : \eng{If I were..., I would work...} = type 2 ; \eng{are hidden}, \eng{should be told}, \eng{are never heard} = passifs ; \eng{My editor once said that courage was...} = discours rapporté ; \eng{people whose voices...}, \eng{a job which requires...} = relatives.)

\textbf{Sujet 2 — Production modèle :}
\begin{textebox}[Sujet 2 : Social media and young Cameroonians]
Nowadays, social media are everywhere in the lives of young Cameroonians. While some people believe they are useful tools, others think they do more harm than good.

On the one hand, social media allow teenagers to communicate easily and to follow the news quickly. For example, students in Yaoundé can share lessons and help each other with homework. On the other hand, false information is spread very fast online, and many young people believe everything they read. Moreover, spending several hours a day on the phone can damage their studies and their health.

In my opinion, social media are neither good nor bad in themselves: everything depends on how they are used. If media education were taught in every school, young people would benefit from social media without suffering from their dangers. For instance, I personally learnt how to spot fake news thanks to a workshop organised by my school last year.
\end{textebox}
(Environ 155 mots — longueur idéale.)

\textbf{Points de vigilance :}
\begin{itemize}
\item Respecter la structure en trois parties (introduction / développement / conclusion) et annoncer clairement sa position.
\item Utiliser des connecteurs : \eng{on the one hand / on the other hand}, \eng{moreover}, \eng{for example}, \eng{in my opinion}, \eng{for instance}, \eng{to conclude}.
\item \eng{social media} est pluriel : \eng{social media \textbf{are}...}.
\item Compter ses mots et relire : accords, temps, orthographe (\eng{believe}, \eng{which}, \eng{their}).
\end{itemize}
\end{corrige}

\newpage

\coursec{Fiche bilan}

\begin{popfigure}
\begin{tikzpicture}[
  mindmap,
  grow cyclic,
  every node/.style={concept, align=center, font=\small},
  root concept/.append style={concept color=popblue, font=\bfseries\small, minimum size=3.2cm},
  level 1 concept/.append style={level distance=4.6cm, sibling angle=51.4, font=\footnotesize},
  level 2 concept/.append style={level distance=2.6cm, font=\scriptsize}
]
\node[root concept, align=center]{General\\Revision\\Grammar}
  child[concept color=poporange]{ node[align=center]{Tenses\\present, past,\\future, perfect}
    child { node[concept color=poporangeL, align=center]{for / since\\already / yet} }
  }
  child[concept color=popgreen]{ node[align=center]{Passive\\voice}
    child { node[concept color=popgreenL, align=center]{be + past\\participle} }
  }
  child[concept color=popgold]{ node[align=center]{Reported\\speech}
    child { node[concept color=popgoldL, align=center]{backshift\\say / tell / ask} }
  }
  child[concept color=poppurple]{ node[align=center]{Conditionals\\wishes}
    child { node[concept color=poppurpleL, align=center]{if + past,\\would + V\\I wish...} }
  }
  child[concept color=poppink]{ node[align=center]{Relatives\\modals}
    child { node[concept color=poppinkL, align=center]{who / which /\\whose; must,\\should, may} }
  }
  child[concept color=popblue]{ node[align=center]{Gerund vs\\infinitive}
    child { node[concept color=popblueL, align=center]{enjoy + -ing\\want + to-V} }
  }
  child[concept color=popgreen]{ node[align=center]{Comparison\\quantifiers}
    child { node[concept color=popgreenL, align=center]{more / -er\\the most / -est\\much / many} }
  };
\end{tikzpicture}
\legende{Carte mentale : les sept piliers de la révision grammaticale de Terminale.}
\end{popfigure}

\begin{bilan}
\textbf{1. Lexique clé de la leçon (anglais — français)}
\begin{itemize}
  \item \eng{a news report} : un reportage ; \eng{a headline} : un titre (de presse) ; \eng{a broadcast} : une émission.
  \item \eng{to report} : rapporter (des propos) ; \eng{to claim} : affirmer ; \eng{to deny} : nier.
  \item \eng{reliable} : fiable ; \eng{misleading} : trompeur ; \eng{biased} : partial.
  \item \eng{to broadcast / broadcast / broadcast} : diffuser (verbe irrégulier invariable).
\end{itemize}

\textbf{2. Règles de grammaire à connaître par cœur}
\begin{center}
\begin{tabularx}{\linewidth}{|l|X|X|}
\hline
\rowcolor{popblueL} \textbf{Point} & \textbf{Règle} & \textbf{Exemple} \\
\hline
Present perfect & have/has + participe passé ; bilan présent, avec \eng{for, since, already, yet} & \eng{She has worked here since 2020.} \\
\hline
Past simple & action datée et terminée (\eng{yesterday, ago, last}) & \eng{He left Douala two days ago.} \\
\hline
Passive voice & be (au temps voulu) + participe passé (+ by) & \eng{The match was watched by millions.} \\
\hline
Reported speech & recul des temps ; \eng{say (to), tell + COI, ask, whether/if} & \eng{He said (that) he was tired.} \\
\hline
Conditionals & 1 : if + present, will + V ; 2 : if + past, would + V ; 3 : if + past perfect, would have + pp & \eng{If I had money, I would travel.} \\
\hline
Wishes & \eng{I wish} + past (présent) ; + past perfect (regret) & \eng{I wish I had studied harder.} \\
\hline
Relatives & \eng{who} (personnes), \eng{which} (choses), \eng{whose} (possession), \eng{where/when} & \eng{The reporter who interviewed me...} \\
\hline
Modaux & modal + base verbale ; \eng{must} (obligation), \eng{should} (conseil), \eng{may/might} (possibilité), \eng{can't} (impossibilité) & \eng{You should check your sources.} \\
\hline
Comparison & court : -er / the -est ; long : more / the most ; \eng{as...as} & \eng{Radio is cheaper than TV.} \\
\hline
Gerund / infinitive & \eng{enjoy, avoid, suggest} + -ing ; \eng{want, decide, hope} + to-V & \eng{I enjoy reading the news.} \\
\hline
\end{tabularx}
\end{center}

\textbf{3. Méthodes express}
\begin{enumerate}
  \item \textbf{Choisir un temps} : repérer le marqueur (\eng{ago} $\rightarrow$ past simple ; \eng{since/for} $\rightarrow$ present perfect ; \eng{tomorrow} $\rightarrow$ future).
  \item \textbf{Passer au passif} : COD devient sujet ; be au temps du verbe actif ; participe passé ; agent avec \eng{by} (souvent omis).
  \item \textbf{Rapporter une parole} : reculer le temps d'un cran, adapter pronoms et marqueurs (\eng{now} $\rightarrow$ \eng{then}, \eng{today} $\rightarrow$ \eng{that day}).
  \item \textbf{Phrase conditionnelle} : jamais de \eng{will/would} dans la proposition en \eng{if}.
  \item \textbf{Avant d'écrire} : verbe $\rightarrow$ temps ? voix ? modal ? ; nom $\rightarrow$ article, pluriel ; adjectif $\rightarrow$ invariable.
\end{enumerate}
\end{bilan}

\begin{aretenir}[Checklist avant le Bac]
Cochez chaque point maîtrisé :
\begin{itemize}
  \item $\square$ Je conjugue sans erreur le present perfect et je le distingue du past simple.
  \item $\square$ Je connais les participes passés des verbes irréguliers courants (\eng{go--went--gone}, \eng{write--wrote--written}, \eng{take--took--taken}).
  \item $\square$ Je transforme une phrase active en phrase passive correctement.
  \item $\square$ J'applique le recul des temps au discours rapporté.
  \item $\square$ Je choisis le bon conditionnel (1, 2 ou 3) selon le degré de réalité.
  \item $\square$ J'emploie \eng{I wish} + past / past perfect pour exprimer un souhait ou un regret.
  \item $\square$ Je choisis le bon pronom relatif (\eng{who, which, whose, where}).
  \item $\square$ Je place correctement les modaux suivis de la base verbale (jamais de \eng{to} après un modal).
  \item $\square$ Je sais si un verbe est suivi du gérondif (-ing) ou de l'infinitif (to-V).
  \item $\square$ Je forme les comparatifs et superlatifs sans faute (\eng{better, the best, more interesting}).
  \item $\square$ J'évite les pièges francophones : \eng{informations} $\rightarrow$ \eng{information} (indénombrable), \eng{to be agree} $\rightarrow$ \eng{to agree}.
  \item $\square$ Je relis ma production : sujet--verbe accordés, temps cohérents, orthographe vérifiée.
\end{itemize}
\end{aretenir}

\findecours

\end{document}
